\chapter{Behaviour, Fibrations, and the Categorical Myhill--Nerode Theorem}
\label{ch:behaviour-categories-myhill-nerode}

\section{Fibred orientation and standing assumptions}
\label{sec:behaviour-standing-assumptions}

The Myhill--Nerode theorem characterizes regular behaviour by residuals. For a language, the residuals are the quotients of the language by input words. For a Mealy morphism between internal categories, the residual inputs over an object are organized by a slice internal category. This chapter develops the corresponding residual semantics for the Mealy morphisms of Chapter~\ref{ch:spans-internal-categories-mealy}. The language-theoretic residual quotient goes back to Myhill and Nerode, and the derivative viewpoint is closely related to Brzozowski's calculus of regular-expression derivatives \cite{Myhill1957,Nerode1958,Brzozowski1964}. The categorical formulation below is also related in spirit to the final-semantics and realization-theoretic approaches to automata \cite{Goguen1973,Rutten1998,Rutten2000}.

The input internal category \(\mathbb A\) determines a residual syntax opfibration
\[
    \operatorname{Res}_{\mathbb A}\longrightarrow \mathbb A.
\]
The fibre over an object \(v\in A_0\) is the residual slice
\[
    \mathbb A/v.
\]
An input arrow
\[
    a:u\to v
\]
acts on residual syntax by cocartesian transport, given by postcomposition:
\[
    a_*:\mathbb A/u\to\mathbb A/v.
\]
Thus residual syntax is covariantly indexed by \(\mathbb A\).

The behaviour assignment is contravariant in residual syntax. A residual behaviour over \(v\) is an internal functor
\[
    H:\mathbb A/v\to\mathbb B.
\]
Reindexing along
\[
    a:u\to v
\]
is precomposition with residual transport:
\[
    a^*H:=H\circ a_*.
\]
This reindexing is the derivative of \(H\) by \(a\):
\[
    \partial_aH:=a^*H=H\circ a_*.
\]
Hence the behaviour fibration is the contravariantly indexed internal category
\[
    \mathbf{Beh}_{\mathbb A,\mathbb B}\longrightarrow \mathbb A
\]
whose fibre over \(v\) is
\[
    \mathbf{Beh}_{\mathbb A,\mathbb B}(v)
    =
    [\mathbb A/v,\mathbb B].
\]

The behaviour fibration carries a canonical Mealy morphism. Its transition is cartesian reindexing,
\[
    H\longmapsto a^*H=\partial_aH,
\]
and its output is the anchor comparison obtained from the residual morphism
\[
    a_*1_u=a\longrightarrow 1_v
\]
in the slice \(\mathbb A/v\).

\begin{assumption}
\label{ass:behaviour-base}
    Throughout this chapter, \(\mathcal E\) is a Barr-exact locally cartesian closed category with chosen pullbacks. Thus \(\mathcal E\) has finite limits, dependent products in slices, stable regular epi--mono factorizations, and effective quotients of internal equivalence relations.
\end{assumption}

\begin{notation}
\label{not:behaviour-internal-categories}
    Let
    \[
        \mathbb A=(A_0,A_1,s_A,t_A,e_A,m_A),
        \qquad
        \mathbb B=(B_0,B_1,s_B,t_B,e_B,m_B)
    \]
    be internal categories in \(\mathcal E\). We use the composition convention of Chapter~\ref{ch:spans-internal-categories-mealy}: if
    \[
        a:u\to v,
        \qquad
        b:v\to w,
    \]
    then
    \[
        m_A(a,b)=b\circ a.
    \]
    Thus a Mealy state over \(w\) processes the input \(b\circ a\) by first processing \(b\), then processing \(a\). We write
    \[
        1_v:=e_A(v)
    \]
    for the identity arrow at \(v\), and similarly in \(\mathbb B\).
\end{notation}

\begin{notation}
\label{not:behaviour-generalized-elements}
    Generalized-element notation abbreviates morphism equations in \(\mathcal E\). When variables such as \(a,r,x,H\) appear below, the corresponding formal assertion is the equality of maps out of the pullback object expressing the displayed typing constraints.
\end{notation}

\begin{notation}
\label{not:mealy-recall-behaviour}
    A Mealy morphism
    \[
        P:\mathbb A\rightsquigarrow\mathbb B
    \]
    has carrier span
    \[
        A_0\xleftarrow{p}P\xrightarrow{q}B_0
    \]
    and structure maps
    \[
        \delta:
        A_1\times_{t_A,A_0,p}P\to P,
        \qquad
        \omega:
        A_1\times_{t_A,A_0,p}P\to B_1.
    \]
    For \(a:u\to v\) and \(p(x)=v\), the input is processed target-to-source:
    \[
        p(\delta(a,x))=u,
    \]
    and the output arrow is
    \[
        \omega(a,x):q(\delta(a,x))\to q(x).
    \]
    The Mealy unit laws are
    \[
        \delta(1_{p(x)},x)=x,
        \qquad
        \omega(1_{p(x)},x)=1_{q(x)},
    \]
    and the Mealy multiplication laws are
    \[
        \delta(b\circ a,x)=\delta(a,\delta(b,x)),
    \]
    \[
        \omega(b\circ a,x)
        =
        \omega(b,x)\circ\omega(a,\delta(b,x)).
    \]
\end{notation}

\begin{notation}
\label{not:simulation-recall-behaviour}
    A Mealy simulation
    \[
        (\theta,\alpha):P\Rightarrow Q
    \]
    has underlying state map in the opposite direction
    \[
        \theta:Q\to P,
    \]
    and an output-comparison component
    \[
        \alpha:Q\to B_1
    \]
    satisfying
    \[
        p_P\theta=p_Q,
        \qquad
        s_B\alpha=q_P\theta,
        \qquad
        t_B\alpha=q_Q.
    \]
    Thus
    \[
        \alpha_y:q_P(\theta y)\to q_Q(y).
    \]
    The simulation laws are
    \[
        \theta(\delta_Q(a,y))
        =
        \delta_P(a,\theta y),
    \]
    and
    \[
        \alpha_y\circ\omega_P(a,\theta y)
        =
        \omega_Q(a,y)\circ\alpha_{\delta_Q(a,y)}.
    \]
\end{notation}

\begin{remark}
\label{rem:simulation-direction-warning-behaviour}
    A displayed simulation
    \[
        P\Rightarrow Q
    \]
    has state map \(Q\to P\). The strict semantic homomorphisms used in Section~\ref{sec:terminal-residual-semantics} have covariantly directed state maps \(P\to Q\).
\end{remark}

\section{Residual syntax and the behaviour fibration}
\label{sec:residual-syntax-behaviour-fibration}

The residual syntax of an input category is the indexed-category presentation of the codomain opfibration. We use the split cleavage given by postcomposition.

\subsection{Residual syntax as a split codomain opfibration}
\label{subsec:residual-opfibration}

\begin{definition}[Residual syntax]
\label{def:residual-syntax}
    The \textbf{vertical residual category} of \(\mathbb A\) is the internal
    category
    \[
        \mathbb R_{\mathbb A}
    \]
    in the slice \(\mathcal E/A_0\) defined as follows.

    Its object object is
    \[
        (\mathbb R_{\mathbb A})_0
        :=
        \left(
            A_1\xrightarrow{t_A}A_0
        \right).
    \]
    Its arrow object is
    \[
        (\mathbb R_{\mathbb A})_1
        :=
        \left(
            T_A\xrightarrow{t_A\pi_2}A_0
        \right),
    \]
    where
    \[
        T_A
        :=
        A_1\times_{t_A,A_0,s_A}A_1.
    \]
    Thus a generalized element of \(T_A\) is a composable pair
    \[
        z\xrightarrow{c}w\xrightarrow{r}v.
    \]

    The source and target maps are
    \[
        s_{\mathbb R}
        :=
        m_A:T_A\to A_1,
        \qquad
        t_{\mathbb R}
        :=
        \pi_2:T_A\to A_1.
    \]
    In generalized-element notation,
    \[
        s_{\mathbb R}(c,r)=r\circ c,
        \qquad
        t_{\mathbb R}(c,r)=r.
    \]
    Hence \((c,r)\) represents the residual arrow
    \[
        r\circ c\longrightarrow r.
    \]

    The identity map is
    \[
        e_{\mathbb R}
        :=
        \left\langle
            e_As_A,
            1_{A_1}
        \right\rangle:
        A_1\to T_A,
    \]
    so that
    \[
        e_{\mathbb R}(r)=(1_{s_A(r)},r).
    \]

    Composition is induced by composition in \(\mathbb A\). On a residual
    triple
    \[
        y\xrightarrow{d}z\xrightarrow{c}w\xrightarrow{r}v,
    \]
    it is
    \[
        (c,r)\circ(d,r\circ c)
        =
        (c\circ d,r).
    \]

    Residual postcomposition is represented internally as follows. Pull back
    \(\mathbb R_{\mathbb A}\) along
    \[
        s_A,t_A:A_1\to A_0.
    \]
    There is an internal functor in \(\mathcal E/A_1\)
    \[
        \mathsf{post}:
        s_A^*\mathbb R_{\mathbb A}
        \longrightarrow
        t_A^*\mathbb R_{\mathbb A}.
    \]
    Its object map is
    \[
        (a,r)\longmapsto(a,a\circ r),
    \]
    where
    \[
        r:w\to s_A(a),
    \]
    and its arrow map is
    \[
        (a,c,r)\longmapsto(a,c,a\circ r).
    \]
    Fibrewise, over an arrow
    \[
        a:u\to v,
    \]
    this internal functor is the postcomposition functor
    \[
        a_*:\mathbb A/u\to\mathbb A/v.
    \]

    The pair
    \[
        \bigl(
            \mathbb R_{\mathbb A},
            \mathsf{post}
        \bigr)
    \]
    is called the \textbf{residual syntax} of \(\mathbb A\).
\end{definition}

The defining pullback for the residual arrow object is
\[
    \begin{tikzcd}[column sep=large, row sep=large]
            T_A
                \arrow[r, "\pi_2"]
                \arrow[d, "\pi_1"']
                \arrow[dr, phantom, "\lrcorner", very near start]
        &
            A_1
                \arrow[d, "s_A"]
        \\
            A_1
                \arrow[r, "t_A"']
        &
            A_0.
    \end{tikzcd}
\]
A generalized element is a pair
\[
    (c,r)
\]
with
\[
    z\xrightarrow{c}w\xrightarrow{r}v.
\]
It represents the residual triangle
\[
    \begin{tikzcd}[column sep=large, row sep=large]
            z
                \arrow[rr, "r\circ c", bend left=15]
                \arrow[r, "c"']
        &
            w
                \arrow[r, "r"']
        &
            v.
    \end{tikzcd}
\]

\begin{proposition}[Residual fibres]
\label{prop:residual-fibres-in-slice}
    The data defining \(\mathbb R_{\mathbb A}\) form an internal category in
    \[
        \mathcal E/A_0.
    \]
    Its fibre over an object \(v\) of \(\mathbb A\) is the slice internal
    category
    \[
        \mathbb A/v.
    \]
\end{proposition}
\begin{proof}
    \leavevmode
    \begin{enumerate}
        \item \textbf{Verify the structure maps lie over \(A_0\).}

        Since
        \[
            t_Am_A=t_A\pi_2
        \]
        on \(T_A\), the source map
        \[
            s_{\mathbb R}=m_A
        \]
        is a morphism in \(\mathcal E/A_0\). The target map
        \[
            t_{\mathbb R}=\pi_2
        \]
        is also a morphism over \(A_0\).

        The identity map lies over \(A_0\) because
        \[
            t_Ae_As_A=s_A
        \]
        and
        \[
            t_A\pi_2e_{\mathbb R}
            =
            t_A.
        \]

        \item \textbf{Verify source and target of identities.}

        For \(r:w\to v\),
        \[
            s_{\mathbb R}e_{\mathbb R}(r)
            =
            r\circ1_w
            =
            r,
        \]
        and
        \[
            t_{\mathbb R}e_{\mathbb R}(r)=r.
        \]

        \item \textbf{Verify composition.}

        A composable pair of residual arrows is represented by a diagram
        \[
            y\xrightarrow{d}z\xrightarrow{c}w\xrightarrow{r}v.
        \]
        The two residual arrows are
        \[
            r\circ c\circ d\longrightarrow r\circ c
        \]
        and
        \[
            r\circ c\longrightarrow r.
        \]
        Their composite is represented by
        \[
            (c\circ d,r):
            r\circ(c\circ d)\longrightarrow r.
        \]
        Associativity in \(\mathbb A\) identifies
        \[
            r\circ(c\circ d)
            =
            (r\circ c)\circ d,
        \]
        so this has the required source.

        \item \textbf{Verify the category laws.}

        The two unit equations reduce to the unit equations in \(\mathbb A\).
        The associativity equation reduces to associativity in \(\mathbb A\)
        on a string
        \[
            x\longrightarrow y\longrightarrow z
            \longrightarrow w\longrightarrow v.
        \]
        Therefore \(\mathbb R_{\mathbb A}\) is an internal category in
        \(\mathcal E/A_0\).

        \item \textbf{Identify the fibres.}

        Pulling \(\mathbb R_{\mathbb A}\) back to an object \(v\) gives
        objects
        \[
            r:w\to v
        \]
        and arrows
        \[
            r\circ c\to r
        \]
        represented by commuting triangles over \(v\). These are exactly the
        objects and arrows of
        \[
            \mathbb A/v.
        \]
    \end{enumerate}
\end{proof}

\begin{definition}[Cocartesian residual transport]
\label{def:residual-transport}
    For an arrow
    \[
        a:u\to v
    \]
    of \(\mathbb A\), the \textbf{cocartesian residual transport} functor
    \[
        a_*:\mathbb A/u\to\mathbb A/v
    \]
    is the fibre of
    \[
        \mathsf{post}:
        s_A^*\mathbb R_{\mathbb A}
        \longrightarrow
        t_A^*\mathbb R_{\mathbb A}
    \]
    over \(a\).

    Explicitly,
    \[
        a_*(r)=a\circ r
    \]
    for \(r:w\to u\), and
    \[
        a_*(c,r)=(c,a\circ r)
    \]
    for a residual arrow
    \[
        r\circ c\to r.
    \]
\end{definition}

The chosen cocartesian residual square is
\[
    \begin{tikzcd}[column sep=large, row sep=large]
            w
                \arrow[r, equal]
                \arrow[d, "r"']
        &
            w
                \arrow[d, "a\circ r"]
        \\
            u
                \arrow[r, "a"']
        &
            v.
    \end{tikzcd}
\]

\begin{proposition}[Split opfibration laws]
\label{prop:residual-transport-laws}
    Residual transport satisfies
    \[
        (1_v)_*
        =
        1_{\mathbb A/v},
    \]
    and, for composable arrows
    \[
        a:u\to v,
        \qquad
        b:v\to w,
    \]
    \[
        (b\circ a)_*
        =
        b_*\circ a_*.
    \]
    These equations hold on both residual objects and residual arrows.
\end{proposition}
\begin{proof}
    On a residual object \(r:z\to u\),
    \[
        (1_u)_*(r)
        =
        1_u\circ r
        =
        r.
    \]
    For composable \(a\) and \(b\),
    \[
        (b\circ a)_*(r)
        =
        (b\circ a)\circ r
        =
        b\circ(a\circ r)
        =
        b_*(a_*(r)).
    \]
    The corresponding equations on residual arrows follow from the same
    associativity equation in \(\mathbb A\), with the residual triangle map
    left unchanged. Hence the equations hold as equations of internal
    functors.
\end{proof}

\begin{theorem}[Total residual syntax and the codomain opfibration]
\label{thm:total-residual-opfibration}
    Let
    \[
        T_A
        =
        A_1\times_{t_A,A_0,s_A}A_1
    \]
    be the object of composable pairs in \(\mathbb A\), and define
    \[
        \operatorname{Sq}_A
        :=
        T_A\times_{m_A,A_1,m_A}T_A.
    \]
    Write
    \[
        \lambda,\rho:
        \operatorname{Sq}_A\rightrightarrows T_A
    \]
    for the two pullback projections, and put
    \[
        \lambda_i:=\pi_i\lambda,
        \qquad
        \rho_i:=\pi_i\rho
        \qquad
        (i=1,2).
    \]

    A generalized element of \(\operatorname{Sq}_A\) is written
    \[
        \bigl((c,r'),(r,a)\bigr),
    \]
    where
    \[
        r:w\to u,
        \qquad
        r':w'\to v,
        \qquad
        c:w\to w',
        \qquad
        a:u\to v,
    \]
    and the defining pullback equation is
    \[
        r'\circ c=a\circ r.
    \]

    There is an internal category
    \[
        \operatorname{Res}_{\mathbb A}
        :=
        \mathbb A^\to
    \]
    whose object object is \(A_1\), whose arrow object is
    \(\operatorname{Sq}_A\), and whose source and target maps are
    \[
        s_{\operatorname{Res}}
        :=
        \rho_1:
        \operatorname{Sq}_A\to A_1,
    \]
    \[
        t_{\operatorname{Res}}
        :=
        \lambda_2:
        \operatorname{Sq}_A\to A_1.
    \]
    Thus the commutative square
    \[
        \bigl((c,r'),(r,a)\bigr)
    \]
    is an arrow
    \[
        r\longrightarrow r'
    \]
    in \(\operatorname{Res}_{\mathbb A}\).

    The identity at \(r:w\to u\) is the square
    \[
        e_{\operatorname{Res}}(r)
        =
        \bigl(
            (1_w,r),
            (r,1_u)
        \bigr).
    \]
    In morphism notation,
    \[
        e_{\operatorname{Res}}
        =
        \left\langle
            \left\langle
                e_As_A,
                1_{A_1}
            \right\rangle,
            \left\langle
                1_{A_1},
                e_At_A
            \right\rangle
        \right\rangle.
    \]

    Composition is induced by horizontal and vertical composition in
    \(\mathbb A\). Explicitly, for composable squares
    \[
        \bigl((c,r'),(r,a)\bigr):
        r\to r',
    \]
    and
    \[
        \bigl((d,r''),(r',b)\bigr):
        r'\to r'',
    \]
    their composite is
    \[
        \bigl(
            (d\circ c,r''),
            (r,b\circ a)
        \bigr).
    \]

    The maps
    \[
        \operatorname{cod}_0
        :=
        t_A:A_1\to A_0
    \]
    and
    \[
        \operatorname{cod}_1
        :=
        \rho_2:
        \operatorname{Sq}_A\to A_1
    \]
    define an internal functor
    \[
        \operatorname{cod}:
        \operatorname{Res}_{\mathbb A}
        \longrightarrow
        \mathbb A.
    \]

    This internal functor is a split internal opfibration. Its chosen
    cocartesian lift of
    \[
        a:u\to v
    \]
    at an object
    \[
        r:w\to u
    \]
    is the commutative square
    \[
        \chi(r,a)
        :=
        \bigl(
            (1_w,a\circ r),
            (r,a)
        \bigr):
        r\longrightarrow a\circ r,
    \]
    lying over \(a\).

    The vertical internal category of
    \[
        \operatorname{cod}:
        \operatorname{Res}_{\mathbb A}\to\mathbb A
    \]
    is canonically isomorphic, as an internal category in
    \(\mathcal E/A_0\), to
    \[
        \mathbb R_{\mathbb A}.
    \]
    Under this identification, the transport induced by the chosen
    cocartesian lifts is exactly the residual postcomposition functor
    \[
        \mathsf{post}:
        s_A^*\mathbb R_{\mathbb A}
        \longrightarrow
        t_A^*\mathbb R_{\mathbb A}.
    \]
    Fibrewise, the induced transport along
    \[
        a:u\to v
    \]
    is
    \[
        a_*:\mathbb A/u\to\mathbb A/v.
    \]
\end{theorem}
\begin{proof}
    \leavevmode
    \begin{enumerate}
        \item \textbf{Construct the internal arrow category.}

        By definition,
        \[
            \operatorname{Sq}_A
            =
            T_A\times_{m_A,A_1,m_A}T_A.
        \]
        Hence a generalized element
        \[
            \bigl((c,r'),(r,a)\bigr)
        \]
        consists of two composable pairs satisfying
        \[
            m_A(c,r')=m_A(r,a).
        \]
        In categorical notation this is
        \[
            r'\circ c=a\circ r,
        \]
        so the data are precisely a commutative square in
        \(\mathbb A\).

        The identity assignment is well-defined because
        \[
            r\circ1_{s_A(r)}=r
        \]
        and
        \[
            1_{t_A(r)}\circ r=r.
        \]

        To check that composition is well-defined, suppose
        \[
            r'\circ c=a\circ r
        \]
        and
        \[
            r''\circ d=b\circ r'.
        \]
        Then
        \[
            \begin{aligned}
                r''\circ(d\circ c)
                &=
                (r''\circ d)\circ c
                \\
                &=
                (b\circ r')\circ c
                \\
                &=
                b\circ(r'\circ c)
                \\
                &=
                b\circ(a\circ r)
                \\
                &=
                (b\circ a)\circ r.
            \end{aligned}
        \]
        Thus
        \[
            \bigl(
                (d\circ c,r''),
                (r,b\circ a)
            \bigr)
        \]
        is again an element of \(\operatorname{Sq}_A\).

        The source and target of the composite are \(r\) and \(r''\).
        The unit and associativity equations follow componentwise from the
        unit and associativity equations in \(\mathbb A\). Therefore
        \[
            \operatorname{Res}_{\mathbb A}
        \]
        is an internal category.

        \item \textbf{Construct the codomain functor.}

        For a commutative square
        \[
            \bigl((c,r'),(r,a)\bigr),
        \]
        the pair \((r,a)\in T_A\) gives
        \[
            t_A(r)=s_A(a).
        \]
        Moreover,
        \[
            t_A(r')
            =
            t_A(r'\circ c)
            =
            t_A(a\circ r)
            =
            t_A(a).
        \]
        Hence
        \[
            s_A\operatorname{cod}_1
            =
            \operatorname{cod}_0s_{\operatorname{Res}},
        \]
        and
        \[
            t_A\operatorname{cod}_1
            =
            \operatorname{cod}_0t_{\operatorname{Res}}.
        \]

        The bottom edge of the identity square at \(r\) is
        \(1_{t_A(r)}\), and the bottom edge of the composite of squares over
        \(a\) and \(b\) is \(b\circ a\). Therefore
        \[
            \operatorname{cod}
        \]
        preserves identities and composition.

        \item \textbf{Define the chosen cocartesian lifts.}

        Let
        \[
            r:w\to u,
            \qquad
            a:u\to v.
        \]
        Define
        \[
            \chi(r,a)
            =
            \bigl(
                (1_w,a\circ r),
                (r,a)
            \bigr).
        \]
        This lies in \(\operatorname{Sq}_A\) because
        \[
            (a\circ r)\circ1_w
            =
            a\circ r.
        \]
        Its source is \(r\), its target is \(a\circ r\), and its image under
        \(\operatorname{cod}\) is \(a\).

        Internally, the chosen-lift map is
        \[
            \chi:
            T_A\to\operatorname{Sq}_A
        \]
        given by
        \[
            \chi
            =
            \left\langle
                \left\langle
                    e_As_A\pi_1,
                    m_A
                \right\rangle,
                1_{T_A}
            \right\rangle,
        \]
        where a generalized element of \(T_A\) is written \((r,a)\).

        \item \textbf{Verify the cocartesian universal property.}

        Let
        \[
            a:u\to v,
            \qquad
            b:v\to z,
        \]
        and suppose that
        \[
            g:r\to s
        \]
        is a square lying over \(b\circ a\). Write
        \[
            g
            =
            \bigl(
                (d,s),
                (r,b\circ a)
            \bigr).
        \]
        Its commutativity equation is
        \[
            s\circ d=(b\circ a)\circ r.
        \]

        Define
        \[
            h
            :=
            \bigl(
                (d,s),
                (a\circ r,b)
            \bigr).
        \]
        This is a commutative square because
        \[
            \begin{aligned}
                s\circ d
                &=
                (b\circ a)\circ r
                \\
                &=
                b\circ(a\circ r).
            \end{aligned}
        \]
        It is an arrow
        \[
            h:a\circ r\to s
        \]
        lying over \(b\).

        By the composition formula in
        \(\operatorname{Res}_{\mathbb A}\),
        \[
            h\circ\chi(r,a)
            =
            g.
        \]

        This factorization is unique. Indeed, composition with
        \(\chi(r,a)\) leaves the top edge and the target vertical edge
        unchanged. Thus any square over \(b\) whose composite with
        \(\chi(r,a)\) is \(g\) must have top edge \(d\), target edge \(s\),
        source edge \(a\circ r\), and bottom edge \(b\), and is therefore
        equal to \(h\).

        By Notation~\ref{not:behaviour-generalized-elements}, this
        generalized-element calculation is the equality of the corresponding
        internal morphisms from the pullback object classifying such
        factorization problems. Hence \(\chi(r,a)\) is internally
        cocartesian.

        \item \textbf{Verify splitness.}

        For an identity arrow,
        \[
            \chi(r,1_u)
            =
            \bigl(
                (1_w,r),
                (r,1_u)
            \bigr)
            =
            1_r.
        \]

        For composable arrows
        \[
            a:u\to v,
            \qquad
            b:v\to z,
        \]
        one has
        \[
            \chi(a\circ r,b)\circ\chi(r,a)
            =
            \chi(r,b\circ a).
        \]
        Both sides are the square with top edge \(1_w\), source edge \(r\),
        target edge
        \[
            b\circ(a\circ r)
            =
            (b\circ a)\circ r,
        \]
        and bottom edge \(b\circ a\).

        Therefore the chosen cocartesian lifts preserve identities and
        composition strictly.

        \item \textbf{Identify the vertical category.}

        A vertical arrow is a commutative square whose bottom edge is an
        identity. Given a composable pair
        \[
            (c,r)\in T_A,
        \]
        define the vertical square
        \[
            \Phi(c,r)
            :=
            \bigl(
                (c,r),
                (r\circ c,1_{t_A(r)})
            \bigr).
        \]
        Its source object is \(r\circ c\), its target object is \(r\), and it
        lies over \(1_{t_A(r)}\).

        Conversely, let
        \[
            \bigl(
                (c,r'),
                (r,1_v)
            \bigr)
        \]
        be a vertical square. Its commutativity equation gives
        \[
            r'\circ c=1_v\circ r=r.
        \]
        It is therefore uniquely determined by the composable pair
        \[
            (c,r').
        \]

        These two assignments are mutually inverse. Under this
        identification, a vertical arrow is precisely
        \[
            r\circ c\longrightarrow r,
        \]
        its identity is
        \[
            (1_{s_A(r)},r),
        \]
        and its composition is
        \[
            (c,r)\circ(d,r\circ c)
            =
            (c\circ d,r).
        \]
        These are exactly the structure maps of
        \[
            \mathbb R_{\mathbb A}.
        \]

        Hence the vertical category of \(\operatorname{cod}\) is canonically
        isomorphic to \(\mathbb R_{\mathbb A}\) in \(\mathcal E/A_0\).

        \item \textbf{Identify the induced transport.}

        Let
        \[
            a:u\to v.
        \]
        The chosen cocartesian lift sends a residual object
        \[
            r:w\to u
        \]
        to
        \[
            a\circ r:w\to v.
        \]

        A vertical residual arrow represented by
        \[
            (c,r):
            r\circ c\longrightarrow r
        \]
        is transported to the vertical square represented by
        \[
            (c,a\circ r):
            (a\circ r)\circ c
            \longrightarrow
            a\circ r.
        \]
        By associativity,
        \[
            (a\circ r)\circ c
            =
            a\circ(r\circ c).
        \]
        Thus the induced transport is exactly
        \[
            r\longmapsto a\circ r,
            \qquad
            (c,r)\longmapsto(c,a\circ r),
        \]
        which is the internal functor
        \[
            \mathsf{post}:
            s_A^*\mathbb R_{\mathbb A}
            \longrightarrow
            t_A^*\mathbb R_{\mathbb A}.
        \]
        Fibrewise, this is
        \[
            a_*:\mathbb A/u\to\mathbb A/v.
        \]
    \end{enumerate}
\end{proof}

\subsection{The behaviour fibration}
\label{subsec:behaviour-fibration}

Behaviour is obtained from residual syntax by forming the internal functor
category from the vertical residual category into the constant internal
category \(\mathbb B\), both regarded as internal categories in the slice
\(\mathcal E/A_0\).

\begin{notation}[The vertical residual category]
\label{not:vertical-residual-category}
    Let
    \[
        \mathbb R_{\mathbb A}
    \]
    denote the internal category in \(\mathcal E/A_0\) constructed in
    Proposition~\ref{prop:residual-fibres-in-slice}. Thus its object object is
    \[
        A_1\xrightarrow{t_A}A_0,
    \]
    and its arrow object is
    \[
        T_A=A_1\times_{t_A,A_0,s_A}A_1
        \xrightarrow{t_A\pi_2}A_0.
    \]
    Its fibre over \(v\in A_0\) is the slice internal category
    \[
        \mathbb A/v.
    \]
\end{notation}

\begin{definition}[Constant target category over the residual base]
\label{def:constant-target-category-over-base}
    Let
    \[
        \Delta_{A_0}\mathbb B
    \]
    denote the constant internal category in \(\mathcal E/A_0\) whose object
    and arrow objects are
    \[
        A_0\times B_0\longrightarrow A_0,
        \qquad
        A_0\times B_1\longrightarrow A_0,
    \]
    with structure maps
    \[
        1_{A_0}\times s_B,
        \qquad
        1_{A_0}\times t_B,
        \qquad
        1_{A_0}\times e_B,
        \qquad
        1_{A_0}\times m_B.
    \]
\end{definition}

\begin{lemma}[Base change for internal functor categories]
\label{lem:base-change-internal-functor-categories}
    Let
    \[
        f:X\to Y
    \]
    be a morphism in a locally cartesian closed category \(\mathcal E\), and
    let \(\mathbb C,\mathbb D\) be internal categories in
    \[
        \mathcal E/Y.
    \]
    There is a canonical isomorphism of internal categories in
    \(\mathcal E/X\)
    \[
        f^*
        [\mathbb C,\mathbb D]_{\mathcal E/Y}
        \cong
        [f^*\mathbb C,f^*\mathbb D]_{\mathcal E/X}.
    \]
    Under this isomorphism, the pullback of the evaluation functor agrees
    with the evaluation functor of the pulled-back internal categories.
\end{lemma}
\begin{proof}
    The strict internal functor category
    \[
        [\mathbb C,\mathbb D]_{\mathcal E/Y}
    \]
    is constructed from finite limits and exponentials in the slice
    \(\mathcal E/Y\).

    Reindexing
    \[
        f^*:\mathcal E/Y\to\mathcal E/X
    \]
    preserves finite limits. Since \(\mathcal E\) is locally cartesian
    closed, reindexing also preserves slice exponentials. Concretely, this
    follows by expressing an exponential as a dependent product along a
    projection and applying the Beck--Chevalley isomorphism for the
    corresponding pullback square.

    Applying \(f^*\) to the finite-limit and exponential presentation of the
    object of internal functors therefore gives the corresponding
    presentation of
    \[
        [f^*\mathbb C,f^*\mathbb D]_{\mathcal E/X}.
    \]
    The same argument applies to the object of internal natural
    transformations and to the source, target, identity, and composition
    maps.

    The evaluation functor is induced by the evaluation maps of the
    exponentials used in the construction. These evaluation maps are
    preserved by the same base-change isomorphisms. Hence the stated
    isomorphism is compatible with evaluation.
\end{proof}

\begin{theorem}[Internal category of residual behaviours]
\label{thm:object-of-residual-behaviours}
    The strict internal functor category
    \[
        \mathbf{Beh}^{\mathrm{vert}}_{\mathbb A,\mathbb B}
        :=
        [\mathbb R_{\mathbb A},
          \Delta_{A_0}\mathbb B]_{\mathcal E/A_0}
    \]
    exists as an internal category in
    \[
        \mathcal E/A_0.
    \]

    Write its object and arrow objects as
    \[
        p_{\mathbf{Beh}}:
        \mathbf{Beh}_0\to A_0,
    \]
    and
    \[
        p_{\mathbf{Beh},1}:
        \mathbf{Beh}_1\to A_0.
    \]
    Its fibre over an object \(v\) of \(\mathbb A\) is canonically the strict
    internal functor category
    \[
        [\mathbb A/v,\mathbb B].
    \]
\end{theorem}
\begin{proof}
    Since \(\mathcal E\) is locally cartesian closed, the slice
    \[
        \mathcal E/A_0
    \]
    is cartesian closed and finitely complete.

    In any finitely complete cartesian closed category, the strict internal
    functor category between two internal categories can be constructed from
    exponentials and equalizers. Its object object is the subobject of the
    object of pairs
    \[
        (H_0,H_1)
    \]
    of maps
    \[
        H_0:
        (\mathbb R_{\mathbb A})_0
        \to
        (\Delta_{A_0}\mathbb B)_0,
    \]
    \[
        H_1:
        (\mathbb R_{\mathbb A})_1
        \to
        (\Delta_{A_0}\mathbb B)_1
    \]
    satisfying the source, target, identity, and composition equations.

    Its arrow object is the object of triples
    \[
        (H,K,\eta)
    \]
    in which \(H\) and \(K\) are such internal functors and
    \[
        \eta:H\Rightarrow K
    \]
    is an internal natural transformation. The component-typing and
    naturality equations are again imposed by finite limits.

    This constructs
    \[
        \mathbf{Beh}^{\mathrm{vert}}_{\mathbb A,\mathbb B}
    \]
    in \(\mathcal E/A_0\).

    To identify its fibres, pull the construction back along an object
    \[
        v:1\to A_0
    \]
    or, more generally, along a generalized element of \(A_0\). By
    Lemma~\ref{lem:base-change-internal-functor-categories},
    \[
        v^*
        [\mathbb R_{\mathbb A},
          \Delta_{A_0}\mathbb B]_{\mathcal E/A_0}
        \cong
        [v^*\mathbb R_{\mathbb A},
          v^*\Delta_{A_0}\mathbb B].
    \]
    The pulled-back residual category is
    \[
        v^*\mathbb R_{\mathbb A}
        =
        \mathbb A/v,
    \]
    and the pulled-back constant target category is \(\mathbb B\).
    Therefore the fibre is
    \[
        [\mathbb A/v,\mathbb B].
    \]
\end{proof}

\begin{definition}[Behaviour comparison]
\label{def:behaviour-comparison}
    Let
    \[
        H,K:\mathbb A/v\to\mathbb B
    \]
    be residual behaviours over the same object \(v\). A
    \textbf{behaviour comparison}
    \[
        \eta:H\Rightarrow K
    \]
    is an internal natural transformation.

    Its component at a residual object
    \[
        r:u\to v
    \]
    is an arrow
    \[
        \eta_r:H(r)\to K(r)
    \]
    satisfying, for every residual triangle
    \[
        w\xrightarrow{c}u\xrightarrow{r}v,
    \]
    the naturality equation
    \[
        \eta_r\circ H(c,r)
        =
        K(c,r)\circ\eta_{r\circ c}.
    \]
\end{definition}

The naturality square is
\[
    \begin{tikzcd}[column sep=large, row sep=large]
            H(r\circ c)
                \arrow[r, "{H(c,r)}"]
                \arrow[d, "{\eta_{r\circ c}}"']
        &
            H(r)
                \arrow[d, "{\eta_r}"]
        \\
            K(r\circ c)
                \arrow[r, "{K(c,r)}"']
        &
            K(r).
    \end{tikzcd}
\]

\begin{theorem}[Object of behaviour comparisons]
\label{thm:object-of-behaviour-comparisons}
    There is an object
    \[
        \mathbf{Beh}_1\to A_0
    \]
    representing triples
    \[
        (H,K,\eta),
    \]
    where \(H\) and \(K\) are residual behaviours over the same object of
    \(\mathbb A\), and
    \[
        \eta:H\Rightarrow K
    \]
    is a behaviour comparison.

    The source and target maps
    \[
        s_{\mathbf{Beh}},t_{\mathbf{Beh}}:
        \mathbf{Beh}_1\rightrightarrows\mathbf{Beh}_0
    \]
    send
    \[
        \eta:H\Rightarrow K
    \]
    to \(H\) and \(K\), respectively, and factor through
    \[
        \mathbf{Beh}_0\times_{A_0}\mathbf{Beh}_0.
    \]
\end{theorem}
\begin{proof}
    The arrow object of the internal functor category
    \[
        [\mathbb R_{\mathbb A},\Delta_{A_0}\mathbb B]_{\mathcal E/A_0}
    \]
    is the object of internal natural transformations.

    Over a pair
    \[
        (H,K)\in
        \mathbf{Beh}_0\times_{A_0}\mathbf{Beh}_0,
    \]
    a component assignment is a family of arrows of \(\mathbb B\)
    \[
        \eta_r:H(r)\to K(r)
    \]
    indexed by residual objects \(r\). Such assignments are represented by
    exponentiation in the slice. Their source and target equations, and the
    naturality equation indexed by the residual triangle object \(T_A\), are
    imposed by pullbacks and equalizers. This gives the representing object
    \(\mathbf{Beh}_1\).
\end{proof}

\begin{theorem}[Vertical behaviour category]
\label{thm:internal-behaviour-category}
    The data
    \[
        \mathbf{Beh}^{\mathrm{vert}}_{\mathbb A,\mathbb B}
        =
        \left(
            \mathbf{Beh}_0,
            \mathbf{Beh}_1,
            s_{\mathbf{Beh}},
            t_{\mathbf{Beh}},
            e_{\mathbf{Beh}},
            m_{\mathbf{Beh}}
        \right)
    \]
    form an internal category in the slice
    \[
        \mathcal E/A_0.
    \]
    In particular, every vertical behaviour comparison has source and target
    in the same fibre over \(A_0\).
\end{theorem}
\begin{proof}
    Identities and composition are those of the internal functor category
    \[
        [\mathbb R_{\mathbb A},\Delta_{A_0}\mathbb B]_{\mathcal E/A_0}.
    \]
    Fibrewise, they are identity natural transformations and vertical
    composition of natural transformations. Their laws hold pointwise in
    \(\mathbb B\), and hence hold internally in \(\mathcal E/A_0\).
\end{proof}

\begin{theorem}[Parametrized precomposition]
\label{thm:parametrized-precomposition}
    Residual postcomposition induces an internal functor in
    \(\mathcal E/A_1\)
    \[
        \mathsf{pre}_A:
        t_A^*
        \mathbf{Beh}^{\mathrm{vert}}_{\mathbb A,\mathbb B}
        \longrightarrow
        s_A^*
        \mathbf{Beh}^{\mathrm{vert}}_{\mathbb A,\mathbb B}.
    \]

    Fibrewise, over an arrow
    \[
        a:u\to v,
    \]
    this functor is precomposition with
    \[
        a_*:\mathbb A/u\to\mathbb A/v.
    \]
    Thus it sends
    \[
        H:\mathbb A/v\to\mathbb B
    \]
    to
    \[
        H\circ a_*:\mathbb A/u\to\mathbb B,
    \]
    and sends a natural transformation
    \[
        \eta:H\Rightarrow K
    \]
    to
    \[
        \eta\ast a_*:
        H\circ a_*
        \Rightarrow
        K\circ a_*.
    \]
\end{theorem}
\begin{proof}
    By Lemma~\ref{lem:base-change-internal-functor-categories}, there are
    canonical isomorphisms of internal categories in \(\mathcal E/A_1\)
    \[
        t_A^*
        \mathbf{Beh}^{\mathrm{vert}}_{\mathbb A,\mathbb B}
        \cong
        [t_A^*\mathbb R_{\mathbb A},
         \Delta_{A_1}\mathbb B]_{\mathcal E/A_1},
    \]
    and
    \[
        s_A^*
        \mathbf{Beh}^{\mathrm{vert}}_{\mathbb A,\mathbb B}
        \cong
        [s_A^*\mathbb R_{\mathbb A},
         \Delta_{A_1}\mathbb B]_{\mathcal E/A_1}.
    \]

    Residual syntax supplies an internal functor
    \[
        \mathsf{post}:
        s_A^*\mathbb R_{\mathbb A}
        \longrightarrow
        t_A^*\mathbb R_{\mathbb A}.
    \]
    Precomposition with \(\mathsf{post}\) defines an internal functor
    \[
        \mathsf{post}^*:
        [t_A^*\mathbb R_{\mathbb A},
         \Delta_{A_1}\mathbb B]_{\mathcal E/A_1}
        \longrightarrow
        [s_A^*\mathbb R_{\mathbb A},
         \Delta_{A_1}\mathbb B]_{\mathcal E/A_1}.
    \]
    On objects it sends an internal functor \(H\) to
    \[
        H\circ\mathsf{post},
    \]
    and on arrows it sends an internal natural transformation \(\eta\) to
    its whiskering
    \[
        \eta\ast\mathsf{post}.
    \]

    Transporting this internal functor across the two canonical
    base-change isomorphisms defines
    \[
        \mathsf{pre}_A.
    \]
    Equivalently, \(\mathsf{pre}_A\) is the unique internal functor whose
    pullback along any generalized arrow of \(\mathbb A\) classifies
    precomposition by the corresponding residual postcomposition functor.

    Over a generalized arrow
    \[
        a:u\to v,
    \]
    the pullback of \(\mathsf{post}\) is
    \[
        a_*:\mathbb A/u\to\mathbb A/v.
    \]
    Consequently, the pullback of \(\mathsf{pre}_A\) sends
    \[
        H\longmapsto H\circ a_*,
    \]
    and
    \[
        \eta\longmapsto\eta\ast a_*,
    \]
    as claimed.

    Since internal precomposition is defined by the representing property
    of the strict internal functor category, equal residual transport
    functors induce equal maps on both the object and arrow objects of the
    behaviour category. This observation will be used in
    Proposition~\ref{prop:derivative-laws-behaviour-category}.
\end{proof}

\begin{definition}[Behaviour indexed category]
\label{def:behaviour-indexed-category}
    The \textbf{behaviour indexed category} is the vertical behaviour
    category
    \[
        \mathbf{Beh}^{\mathrm{vert}}_{\mathbb A,\mathbb B}
        \longrightarrow A_0
    \]
    together with the parametrized reindexing functor
    \[
        \mathsf{pre}_A:
        t_A^*
        \mathbf{Beh}^{\mathrm{vert}}_{\mathbb A,\mathbb B}
        \longrightarrow
        s_A^*
        \mathbf{Beh}^{\mathrm{vert}}_{\mathbb A,\mathbb B}.
    \]

    The object map of \(\mathsf{pre}_A\), followed by projection from the
    pullback over \(s_A\), is denoted
    \[
        \partial_0:
        A_1
        \times_{t_A,A_0,p_{\mathbf{Beh}}}
        \mathbf{Beh}_0
        \longrightarrow
        \mathbf{Beh}_0.
    \]
    The arrow map is denoted
    \[
        \partial_1:
        A_1
        \times_{t_A,A_0,p_{\mathbf{Beh}}s_{\mathbf{Beh}}}
        \mathbf{Beh}_1
        \longrightarrow
        \mathbf{Beh}_1.
    \]
    We also write
    \[
        \partial:=\partial_0.
    \]

    These maps satisfy
    \[
        p_{\mathbf{Beh}}\partial_0
        =
        s_A\pi_{A_1},
    \]
    and
    \[
        s_{\mathbf{Beh}}\partial_1
        =
        \partial_0
        (1_{A_1}\times s_{\mathbf{Beh}}),
    \]
    \[
        t_{\mathbf{Beh}}\partial_1
        =
        \partial_0
        (1_{A_1}\times t_{\mathbf{Beh}}).
    \]
\end{definition}

\begin{definition}[Derivative]
\label{def:derivative-behaviour}
    Let
    \[
        a:u\to v
    \]
    and let
    \[
        H:\mathbb A/v\to\mathbb B
    \]
    be a residual behaviour. The \textbf{derivative} of \(H\) by \(a\) is
    \[
        \partial_aH
        :=
        H\circ a_*.
    \]
    Thus, for a residual object
    \[
        r:w\to u,
    \]
    \[
        (\partial_aH)(r)
        =
        H(a\circ r).
    \]
\end{definition}

\begin{definition}[Derivative of a comparison]
\label{def:derivative-comparison}
    Let
    \[
        \eta:H\Rightarrow K
    \]
    be a behaviour comparison over \(v\). Its \textbf{derivative} by
    \(a:u\to v\) is
    \[
        \partial_a\eta
        :=
        \eta\ast a_*:
        \partial_aH\Rightarrow\partial_aK.
    \]
\end{definition}

\begin{proposition}[Split fibration laws for derivatives]
\label{prop:derivative-laws-behaviour-category}
    Derivatives satisfy
    \[
        \partial_{1_v}H=H,
    \]
    and, for composable arrows
    \[
        a:u\to v,
        \qquad
        b:v\to w,
    \]
    \[
        \partial_{b\circ a}H
        =
        \partial_a(\partial_bH).
    \]
    The same equations hold for behaviour comparisons.

    Derivative reindexing also preserves vertical identities and vertical
    composition:
    \[
        \partial_a(1_H)
        =
        1_{\partial_aH},
    \]
    and
    \[
        \partial_a(\mu\circ\eta)
        =
        \partial_a\mu\circ\partial_a\eta.
    \]
\end{proposition}
\begin{proof}
    Residual transport satisfies the strict internal-functor equations
    \[
        (1_v)_*
        =
        1_{\mathbb A/v},
    \]
    and
    \[
        (b\circ a)_*
        =
        b_*\circ a_*.
    \]
    Therefore, on residual behaviours,
    \[
        \begin{aligned}
            \partial_{1_v}H
            &=
            H\circ(1_v)_*
            \\
            &=
            H\circ1_{\mathbb A/v}
            \\
            &=
            H,
        \end{aligned}
    \]
    and
    \[
        \begin{aligned}
            \partial_{b\circ a}H
            &=
            H\circ(b\circ a)_*
            \\
            &=
            H\circ(b_*\circ a_*)
            \\
            &=
            (H\circ b_*)\circ a_*
            \\
            &=
            \partial_a(\partial_bH).
        \end{aligned}
    \]

    These are equalities of internal morphisms, rather than merely
    fibrewise canonical isomorphisms. Indeed, the two sides classify the
    same parametrized internal functor on residual objects and residual
    arrows. Since the object of residual behaviours represents such
    parametrized internal functors, equality of the classified functors
    gives equality of the corresponding maps into
    \(\mathbf{Beh}_0\).

    The same argument applies to behaviour comparisons. Explicitly,
    \[
        \partial_{1_v}\eta
        =
        \eta\ast(1_v)_*
        =
        \eta,
    \]
    and
    \[
        \begin{aligned}
            \partial_{b\circ a}\eta
            &=
            \eta\ast(b\circ a)_*
            \\
            &=
            \eta\ast(b_*\circ a_*)
            \\
            &=
            (\eta\ast b_*)\ast a_*
            \\
            &=
            \partial_a(\partial_b\eta).
        \end{aligned}
    \]
    Representability of the object of internal natural transformations
    makes these equalities of maps into \(\mathbf{Beh}_1\).

    Finally, \(\mathsf{pre}_A\) is an internal functor. It therefore
    preserves identity arrows and vertical composition:
    \[
        \partial_a(1_H)
        =
        1_{\partial_aH},
    \]
    and
    \[
        \partial_a(\mu\circ\eta)
        =
        \partial_a\mu\circ\partial_a\eta.
    \]
\end{proof}

\begin{theorem}[Internal evaluation of residual behaviours]
\label{thm:internal-evaluation-residual-behaviours}
    The internal functor category
    \[
        \mathbf{Beh}^{\mathrm{vert}}_{\mathbb A,\mathbb B}
        =
        [\mathbb R_{\mathbb A},
         \Delta_{A_0}\mathbb B]_{\mathcal E/A_0}
    \]
    carries a canonical internal evaluation functor
    \[
        \operatorname{ev}:
        \mathbf{Beh}^{\mathrm{vert}}_{\mathbb A,\mathbb B}
        \times_{A_0}
        \mathbb R_{\mathbb A}
        \longrightarrow
        \Delta_{A_0}\mathbb B.
    \]

    Write its object and arrow maps, after suppressing the evident
    \(A_0\)-coordinates, as
    \[
        \operatorname{ev}_0:
        \mathbf{Beh}_0
        \times_{A_0}
        A_1
        \longrightarrow
        B_0,
    \]
    and
    \[
        \operatorname{ev}_1:
        \mathbf{Beh}_1
        \times_{A_0}
        T_A
        \longrightarrow
        B_1.
    \]

    In generalized-element notation,
    \[
        \operatorname{ev}_0(H,r)=H(r).
    \]
    If
    \[
        \eta:H\Rightarrow K
    \]
    and
    \[
        \xi:r\to r'
    \]
    are arrows in the two factors, then
    \[
        \operatorname{ev}_1(\eta,\xi):
        H(r)\to K(r')
    \]
    is the common composite
    \[
        K(\xi)\circ\eta_r
        =
        \eta_{r'}\circ H(\xi).
    \]
\end{theorem}
\begin{proof}
    This is the evaluation functor associated to the internal hom
    \[
        [\mathbb R_{\mathbb A},
         \Delta_{A_0}\mathbb B]_{\mathcal E/A_0}.
    \]
    Its object map is exponential evaluation. Its arrow map is induced by
    evaluation of the component of an internal natural transformation
    followed by evaluation of the relevant residual arrow.

    The equality
    \[
        K(\xi)\circ\eta_r
        =
        \eta_{r'}\circ H(\xi)
    \]
    is precisely the internal naturality equation for \(\eta\). Preservation
    of identities and composition is the standard evaluation-functor
    calculation in the internal functor category.
\end{proof}

\begin{definition}[Anchor evaluation]
\label{def:anchor-evaluation}
    The identity residual section is
    \[
        e_A:A_0\to A_1.
    \]
    It selects
    \[
        1_v:v\to v
    \]
    in the residual fibre over \(v\).

    The \textbf{anchor map} is
    \[
        q_{\mathbf{Beh}}
        :=
        \operatorname{ev}_0
        \circ
        \left\langle
            1_{\mathbf{Beh}_0},
            e_Ap_{\mathbf{Beh}}
        \right\rangle:
        \mathbf{Beh}_0\to B_0.
    \]
    Thus, for
    \[
        H:\mathbb A/v\to\mathbb B,
    \]
    \[
        q_{\mathbf{Beh}}(H)
        =
        H(1_v).
    \]
\end{definition}

Thus the object of residual behaviours carries the anchor span
\[
    \begin{tikzcd}[column sep=large]
        &
            \mathbf{Beh}_0
                \arrow[dl, "p_{\mathbf{Beh}}"']
                \arrow[dr, "q_{\mathbf{Beh}}"]
        &
        \\
            A_0
        &
        &
            B_0.
    \end{tikzcd}
\]

\begin{definition}[Total behaviour category]
\label{def:total-behaviour-category}
    Put
    \[
        X_{\mathbf{Beh}}
        :=
        A_1\times_{t_A,A_0,p_{\mathbf{Beh}}}\mathbf{Beh}_0.
    \]
    A generalized element is a pair
    \[
        (a,H),
        \qquad
        a:u\to v,
        \qquad
        H:\mathbb A/v\to\mathbb B.
    \]

    Define the total arrow object
    \[
        \mathbf{Beh}^{\mathrm{tot}}_1
        :=
        X_{\mathbf{Beh}}
        \times_{\partial,\mathbf{Beh}_0,t_{\mathbf{Beh}}}
        \mathbf{Beh}_1.
    \]
    Thus a generalized total arrow is a triple
    \[
        (a,H,\eta)
    \]
    where
    \[
        \eta:K\Rightarrow\partial_aH
    \]
    is a vertical behaviour comparison over \(u\).

    Its source and target are
    \[
        s_{\mathrm{tot}}(a,H,\eta)
        :=
        s_{\mathbf{Beh}}(\eta)=K,
    \]
    \[
        t_{\mathrm{tot}}(a,H,\eta)
        :=
        H.
    \]
\end{definition}

\begin{theorem}[The total behaviour fibration]
\label{thm:total-behaviour-fibration}
    The object
    \[
        \mathbf{Beh}_0
    \]
    and the arrow object
    \[
        \mathbf{Beh}^{\mathrm{tot}}_1
    \]
    of Definition~\ref{def:total-behaviour-category} form an internal
    category
    \[
        \mathbf{Beh}^{\mathrm{tot}}_{\mathbb A,\mathbb B}.
    \]

    The maps
    \[
        p_{\mathbf{Beh}}:
        \mathbf{Beh}_0\to A_0
    \]
    and
    \[
        p_{\mathbf{Beh},1}^{\mathrm{tot}}:
        \mathbf{Beh}^{\mathrm{tot}}_1\to A_1,
        \qquad
        (a,H,\eta)\longmapsto a,
    \]
    define an internal functor
    \[
        \pi_{\mathbf{Beh}}:
        \mathbf{Beh}^{\mathrm{tot}}_{\mathbb A,\mathbb B}
        \to
        \mathbb A.
    \]

    This internal functor is a split internal fibration. Its chosen
    cartesian lift of
    \[
        a:u\to v
    \]
    at a behaviour \(H\) over \(v\) is
    \[
        \chi(a,H)
        :=
        (a,H,1_{\partial_aH}):
        \partial_aH\longrightarrow H.
    \]
\end{theorem}
\begin{proof}
    \leavevmode
    \begin{enumerate}
        \item \textbf{Define identities.}

        Let \(H\) lie over \(v\). Define
        \[
            1_H^{\mathrm{tot}}
            :=
            (1_v,H,1_H).
        \]
        This is well-defined because
        \[
            \partial_{1_v}H=H.
        \]

        \item \textbf{Define composition.}

        Suppose
        \[
            (a,H,\eta):K\to H
        \]
        lies over
        \[
            a:u\to v,
        \]
        so that
        \[
            \eta:K\Rightarrow\partial_aH,
        \]
        and suppose
        \[
            (b,L,\mu):H\to L
        \]
        lies over
        \[
            b:v\to w,
        \]
        so that
        \[
            \mu:H\Rightarrow\partial_bL.
        \]
        Define
        \[
            (b,L,\mu)\circ(a,H,\eta)
            :=
            \left(
                b\circ a,
                L,
                (\partial_a\mu)\circ\eta
            \right).
        \]
        This is well-typed because
        \[
            \partial_a\mu:
            \partial_aH
            \Rightarrow
            \partial_a\partial_bL
            =
            \partial_{b\circ a}L.
        \]

        \item \textbf{Verify the category laws.}

        The left and right unit equations follow from
        \[
            \partial_{1_v}=1
        \]
        and preservation of identity comparisons by derivative
        reindexing.

        For associativity, let a third total arrow be given. Expanding both
        parenthesizations reduces their comparison components to the same
        composite by:
        \[
            \partial_a(\nu\circ\mu)
            =
            \partial_a\nu\circ\partial_a\mu,
        \]
        \[
            \partial_a\partial_b
            =
            \partial_{b\circ a},
        \]
        and associativity of vertical composition in
        \(\mathbf{Beh}^{\mathrm{vert}}_{\mathbb A,\mathbb B}\).

        Hence the displayed data form an internal category.

        \item \textbf{Verify that the projection is an internal functor.}

        The identity total arrow at \(H\) lies over \(1_{p_{\mathbf{Beh}}H}\).
        The composite of arrows over \(a\) and \(b\) lies over \(b\circ a\).
        Therefore
        \[
            \pi_{\mathbf{Beh}}
        \]
        preserves identities and composition.

        \item \textbf{Define the cleavage.}

        Put
        \[
            X_{\mathbf{Beh}}
            :=
            A_1
            \times_{t_A,A_0,p_{\mathbf{Beh}}}
            \mathbf{Beh}_0.
        \]
        Define
        \[
            \chi:
            X_{\mathbf{Beh}}
            \to
            \mathbf{Beh}^{\mathrm{tot}}_1
        \]
        by
        \[
            \chi(a,H)
            =
            (a,H,1_{\partial_aH}).
        \]
        Its target is \(H\), its source is \(\partial_aH\), and its image
        under the projection functor is \(a\).

        \item \textbf{Form the internal object of composite-factorization
        problems.}

        Let
        \[
            A_2
            :=
            A_1\times_{t_A,A_0,s_A}A_1.
        \]
        We write a generalized element as
        \[
            (c,a),
        \]
        where
        \[
            c:w\to u,
            \qquad
            a:u\to v.
        \]
        Thus
        \[
            m_A(c,a)=a\circ c.
        \]

        Let \(\mathcal L\) be the pullback
        \[
            \mathcal L
            :=
            A_2
            \times_{m_A,A_1,p_{\mathbf{Beh},1}^{\mathrm{tot}}}
            \mathbf{Beh}^{\mathrm{tot}}_1.
        \]
        A generalized element of \(\mathcal L\) is a triple
        \[
            (c,a,g),
        \]
        where \(g:K\to H\) is a total arrow lying over
        \[
            a\circ c.
        \]

        Next put
        \[
            Y
            :=
            A_2
            \times_{t_A\pi_2,A_0,p_{\mathbf{Beh}}}
            \mathbf{Beh}_0.
        \]
        A generalized element of \(Y\) is
        \[
            (c,a,H)
        \]
        with \(H\) lying over \(t_A(a)\).

        There is a map
        \[
            \vartheta:
            Y
            \to
            \mathbf{Beh}_0
            \times_{p_{\mathbf{Beh}},A_0,t_A}
            A_1
        \]
        defined by
        \[
            \vartheta(c,a,H)
            =
            (\partial_aH,c).
        \]
        This is well-defined because
        \[
            p_{\mathbf{Beh}}(\partial_aH)
            =
            s_A(a)
            =
            t_A(c).
        \]

        Define
        \[
            \mathcal R
            :=
            Y
            \times_{\vartheta,\,
                    \mathbf{Beh}_0
                    \times_{A_0}A_1,\,
                    (t_{\mathrm{tot}},
                     p_{\mathbf{Beh},1}^{\mathrm{tot}})}
            \mathbf{Beh}^{\mathrm{tot}}_1.
        \]
        A generalized element of \(\mathcal R\) is a quadruple
        \[
            (c,a,H,h),
        \]
        where
        \[
            h:K\to\partial_aH
        \]
        is a total arrow lying over \(c\).

        \item \textbf{Construct the cartesian-factorization isomorphism.}

        Composition with the chosen lift defines a map
        \[
            \Gamma:
            \mathcal R\to\mathcal L
        \]
        by
        \[
            \Gamma(c,a,H,h)
            =
            \bigl(
                c,a,
                \chi(a,H)\circ h
            \bigr).
        \]

        To compute this map, write
        \[
            h=(c,\partial_aH,\nu),
        \]
        where
        \[
            \nu:
            K\Rightarrow
            \partial_c\partial_aH
            =
            \partial_{a\circ c}H.
        \]
        Since
        \[
            \chi(a,H)
            =
            (a,H,1_{\partial_aH}),
        \]
        the composition formula gives
        \[
            \chi(a,H)\circ h
            =
            (a\circ c,H,\nu).
        \]

        Conversely, let
        \[
            g=(a\circ c,H,\nu)
        \]
        represent an element of \(\mathcal L\). Define
        \[
            \Gamma^{-1}(c,a,g)
            :=
            \left(
                c,a,H,
                (c,\partial_aH,\nu)
            \right).
        \]
        This is well-defined because
        \[
            \partial_c\partial_aH
            =
            \partial_{a\circ c}H.
        \]

        The two constructions leave the vertical comparison \(\nu\)
        unchanged. Therefore
        \[
            \Gamma^{-1}\Gamma=1_{\mathcal R},
            \qquad
            \Gamma\Gamma^{-1}=1_{\mathcal L}.
        \]
        Hence
        \[
            \Gamma:\mathcal R\overset{\cong}{\longrightarrow}\mathcal L
        \]
        is an isomorphism.

        This is the internal cartesian universal property: every total arrow
        \[
            g:K\to H
        \]
        lying over \(a\circ c\) factors uniquely as
        \[
            K
            \xrightarrow{h}
            \partial_aH
            \xrightarrow{\chi(a,H)}
            H,
        \]
        where \(h\) lies over \(c\).

        \item \textbf{Verify splitness.}

        For an identity arrow,
        \[
            \chi(1_v,H)
            =
            (1_v,H,1_H)
            =
            1_H^{\mathrm{tot}}.
        \]

        For composable
        \[
            c:w\to u,
            \qquad
            a:u\to v,
        \]
        one has
        \[
            \chi(a\circ c,H)
            =
            \chi(a,H)
            \circ
            \chi(c,\partial_aH).
        \]
        Indeed, both sides are the total arrow
        \[
            \partial_{a\circ c}H\to H
        \]
        over \(a\circ c\) with identity vertical comparison, using
        \[
            \partial_c\partial_aH
            =
            \partial_{a\circ c}H.
        \]

        Thus the chosen cartesian lifts are strictly compatible with
        identities and composition, and
        \[
            \pi_{\mathbf{Beh}}
        \]
        is a split internal fibration.
    \end{enumerate}
\end{proof}

\begin{remark}
\label{rem:vertical-not-total-fibration}
    The internal category
    \[
        \mathbf{Beh}^{\mathrm{vert}}_{\mathbb A,\mathbb B}
    \]
    contains only comparisons inside a fixed residual fibre. The total
    category
    \[
        \mathbf{Beh}^{\mathrm{tot}}_{\mathbb A,\mathbb B}
    \]
    also contains arrows over non-identity arrows of \(\mathbb A\).

    An arrow over
    \[
        a:u\to v
    \]
    from a behaviour \(K\) over \(u\) to a behaviour \(H\) over \(v\) is a
    vertical comparison
    \[
        K\Rightarrow\partial_aH.
    \]
    The canonical behaviour Mealy morphism uses the chosen cartesian
    cleavage of this total fibration.
\end{remark}

\begin{remark}[Equality of behaviours]
\label{rem:equality-of-behaviours}
    The object \(\mathbf{Beh}_0\) represents residual functors themselves. Thus equality in \(\mathbf{Beh}_0\) is equality of represented residual functors, not isomorphism of residual functors and not equivalence inside the vertical behaviour category.
\end{remark}

\section{The cleavage Mealy morphism}
\label{sec:cleavage-mealy-morphism}

The behaviour indexed category has a split cleavage given by derivative reindexing. Anchor evaluation turns this cleavage into a Mealy morphism.

\subsection{Internal derivative and anchor-output maps}
\label{subsec:internal-derivative-anchor-output}

\begin{lemma}[Internal derivative and anchor-output maps]
\label{lem:internal-derivative-output-maps}
    Put
    \[
        X_{\mathbf{Beh}}
        :=
        A_1
        \times_{t_A,A_0,p_{\mathbf{Beh}}}
        \mathbf{Beh}_0,
    \]
    and write its projections as
    \[
        \pi_A:X_{\mathbf{Beh}}\to A_1,
        \qquad
        \pi_H:X_{\mathbf{Beh}}\to\mathbf{Beh}_0.
    \]

    The derivative map
    \[
        \partial:
        X_{\mathbf{Beh}}\to\mathbf{Beh}_0
    \]
    satisfies
    \[
        p_{\mathbf{Beh}}\partial
        =
        s_A\pi_A.
    \]

    Define
    \[
        \sigma:
        A_1\to T_A
    \]
    by
    \[
        \sigma
        :=
        \left\langle
            1_{A_1},
            e_At_A
        \right\rangle.
    \]
    Then \(\sigma(a)\) represents the residual arrow
    \[
        a\longrightarrow1_{t_Aa}.
    \]

    The anchor-output map is
    \[
        \mathrm{out}_{\mathbf{Beh}}
        :=
        \operatorname{ev}_1
        \circ
        \left\langle
            e_{\mathbf{Beh}}\pi_H,
            \sigma\pi_A
        \right\rangle:
        X_{\mathbf{Beh}}\to B_1.
    \]
    It satisfies
    \[
        s_B\mathrm{out}_{\mathbf{Beh}}
        =
        q_{\mathbf{Beh}}\partial,
    \]
    and
    \[
        t_B\mathrm{out}_{\mathbf{Beh}}
        =
        q_{\mathbf{Beh}}\pi_H.
    \]
\end{lemma}
\begin{proof}
    The derivative equation
    \[
        p_{\mathbf{Beh}}\partial
        =
        s_A\pi_A
    \]
    is the base equation for parametrized precomposition.

    To define \(\sigma\), observe that
    \[
        t_A
        =
        s_Ae_At_A.
    \]
    Hence the pair
    \[
        \left(
            1_{A_1},
            e_At_A
        \right)
    \]
    induces a map into
    \[
        T_A
        =
        A_1\times_{t_A,A_0,s_A}A_1.
    \]
    For an arrow \(a:u\to v\),
    \[
        \sigma(a)=(a,1_v).
    \]
    Its source and target in the residual category are
    \[
        1_v\circ a=a,
        \qquad
        1_v,
    \]
    so it represents
    \[
        a\to1_v.
    \]

    The pair
    \[
        \left(
            e_{\mathbf{Beh}}\pi_H,
            \sigma\pi_A
        \right)
    \]
    is a map into the arrow object of
    \[
        \mathbf{Beh}^{\mathrm{vert}}_{\mathbb A,\mathbb B}
        \times_{A_0}
        \mathbb R_{\mathbb A}.
    \]
    Its first component is the identity natural transformation on \(H\),
    while its second component is the residual arrow
    \[
        a\to1_v.
    \]
    Applying the arrow part of evaluation gives
    \[
        \mathrm{out}_{\mathbf{Beh}}(a,H)
        =
        H(a\to1_v):
        H(a)\to H(1_v).
    \]

    The target is therefore
    \[
        H(1_v)
        =
        q_{\mathbf{Beh}}(H),
    \]
    proving
    \[
        t_B\mathrm{out}_{\mathbf{Beh}}
        =
        q_{\mathbf{Beh}}\pi_H.
    \]

    For the source, let \(a:u\to v\). By definition of derivative,
    \[
        (\partial_aH)(1_u)
        =
        H(a_*1_u).
    \]
    Residual postcomposition gives
    \[
        a_*1_u
        =
        a\circ1_u
        =
        a.
    \]
    Consequently,
    \[
        H(a)
        =
        (\partial_aH)(1_u)
        =
        q_{\mathbf{Beh}}(\partial_aH).
    \]
    Hence
    \[
        s_B\mathrm{out}_{\mathbf{Beh}}
        =
        q_{\mathbf{Beh}}\partial.
    \]
\end{proof}

\subsection{The canonical behaviour carrier}
\label{subsec:canonical-behaviour-carrier}

\begin{definition}[Canonical behaviour carrier]
\label{def:canonical-behaviour-carrier}
    The \textbf{canonical behaviour carrier} is the anchor span
    \[
        A_0
        \xleftarrow{p_{\mathbf{Beh}}}
        \mathbf{Beh}_0
        \xrightarrow{q_{\mathbf{Beh}}}
        B_0.
    \]
\end{definition}

\begin{definition}[Canonical transition and output]
\label{def:canonical-transition-output}
    The \textbf{canonical transition} is derivative reindexing:
    \[
        \delta_{\mathbf{Beh}}:=\partial.
    \]
    The \textbf{canonical output} is the anchor comparison:
    \[
        \omega_{\mathbf{Beh}}:=\mathrm{out}_{\mathbf{Beh}}.
    \]
    In generalized-element notation,
    \[
        \delta_{\mathbf{Beh}}(a,H)=\partial_aH,
    \]
    and
    \[
        \omega_{\mathbf{Beh}}(a,H)=H(a,1_v):H(a)\to H(1_v),
    \]
    where \(a:u\to v\) and \(H:\mathbb A/v\to\mathbb B\).
\end{definition}

The defining residual triangle for the output is
\[
    \begin{tikzcd}[column sep=large, row sep=large]
            u
                \arrow[r, "a"]
                \arrow[rr, "a", bend left=15]
        &
            v
                \arrow[r, "1_v"']
        &
            v.
    \end{tikzcd}
\]

\subsection{The cleavage Mealy morphism}
\label{subsec:cleavage-mealy-morphism}

\begin{theorem}[Cleavage Mealy morphism of the behaviour fibration]
\label{thm:canonical-behaviour-mealy-morphism}
    The anchor span, together with derivative transition and anchor-comparison output, defines a Mealy morphism
    \[
        \mathbf{Beh}^{\mathrm{can}}_{\mathbb A,\mathbb B}:
        \mathbb A\rightsquigarrow\mathbb B.
    \]
\end{theorem}
\begin{proof}
    \leavevmode
    \begin{enumerate}
        \item \textbf{Check typing.} The typing equations are exactly the two parts of Lemma~\ref{lem:internal-derivative-output-maps}.

        \item \textbf{Verify the transition laws.} The transition unit and multiplication laws are the split fibration laws:
        \[
            \partial_{1_v}H=H,
            \qquad
            \partial_{b\circ a}H=\partial_a(\partial_bH).
        \]

        \item \textbf{Verify the output unit law.} The output unit law follows because the residual morphism
        \[
            1_v\to 1_v
        \]
        is the identity in the slice \(\mathbb A/v\), so applying \(H\) gives
        \[
            1_{H(1_v)}.
        \]

        \item \textbf{Factor the residual composite.} For multiplication, let
        \[
            a:u\to v,
            \qquad
            b:v\to w,
            \qquad
            H:\mathbb A/w\to\mathbb B.
        \]
        In the slice \(\mathbb A/w\), the residual morphism
        \[
            b\circ a\to 1_w
        \]
        factors as
        \[
            \begin{tikzcd}[column sep=large, row sep=large]
                    u
                        \arrow[r, "a"]
                        \arrow[rr, "b\circ a", bend left=17]
                        \arrow[rrr, "b\circ a", bend left=32]
                &
                    v
                        \arrow[r, "b"]
                        \arrow[rr, "b", bend left=17]
                &
                    w
                        \arrow[r, "1_w"']
                &
                    w.
            \end{tikzcd}
        \]
        That is, the residual arrow
        \[
            b\circ a\to 1_w
        \]
        factors through
        \[
            b\circ a\longrightarrow b\longrightarrow 1_w.
        \]

        \item \textbf{Apply functoriality of \(H\).} Applying the functor \(H\) gives
        \[
            H(b\circ a,1_w)
            =
            H(b,1_w)\circ H(a,b).
        \]
        This is the required Mealy output law:
        \[
            \omega_{\mathbf{Beh}}(b\circ a,H)
            =
            \omega_{\mathbf{Beh}}(b,H)
            \circ
            \omega_{\mathbf{Beh}}(a,\partial_bH).
        \]
    \end{enumerate}
\end{proof}

\subsection{Derivative-closed subobjects}
\label{subsec:derivative-closed-subobjects}

Put
\[
    Z:=A_0\times B_0.
\]
The behaviour object is regarded as an object of \(\mathcal E/Z\) by
\[
    (p_{\mathbf{Beh}},q_{\mathbf{Beh}}):
    \mathbf{Beh}_0\to Z.
\]

\begin{definition}[Derivative domain of a subobject]
\label{def:derivative-domain-subobject}
    Let
    \[
        i:S\hookrightarrow\mathbf{Beh}_0
    \]
    be a subobject in \(\mathcal E/Z\). Write
    \[
        p_S:=p_{\mathbf{Beh}}i.
    \]
    The \textbf{derivative domain} of \(S\) is
    \[
        \mathsf D(S):=
        A_1\times_{t_A,A_0,p_S}S.
    \]
    It is regarded as an object of \(\mathcal E/Z\) through the derivative composite
    \[
        \mathsf D(S)
            \xrightarrow{\partial(1_{A_1}\times i)}
        \mathbf{Beh}_0
            \xrightarrow{(p_{\mathbf{Beh}},q_{\mathbf{Beh}})}
        Z.
    \]
\end{definition}

The defining pullback is
\[
    \begin{tikzcd}[column sep=large, row sep=large]
            \mathsf D(S)
                \arrow[r, "\pi_S"]
                \arrow[d, "\pi_{A_1}"']
                \arrow[dr, phantom, "\lrcorner", very near start]
        &
            S
                \arrow[d, "p_S"]
        \\
            A_1
                \arrow[r, "t_A"']
        &
            A_0.
    \end{tikzcd}
\]
Thus a generalized element \((a,H)\), with
\[
    a:u\to p_S(H),
\]
has \(Z\)-coordinate
\[
    \bigl(u,\ q_{\mathbf{Beh}}(\partial_a iH)\bigr).
\]
This coordinate is the input interface and anchor of the derivative behaviour.

\begin{definition}[Derivative-closed subobject]
\label{def:derivative-closed-subobject}
    A subobject
    \[
        i:S\hookrightarrow\mathbf{Beh}_0
    \]
    in \(\mathcal E/Z\) is \textbf{derivative-closed} if the derivative map
    \[
        \partial(1_{A_1}\times i):
        \mathsf D(S)\to\mathbf{Beh}_0
    \]
    factors through \(i\) in \(\mathcal E/Z\). Equivalently, there is a unique
    map
    \[
        \bar\partial_S:\mathsf D(S)\to S
    \]
    over \(Z\) such that
    \[
        i\bar\partial_S
        =
        \partial(1_{A_1}\times i).
    \]
\end{definition}

The factorization condition is displayed by
\[
    \begin{tikzcd}[column sep=large, row sep=large]
            \mathsf D(S)
                \arrow[rr, "{\partial(1_{A_1}\times i)}"]
                \arrow[dr, "\bar\partial_S"']
        &
        &
            \mathbf{Beh}_0
        \\
        &
            S
                \arrow[ur, "i"']
        &
    \end{tikzcd}
    \qquad
    \text{in }\mathcal E/Z.
\]

\begin{lemma}[Restriction to derivative-closed subobjects]
\label{lem:restriction-derivative-closed-subobject}
    Let
    \[
        i:S\hookrightarrow\mathbf{Beh}_0
    \]
    be derivative-closed. Then the canonical behaviour Mealy morphism
    restricts to a Mealy morphism carried by
    \[
        A_0\xleftarrow{p_{\mathbf{Beh}}i}S
        \xrightarrow{q_{\mathbf{Beh}}i}B_0.
    \]
\end{lemma}
\begin{proof}
    Write
    \[
        p_S:=p_{\mathbf{Beh}}i,
        \qquad
        q_S:=q_{\mathbf{Beh}}i.
    \]
    Derivative-closedness supplies a unique map
    \[
        \bar\partial_S:
        A_1\times_{t_A,A_0,p_S}S
        \to
        S
    \]
    satisfying
    \[
        i\bar\partial_S
        =
        \partial(1_{A_1}\times i).
    \]
    Define the restricted transition by
    \[
        \delta_S:=\bar\partial_S.
    \]

    Define the restricted output by
    \[
        \omega_S
        :=
        \omega_{\mathbf{Beh}}
        \circ
        (1_{A_1}\times i).
    \]
    Thus, in generalized-element notation,
    \[
        \delta_S(a,s)=\bar\partial_S(a,s),
    \]
    and
    \[
        \omega_S(a,s)
        =
        \omega_{\mathbf{Beh}}(a,i(s)).
    \]

    We first verify the typing equations. Since
    \(\bar\partial_S\) is a morphism over the derivative-induced
    \(Z\)-structure,
    \[
        p_S(\delta_S(a,s))=s_A(a).
    \]
    Moreover,
    \[
        \begin{aligned}
            s_B\omega_S(a,s)
            &=
            s_B\omega_{\mathbf{Beh}}(a,i(s))
            \\
            &=
            q_{\mathbf{Beh}}(\partial_a i(s))
            \\
            &=
            q_{\mathbf{Beh}}(i(\delta_S(a,s)))
            \\
            &=
            q_S(\delta_S(a,s)),
        \end{aligned}
    \]
    while
    \[
        \begin{aligned}
            t_B\omega_S(a,s)
            &=
            t_B\omega_{\mathbf{Beh}}(a,i(s))
            \\
            &=
            q_{\mathbf{Beh}}(i(s))
            \\
            &=
            q_S(s).
        \end{aligned}
    \]

    For the transition unit law, compose with the monomorphism \(i\):
    \[
        \begin{aligned}
            i(\delta_S(1_{p_S(s)},s))
            &=
            \partial_{1_{p_S(s)}}i(s)
            \\
            &=
            i(s).
        \end{aligned}
    \]
    Hence
    \[
        \delta_S(1_{p_S(s)},s)=s.
    \]

    For the transition multiplication law, let
    \[
        a:u\to v,
        \qquad
        b:v\to w,
        \qquad
        p_S(s)=w.
    \]
    After composing with \(i\), one has
    \[
        \begin{aligned}
            i(\delta_S(b\circ a,s))
            &=
            \partial_{b\circ a}i(s)
            \\
            &=
            \partial_a(\partial_bi(s))
            \\
            &=
            \partial_a(i(\delta_S(b,s)))
            \\
            &=
            i(\delta_S(a,\delta_S(b,s))).
        \end{aligned}
    \]
    Since \(i\) is monic,
    \[
        \delta_S(b\circ a,s)
        =
        \delta_S(a,\delta_S(b,s)).
    \]

    The output unit law follows directly by restricting the ambient output
    unit law:
    \[
        \begin{aligned}
            \omega_S(1_{p_S(s)},s)
            &=
            \omega_{\mathbf{Beh}}(1_{p_S(s)},i(s))
            \\
            &=
            1_{q_{\mathbf{Beh}}(i(s))}
            \\
            &=
            1_{q_S(s)}.
        \end{aligned}
    \]

    For the output multiplication law,
    \[
        \begin{aligned}
            \omega_S(b\circ a,s)
            &=
            \omega_{\mathbf{Beh}}(b\circ a,i(s))
            \\
            &=
            \omega_{\mathbf{Beh}}(b,i(s))
            \circ
            \omega_{\mathbf{Beh}}(a,\partial_bi(s))
            \\
            &=
            \omega_{\mathbf{Beh}}(b,i(s))
            \circ
            \omega_{\mathbf{Beh}}
            (a,i(\delta_S(b,s)))
            \\
            &=
            \omega_S(b,s)
            \circ
            \omega_S(a,\delta_S(b,s)).
        \end{aligned}
    \]
    Thus the restricted transition and output satisfy all the Mealy laws.
\end{proof}

\section{Mealy actions and residual semantics}
\label{sec:mealy-actions-residual-semantics}

An arbitrary Mealy morphism determines a residual semantics by sending each state to its residual orbit. This construction is the categorical analogue of sending a state of a deterministic automaton to the language, or output behaviour, observed from that state. Related formulations appear in the coalgebraic account of final semantics and minimization \cite{Rutten1998,Rutten2000}.

\subsection{The execution category}
\label{subsec:execution-category}

The transition part of a Mealy morphism is a contravariant action of \(\mathbb A\), equivalently a covariant action of \(\mathbb A^{\mathrm{op}}\). We use the corresponding category of executions, whose arrows are oriented in the direction in which states evolve.

\begin{definition}[Execution category]
\label{def:execution-category}
    Let
    \[
        P:\mathbb A\rightsquigarrow\mathbb B
    \]
    be a Mealy morphism. Its \textbf{execution category}
    \[
        \int_{\mathbb A}P
    \]
    has object object \(P\) and arrow object
    \[
        A_1\times_{t_A,A_0,p}P.
    \]
    The arrow
    \[
        (a,x)
    \]
    has source and target
    \[
        s_{\int P}(a,x)=x,
        \qquad
        t_{\int P}(a,x)=\delta(a,x).
    \]
    The identity at \(x\) is
    \[
        (1_{p(x)},x).
    \]
    If
    \[
        (a,x):x\to\delta(a,x),
        \qquad
        (c,\delta(a,x)):
        \delta(a,x)\to\delta(c,\delta(a,x))
    \]
    are composable, where
    \[
        c:w\to u,
        \qquad
        a:u\to v,
    \]
    then their composite is
    \[
        (a\circ c,x):x\to\delta(a\circ c,x).
    \]
    Here \(a\circ c\) denotes the ordinary categorical composite
    \[
        w\xrightarrow{c}u\xrightarrow{a}v,
    \]
    so, in the notation of Chapter~\ref{ch:spans-internal-categories-mealy},
    \[
        a\circ c=m_A(c,a).
    \]
\end{definition}

The transition square for an execution arrow is
\[
    \begin{tikzcd}[column sep=large, row sep=large]
            {q(\delta(a,x))}
                \arrow[r, "{\omega(a,x)}"]
        &
            {q(x)}
            \\
            {\delta(a,x)}
                \arrow[u, "q"]
        &
            {x}
                \arrow[l, "{(a,x)}"']
                \arrow[u, "q"']
    \end{tikzcd}
\]
where the horizontal arrow on the bottom is in the execution category and the top arrow is its emitted \(\mathbb B\)-output.

\begin{proposition}[Execution category]
\label{prop:execution-category}
    The data of Definition~\ref{def:execution-category} define an internal category
    \[
        \int_{\mathbb A}P.
    \]
\end{proposition}
\begin{proof}
    \leavevmode
    \begin{enumerate}
        \item \textbf{Use the Mealy laws for \(\delta\).} The identity and composition maps are well-defined by the Mealy unit and multiplication laws for \(\delta\). The identity and associativity laws are inherited from the identity and associativity laws in \(\mathbb A\), using the Mealy laws for \(\delta\).

        \item \textbf{Display the execution order.} The execution category composes in the ordinary categorical order of \(\mathbb A\), while execution itself processes inputs target-to-source. The composite execution square is
        \[
            \begin{tikzcd}[column sep=large, row sep=large]
                    x
                        \arrow[r, "{(a,x)}"]
                        \arrow[rr, "{(a\circ c,x)}"', bend right=18]
                &
                    \delta(a,x)
                        \arrow[r, "{(c,\delta(a,x))}"]
                &
                    \delta(c,\delta(a,x)).
            \end{tikzcd}
        \]
    \end{enumerate}
\end{proof}

\subsection{The output functor}
\label{subsec:output-functor}

\begin{proposition}[Output functor]
\label{prop:output-functor}
    The output map of \(P\) is equivalently an internal functor
    \[
        \Omega_P:\int_{\mathbb A}P\to\mathbb B^{\mathrm{op}}
    \]
    with object map
    \[
        q:P\to B_0.
    \]
    It sends an execution arrow
    \[
        (a,x):x\to\delta(a,x)
    \]
    to
    \[
        \omega(a,x)^{\mathrm{op}}:q(x)\to q(\delta(a,x))
    \]
    in \(\mathbb B^{\mathrm{op}}\).
\end{proposition}
\begin{proof}
    \leavevmode
    \begin{enumerate}
        \item \textbf{Check arrow typing.} The Mealy typing equations give
        \[
            \omega(a,x):q(\delta(a,x))\to q(x)
        \]
        in \(\mathbb B\), hence an arrow
        \[
            q(x)\to q(\delta(a,x))
        \]
        in \(\mathbb B^{\mathrm{op}}\).

        \item \textbf{Check identities.} The Mealy output unit law gives identity preservation.

        \item \textbf{Check composition.} For composition, consider composable execution arrows
        \[
            (a,x):x\to\delta(a,x),
            \qquad
            (c,\delta(a,x)):\delta(a,x)\to\delta(c,\delta(a,x)).
        \]
        Their composite is
        \[
            (a\circ c,x).
        \]
        The output of the composite is
        \[
            \omega(a\circ c,x)
            =
            \omega(a,x)\circ\omega(c,\delta(a,x))
        \]
        in \(\mathbb B\). Passing to \(\mathbb B^{\mathrm{op}}\), this is exactly composition preservation:
        \[
            \omega(c,\delta(a,x))^{\mathrm{op}}
            \circ
            \omega(a,x)^{\mathrm{op}}
            =
            \omega(a\circ c,x)^{\mathrm{op}}.
        \]

        \item \textbf{Display the opposite composition square.} The relevant square in \(\mathbb B^{\mathrm{op}}\) is
        \[
            \begin{tikzcd}[column sep=large, row sep=large]
                    {q(x)}
                        \arrow{r}{\omega(a,x)^{\mathrm{op}}}
                        \arrow{rr}[swap, bend right=15]{\omega(a\circ c,x)^{\mathrm{op}}}
                &
                    {q(\delta(a,x))}
                        \arrow{r}{\omega(c,\delta(a,x))^{\mathrm{op}}}
                &
                    {q(\delta(c,\delta(a,x)))}
            \end{tikzcd}
        \]
    \end{enumerate}
\end{proof}

\subsection{Residual orbit functors}
\label{subsec:orbit-functors}

\begin{definition}[Residual orbit]
\label{def:residual-orbit}
    Let \(x\in P\) with \(p(x)=v\). The \textbf{residual orbit} of \(x\) is the functor
    \[
        \rho_x:\mathbb A/v\to
        \left(\int_{\mathbb A}P\right)^{\mathrm{op}}
    \]
    defined as follows. On a residual object
    \[
        r:u\to v,
    \]
    it is
    \[
        \rho_x(r)=\delta(r,x).
    \]
    On a residual morphism represented by
    \[
        w\xrightarrow{a}u\xrightarrow{r}v,
    \]
    it is the opposite of the execution arrow
    \[
        (a,\delta(r,x)):
        \delta(r,x)\to\delta(a,\delta(r,x)).
    \]
    Since
    \[
        \delta(a,\delta(r,x))=\delta(r\circ a,x),
    \]
    this gives an arrow
    \[
        \delta(r\circ a,x)\to\delta(r,x)
    \]
    in
    \[
        \left(\int_{\mathbb A}P\right)^{\mathrm{op}}.
    \]
\end{definition}

The orbit sends a residual triangle to an opposite execution arrow:
\[
    \begin{tikzcd}[column sep=large, row sep=large]
            w
                \arrow[r, "a"]
                \arrow[rr, "r\circ a", bend left=17]
        &
            u
                \arrow[r, "r"]
        &
            v
    \end{tikzcd}
    \qquad\leadsto\qquad
    \begin{tikzcd}[column sep=large, row sep=large]
            \delta(r\circ a,x)
                \arrow[r]
        &
            \delta(r,x)
    \end{tikzcd}
\]
where the arrow on the right is in \(\left(\int_{\mathbb A}P\right)^{\mathrm{op}}\).

\begin{proposition}[Orbit functor]
\label{prop:residual-orbit-functor}
    For every state \(x\), the residual orbit assignment is a functor
    \[
        \rho_x:\mathbb A/p(x)\to
        \left(\int_{\mathbb A}P\right)^{\mathrm{op}}.
    \]
    Moreover, for \(a:u\to v\) and \(p(x)=v\),
    \[
        \rho_{\delta(a,x)}=\rho_x\circ a_*.
    \]
\end{proposition}
\begin{proof}
    \leavevmode
    \begin{enumerate}
        \item \textbf{Prove functoriality.} Functoriality follows from the identity and composition laws of the execution category.

        \item \textbf{Identify reindexing of orbits.} The displayed compatibility states that executing from the state \(\delta(a,x)\) along a residual \(r\) is the same as executing from \(x\) along the transported residual \(a\circ r\). This is the Mealy multiplication law for \(\delta\).

        \item \textbf{Display the compatibility square.} The compatibility is summarized by
        \[
            \begin{tikzcd}[column sep=huge, row sep=large]
                    \mathbb A/u
                        \arrow[r, "a_*"]
                        \arrow[d, "\rho_{\delta(a,x)}"']
                &
                    \mathbb A/v
                        \arrow[d, "\rho_x"]
                \\
                    \left(\int_{\mathbb A}P\right)^{\mathrm{op}}
                        \arrow[r, equal]
                &
                    \left(\int_{\mathbb A}P\right)^{\mathrm{op}}.
            \end{tikzcd}
        \]
    \end{enumerate}
\end{proof}

\subsection{Residual behaviour of a state}
\label{subsec:residual-behaviour-of-state}

\begin{definition}[Residual behaviour of a state]
\label{def:residual-behaviour-of-state}
    The \textbf{residual behaviour} of a state \(x\in P\) is the composite
    \[
        H_x:=
        \Omega_P^{\mathrm{op}}\circ\rho_x:
        \mathbb A/p(x)\to\mathbb B.
    \]
    Thus
    \[
        H_x(r)=q(\delta(r,x)).
    \]
    For a residual triangle
    \[
        w\xrightarrow{a}u\xrightarrow{r}p(x),
    \]
    the arrow value is
    \[
        H_x(a,r)
        =
        \omega(a,\delta(r,x)):
        q(\delta(r\circ a,x))\to q(\delta(r,x)).
    \]
\end{definition}

The residual behaviour is the composite of the orbit and output functor:
\[
    \begin{tikzcd}[column sep=huge, row sep=large]
            \mathbb A/p(x)
                \arrow[r, "\rho_x"]
                \arrow[rr, "H_x"', bend right=16]
        &
            \left(\int_{\mathbb A}P\right)^{\mathrm{op}}
                \arrow[r, "\Omega_P^{\mathrm{op}}"]
        &
            \mathbb B.
    \end{tikzcd}
\]

\begin{proposition}
\label{prop:state-residual-behaviour-functor}
    For every \(x\in P\), the assignment \(H_x\) is a residual behaviour. Moreover,
    \[
        H_x(1_{p(x)})=q(x).
    \]
\end{proposition}
\begin{proof}
    \leavevmode
    \begin{enumerate}
        \item \textbf{Use functoriality of the composite.} The first claim follows because \(H_x\) is the composite of two functors.

        \item \textbf{Evaluate at the anchor.} The second follows from the Mealy unit law:
        \[
            H_x(1_{p(x)})
            =
            q(\delta(1_{p(x)},x))
            =
            q(x).
        \]
    \end{enumerate}
\end{proof}

\subsection{The behaviour assignment}
\label{subsec:semantic-map-beta}

\begin{definition}[Behaviour assignment]
\label{def:behaviour-assignment-map}
    For a Mealy morphism \(P\), the \textbf{behaviour assignment} is the map
    \[
        \beta_P:P\to\mathbf{Beh}_0
    \]
    classified by the residual behaviours
    \[
        x\longmapsto H_x.
    \]
\end{definition}

\begin{lemma}[Internal construction of \(\beta_P\)]
\label{lem:internal-construction-beta}
    The assignment
    \[
        x\longmapsto H_x
    \]
    defines a morphism
    \[
        \beta_P:P\to\mathbf{Beh}_0.
    \]
\end{lemma}
\begin{proof}
    Pull back the vertical residual category along
    \[
        p:P\to A_0.
    \]
    This gives an internal category
    \[
        p^*\mathbb R_{\mathbb A}
    \]
    in the slice
    \[
        \mathcal E/P.
    \]

    Its object object is
    \[
        P\times_{p,A_0,t_A}A_1
        \longrightarrow
        P,
    \]
    whose generalized elements are pairs
    \[
        (x,r),
        \qquad
        r:u\to p(x).
    \]

    Its arrow object is
    \[
        P\times_{p,A_0,t_A\pi_2}T_A
        \longrightarrow
        P,
    \]
    whose generalized elements are triples
    \[
        (x,a,r),
        \qquad
        w\xrightarrow{a}u\xrightarrow{r}p(x).
    \]
    Such a triple represents the residual arrow
    \[
        r\circ a\longrightarrow r.
    \]

    Put
    \[
        \Delta_P\mathbb B
        :=
        p^*\Delta_{A_0}\mathbb B.
    \]
    This is the constant internal category \(\mathbb B\) in
    \(\mathcal E/P\). Its object and arrow objects are
    \[
        P\times B_0\longrightarrow P,
        \qquad
        P\times B_1\longrightarrow P.
    \]

    Define an object map in \(\mathcal E/P\)
    \[
        \widetilde H_{P,0}:
        P\times_{p,A_0,t_A}A_1
        \longrightarrow
        P\times B_0
    \]
    by
    \[
        \widetilde H_{P,0}(x,r)
        :=
        \bigl(
            x,
            q(\delta(r,x))
        \bigr).
    \]

    Define an arrow map in \(\mathcal E/P\)
    \[
        \widetilde H_{P,1}:
        P\times_{p,A_0,t_A\pi_2}T_A
        \longrightarrow
        P\times B_1
    \]
    by
    \[
        \widetilde H_{P,1}(x,a,r)
        :=
        \bigl(
            x,
            \omega(a,\delta(r,x))
        \bigr).
    \]

    Both maps preserve the displayed \(P\)-coordinate and are therefore
    morphisms in \(\mathcal E/P\).

    We verify that they define an internal functor
    \[
        \widetilde H_P:
        p^*\mathbb R_{\mathbb A}
        \longrightarrow
        \Delta_P\mathbb B.
    \]

    First, the source of the arrow
    \[
        \omega(a,\delta(r,x))
    \]
    is
    \[
        \begin{aligned}
            s_B\omega(a,\delta(r,x))
            &=
            q(\delta(a,\delta(r,x)))
            \\
            &=
            q(\delta(r\circ a,x)),
        \end{aligned}
    \]
    using the Mealy multiplication law for \(\delta\). This is the
    \(B_0\)-coordinate of
    \[
        \widetilde H_{P,0}(x,r\circ a).
    \]

    Its target is
    \[
        t_B\omega(a,\delta(r,x))
        =
        q(\delta(r,x)),
    \]
    which is the \(B_0\)-coordinate of
    \[
        \widetilde H_{P,0}(x,r).
    \]

    Hence the source and target equations hold in \(\mathcal E/P\).

    Preservation of identities follows from the Mealy output unit law. The
    identity residual arrow at
    \[
        r:u\to p(x)
    \]
    is represented by
    \[
        (1_u,r).
    \]
    Therefore
    \[
        \begin{aligned}
            \widetilde H_{P,1}(x,1_u,r)
            &=
            \bigl(
                x,
                \omega(1_u,\delta(r,x))
            \bigr)
            \\
            &=
            \bigl(
                x,
                1_{q(\delta(r,x))}
            \bigr),
        \end{aligned}
    \]
    which is the identity at
    \[
        \widetilde H_{P,0}(x,r).
    \]

    To verify preservation of composition, consider a residual string
    \[
        z\xrightarrow{b}w\xrightarrow{a}u
        \xrightarrow{r}p(x).
    \]
    The two residual arrows are
    \[
        r\circ a\circ b\longrightarrow r\circ a
    \]
    and
    \[
        r\circ a\longrightarrow r,
    \]
    and their composite is represented by
    \[
        (a\circ b,r):
        r\circ(a\circ b)\longrightarrow r.
    \]

    The Mealy output multiplication law gives
    \[
        \begin{aligned}
            \omega(a\circ b,\delta(r,x))
            &=
            \omega(a,\delta(r,x))
            \circ
            \omega(b,\delta(a,\delta(r,x)))
            \\
            &=
            \omega(a,\delta(r,x))
            \circ
            \omega(b,\delta(r\circ a,x)).
        \end{aligned}
    \]
    Hence
    \[
        \widetilde H_{P,1}(x,a\circ b,r)
        =
        \widetilde H_{P,1}(x,a,r)
        \circ
        \widetilde H_{P,1}(x,b,r\circ a).
    \]

    Thus
    \[
        \widetilde H_P:
        p^*\mathbb R_{\mathbb A}
        \longrightarrow
        \Delta_P\mathbb B
    \]
    is an internal functor in \(\mathcal E/P\).

    By Lemma~\ref{lem:base-change-internal-functor-categories}, there is a
    canonical isomorphism of internal categories in \(\mathcal E/P\)
    \[
        p^*
        \mathbf{Beh}^{\mathrm{vert}}_{\mathbb A,\mathbb B}
        \cong
        [p^*\mathbb R_{\mathbb A},
         \Delta_P\mathbb B]_{\mathcal E/P}.
    \]
    The internal functor \(\widetilde H_P\) therefore determines a section of
    the object object
    \[
        p^*\mathbf{Beh}_0\to P.
    \]
    Equivalently, it determines a unique morphism
    \[
        \beta_P:P\to\mathbf{Beh}_0
    \]
    over \(A_0\).

    Fibrewise, this morphism sends a state \(x\) to the residual behaviour
    \[
        H_x:\mathbb A/p(x)\to\mathbb B
    \]
    whose object and arrow assignments are
    \[
        H_x(r)=q(\delta(r,x)),
    \]
    and
    \[
        H_x(a,r)=\omega(a,\delta(r,x)).
    \]
\end{proof}

\begin{proposition}[Semantic homomorphism]
\label{prop:behaviour-assignment}
    The map
    \[
        \beta_P:P\to\mathbf{Beh}_0
    \]
    lies over
    \[
        Z=A_0\times B_0.
    \]
    It satisfies
    \[
        \beta_P(\delta_P(a,x))
        =
        \partial_a(\beta_P(x)),
    \]
    and
    \[
        \omega_P(a,x)
        =
        \omega_{\mathbf{Beh}}(a,\beta_P(x)).
    \]
    Therefore
    \[
        \beta_P:P\to\mathbf{Beh}^{\mathrm{can}}_{\mathbb A,\mathbb B}
    \]
    is a strict semantic Mealy homomorphism.
\end{proposition}
\begin{proof}
    \leavevmode
    \begin{enumerate}
        \item \textbf{Check the map over \(Z\).} Since \(H_x\) lies over \(p(x)\) and satisfies
        \[
            H_x(1_{p(x)})=q(x),
        \]
        the map \(\beta_P\) lies over \(Z\). This is the commutativity of
        \[
            \begin{tikzcd}[column sep=large, row sep=large]
                    P
                        \arrow[r, "\beta_P"]
                        \arrow[dr, "{(p,q)}"']
                &
                    \mathbf{Beh}_0
                        \arrow[d, "{(p_{\mathbf{Beh}},q_{\mathbf{Beh}})}"]
                \\
                &
                    A_0\times B_0.
            \end{tikzcd}
        \]

        \item \textbf{Check transition preservation.} For \(a:u\to v\) and \(p(x)=v\), the residual-orbit identity gives
        \[
            \rho_{\delta(a,x)}=\rho_x\circ a_*.
        \]
        Therefore
        \[
            \beta_P(\delta(a,x))
            =
            H_{\delta(a,x)}
            =
            H_x\circ a_*
            =
            \partial_aH_x
            =
            \partial_a(\beta_P(x)).
        \]

        \item \textbf{Check output preservation.} Finally,
        \[
            \omega_{\mathbf{Beh}}(a,\beta_P(x))
            =
            H_x(a,1_v)
            =
            \omega_P(a,\delta_P(1_v,x))
            =
            \omega_P(a,x).
        \]
    \end{enumerate}
\end{proof}

\section{Terminal residual semantics}
\label{sec:terminal-residual-semantics}

The canonical behaviour machine is terminal among strict semantic Mealy homomorphisms. The terminality is formulated in an ordinary category whose morphisms preserve transition and output strictly. This section is in the same general line as universal-realization and final-semantics formulations of automata behaviour \cite{Goguen1973,Rutten1998,Rutten2000}.

\subsection{Strict semantic homomorphisms}
\label{subsec:strict-semantic-homomorphisms}

\begin{definition}[Strict semantic Mealy homomorphism]
\label{def:strict-semantic-mealy-homomorphism}
    Let
    \[
        P,Q:\mathbb A\rightsquigarrow\mathbb B
    \]
    be Mealy morphisms. A \textbf{strict semantic Mealy homomorphism}
    \[
        f:P\to Q
    \]
    is a map over
    \[
        Z=A_0\times B_0,
    \]
    so that
    \[
        p_Qf=p_P,
        \qquad
        q_Qf=q_P,
    \]
    satisfying
    \[
        f(\delta_P(a,x))=\delta_Q(a,f(x)),
    \]
    and
    \[
        \omega_P(a,x)=\omega_Q(a,f(x)).
    \]
\end{definition}

The strict preservation equations are displayed by
\[
    \begin{tikzcd}[column sep=large, row sep=large]
        x
            \arrow[r, "{a,\omega_P(a,x)}"]
            \arrow[d, "f"']
        &
            \delta_P(a,x)
                \arrow[d, "f"]
        \\
            f(x)
                \arrow[r, "{a,\omega_Q(a,f(x))}"']
        &
            \delta_Q(a,f(x)).
    \end{tikzcd}
\]

\begin{lemma}
\label{lem:strict-semantic-category}
    Strict semantic Mealy homomorphisms are closed under identities and composition.
\end{lemma}
\begin{proof}
    \leavevmode
    \begin{enumerate}
        \item \textbf{Check identities and composition.} This is immediate from the defining preservation equations.
    \end{enumerate}
\end{proof}

\begin{notation}
\label{not:strict-semantic-category}
    Write
    \[
        \mathbf{Mly}^{\mathrm{str}}_{\mathbb A,\mathbb B}
    \]
    for the ordinary category whose objects are Mealy morphisms
    \[
        \mathbb A\rightsquigarrow\mathbb B
    \]
    and whose morphisms are strict semantic Mealy homomorphisms.
\end{notation}

\begin{remark}
\label{rem:terminality-not-bicategorical}
    The terminality statement below is a statement in the ordinary category
    \[
        \mathbf{Mly}^{\mathrm{str}}_{\mathbb A,\mathbb B}.
    \]
    The bicategory \(\mathbf{Mly}(\mathcal E)\) uses Mealy simulations as 2-cells, and those simulations have oppositely directed state maps.
\end{remark}

\subsection{Terminality of the canonical behaviour machine}
\label{subsec:terminality-canonical-behaviour-machine}

\begin{lemma}
\label{lem:canonical-behaviour-map-identity}
    For the canonical behaviour Mealy morphism,
    \[
        \beta_{\mathbf{Beh}^{\mathrm{can}}}=1_{\mathbf{Beh}_0}.
    \]
    If
    \[
        i:S\hookrightarrow\mathbf{Beh}_0
    \]
    is derivative-closed and the canonical behaviour machine is restricted to \(S\), then its behaviour assignment is the inclusion
    \[
        S\hookrightarrow\mathbf{Beh}_0.
    \]
\end{lemma}
\begin{proof}
    \leavevmode
    \begin{enumerate}
        \item \textbf{Compute the residual behaviour of a canonical state.} Let \(H:\mathbb A/v\to\mathbb B\) be a state of the canonical behaviour machine. Its residual behaviour sends
        \[
            r:u\to v
        \]
        to
        \[
            q_{\mathbf{Beh}}(\partial_rH)=H(r),
        \]
        and sends a residual triangle
        \[
            w\xrightarrow{a}u\xrightarrow{r}v
        \]
        to
        \[
            \omega_{\mathbf{Beh}}(a,\partial_rH)=H(a,r).
        \]

        \item \textbf{Conclude identity of the behaviour map.} Thus the residual behaviour determined by \(H\) is \(H\). The restricted statement follows by the same calculation after composing with the subobject inclusion.
    \end{enumerate}
\end{proof}

\begin{theorem}[Terminal residual semantics]
\label{thm:terminal-residual-semantics}
    The canonical behaviour Mealy morphism
    \[
        \mathbf{Beh}^{\mathrm{can}}_{\mathbb A,\mathbb B}
    \]
    is terminal in
    \[
        \mathbf{Mly}^{\mathrm{str}}_{\mathbb A,\mathbb B}.
    \]
    For every Mealy morphism \(P\), the unique strict semantic homomorphism
    \[
        P\to\mathbf{Beh}^{\mathrm{can}}_{\mathbb A,\mathbb B}
    \]
    is
    \[
        \beta_P:x\mapsto H_x.
    \]
\end{theorem}
\begin{proof}
    \leavevmode
    \begin{enumerate}
        \item \textbf{Establish existence.} Existence is Proposition~\ref{prop:behaviour-assignment}.

        \item \textbf{Set up uniqueness.} For uniqueness, let
        \[
            f:P\to\mathbf{Beh}_0
        \]
        be a strict semantic Mealy homomorphism. We show \(f=\beta_P\).

        \item \textbf{Recover the object assignment.} Let \(x\in P\), and let \(r:u\to p(x)\) be a residual input. Since \(f\) preserves transitions,
        \[
            f(\delta_P(r,x))=\partial_r f(x).
        \]
        Since \(f\) lies over \(Z\),
        \[
            q_{\mathbf{Beh}}(f(\delta_P(r,x)))=q(\delta_P(r,x)).
        \]
        But
        \[
            q_{\mathbf{Beh}}(\partial_r f(x))=(f(x))(r).
        \]
        Hence
        \[
            (f(x))(r)=q(\delta_P(r,x))=H_x(r).
        \]

        \item \textbf{Recover the arrow assignment.} Now let
        \[
            w\xrightarrow{a}u\xrightarrow{r}p(x)
        \]
        be a residual triangle. Strict output preservation at the state
        \[
            \delta_P(r,x)
        \]
        gives
        \[
            \omega_P(a,\delta_P(r,x))
            =
            \omega_{\mathbf{Beh}}(a,f(\delta_P(r,x))).
        \]
        Using transition preservation,
        \[
            f(\delta_P(r,x))=\partial_r f(x).
        \]
        Therefore
        \[
            \omega_{\mathbf{Beh}}(a,\partial_r f(x))
            =
            (\partial_r f(x))(a,1_u)
            =
            (f(x))(a,r).
        \]
        Hence
        \[
            (f(x))(a,r)
            =
            \omega_P(a,\delta_P(r,x))
            =
            H_x(a,r).
        \]
        Thus \(f(x)=H_x\) for every \(x\), so \(f=\beta_P\).

        \item \textbf{Display the reconstruction diagrams.} The proof reconstructs the object and arrow assignments of \(f(x)\) by the two diagrams
        \[
            \begin{tikzcd}[column sep=large, row sep=large]
                    \delta_P(r,x)
                        \arrow[r, "f"]
                        \arrow[d, "q"']
                &
                    \partial_r f(x)
                        \arrow[d, "q_{\mathbf{Beh}}"]
                \\
                    q(\delta_P(r,x))
                        \arrow[r, equal]
                &
                    (f(x))(r)
            \end{tikzcd}
            \qquad
            \begin{tikzcd}[column sep=large, row sep=large]
                    q(\delta_P(r\circ a,x))
                        \arrow[r, "{\omega_P(a,\delta_P(r,x))}"]
                        \arrow[d, equal]
                &
                    q(\delta_P(r,x))
                        \arrow[d, equal]
                \\
                    (f(x))(r\circ a)
                        \arrow[r, "{(f(x))(a,r)}"']
                &
                    (f(x))(r).
            \end{tikzcd}
        \]
    \end{enumerate}
\end{proof}

\subsection{Behavioural separatedness}
\label{subsec:behavioural-separatedness}

\begin{definition}[Behavioural separatedness]
\label{def:behavioural-separatedness}
    A Mealy morphism
    \[
        P:\mathbb A\rightsquigarrow\mathbb B
    \]
    is \textbf{behaviourally separated} if
    \[
        \beta_P:P\to\mathbf{Beh}_0
    \]
    is monic.
\end{definition}

\section{Simulations and induced behaviour comparisons}
\label{sec:simulations-and-behaviour-comparisons}

Simulations compare machines by output-comparison arrows. Under residual semantics, every Mealy simulation induces a vertical natural transformation between the corresponding residual behaviours.

\subsection{Simulation-induced vertical behaviour comparisons}
\label{subsec:simulation-induced-behaviour-comparisons}

\begin{theorem}[Simulation-induced vertical behaviour comparison]
\label{thm:simulation-induces-behaviour-comparison}
    Let
    \[
        (\theta,\alpha):P\Rightarrow Q
    \]
    be a Mealy simulation. For each \(y\in Q\), the family
    \[
        (\widehat\alpha_y)_r:=\alpha_{\delta_Q(r,y)}
    \]
    defines a natural transformation
    \[
        \widehat\alpha_y:
        H^P_{\theta y}\Rightarrow H^Q_y.
    \]
    Equivalently, the simulation induces a morphism
    \[
        Q\to\mathbf{Beh}_1
    \]
    whose source is
    \[
        \beta_P\theta
    \]
    and whose target is
    \[
        \beta_Q.
    \]
\end{theorem}
\begin{proof}
    \leavevmode
    \begin{enumerate}
        \item \textbf{Type the components.} The state law gives
        \[
            \theta(\delta_Q(r,y))=\delta_P(r,\theta y).
        \]
        Hence
        \[
            \alpha_{\delta_Q(r,y)}
            :
            q_P(\delta_P(r,\theta y))
            \to
            q_Q(\delta_Q(r,y)),
        \]
        which is exactly a component
        \[
            H^P_{\theta y}(r)\to H^Q_y(r).
        \]

        \item \textbf{Check naturality.} For a residual triangle
        \[
            w\xrightarrow{a}u\xrightarrow{r}p_Q(y),
        \]
        naturality requires
        \[
            (\widehat\alpha_y)_r\circ H^P_{\theta y}(a,r)
            =
            H^Q_y(a,r)\circ(\widehat\alpha_y)_{r\circ a}.
        \]
        This is the simulation output law applied to the state
        \[
            \delta_Q(r,y).
        \]

        \item \textbf{Display the naturality square.} The square is
        \[
            \begin{tikzcd}[column sep=large, row sep=large]
                    H^P_{\theta y}(r\circ a)
                        \arrow[r, "{H^P_{\theta y}(a,r)}"]
                        \arrow[d, "{\alpha_{\delta_Q(r\circ a,y)}}"']
                &
                    H^P_{\theta y}(r)
                        \arrow[d, "{\alpha_{\delta_Q(r,y)}}"]
                \\
                    H^Q_y(r\circ a)
                        \arrow[r, "{H^Q_y(a,r)}"']
                &
                    H^Q_y(r).
            \end{tikzcd}
        \]
    \end{enumerate}
\end{proof}

\begin{remark}
\label{rem:simulation-to-comparison-not-equivalence}
    The construction in Theorem~\ref{thm:simulation-induces-behaviour-comparison} is a semantic consequence of a state-level simulation. A vertical behaviour comparison compares the residual behaviours after applying \(\beta_P\) and \(\beta_Q\). A Mealy simulation also contains the state-level comparison data and satisfies the transition law before residual semantics is applied.
\end{remark}

\section{Specifications and derivative images}
\label{sec:specifications-and-derivative-images}

A specification is a family of behaviours. Its formal derivative orbit is obtained by acting on that family by all residual input arrows through the cleavage of the behaviour fibration.

\subsection{Specifications}
\label{subsec:specifications}

\begin{definition}[Behaviour specification]
\label{def:behaviour-specification}
    A \textbf{behaviour specification} is a morphism
    \[
        k:K\to\mathbf{Beh}_0.
    \]
    Write
    \[
        p_K:=p_{\mathbf{Beh}}k.
    \]
    We regard \(K\) as an object of
    \[
        \mathcal E/Z,
        \qquad
        Z=A_0\times B_0,
    \]
    through the composite
    \[
        K
            \xrightarrow{k}
        \mathbf{Beh}_0
            \xrightarrow{(p_{\mathbf{Beh}},q_{\mathbf{Beh}})}
        Z.
    \]
\end{definition}

\subsection{Formal derivative orbits}
\label{subsec:formal-derivative-orbits}

\begin{definition}[Formal derivative orbit]
\label{def:formal-derivative-orbit}
    The \textbf{formal derivative orbit object} of \(K\) is
    \[
        D_K:=A_1\times_{t_A,A_0,p_K}K.
    \]
    A generalized element is a pair
    \[
        (a,\kappa),
    \]
    where
    \[
        a:u\to p_K(\kappa).
    \]
\end{definition}

Its defining pullback is
\[
    \begin{tikzcd}[column sep=large, row sep=large]
            D_K
                \arrow[r, "\pi_K"]
                \arrow[d, "\pi_{A_1}"']
                \arrow[dr, phantom, "\lrcorner", very near start]
        &
            K
                \arrow[d, "p_K"]
        \\
            A_1
                \arrow[r, "t_A"']
        &
            A_0.
    \end{tikzcd}
\]

\begin{definition}[Residual evaluation]
\label{def:residual-evaluation}
    The \textbf{residual evaluation map} is
    \[
        d_K:D_K\to\mathbf{Beh}_0
    \]
    defined by
    \[
        d_K(a,\kappa)=\partial_a k(\kappa).
    \]
    Equivalently,
    \[
        d_K=
        \partial\circ(1_{A_1}\times k).
    \]
\end{definition}

The object \(D_K\) is regarded as an object over
\[
    Z=A_0\times B_0
\]
through the composite
\[
    D_K
        \xrightarrow{d_K}
    \mathbf{Beh}_0
        \xrightarrow{(p_{\mathbf{Beh}},q_{\mathbf{Beh}})}
    Z.
\]
Thus the \(Z\)-coordinate of \((a,\kappa)\) is the input interface and anchor of the derivative behaviour
\[
    \partial_a k(\kappa).
\]
We write
\[
    p_{D_K}:=p_{\mathbf{Beh}}d_K.
\]
For a generalized element \((a,\kappa)\), with
\[
    a:u\to p_K(\kappa),
\]
one has
\[
    p_{D_K}(a,\kappa)=u=s_A(a).
\]

\begin{definition}[Unit inclusion of syntax]
\label{def:syntax-unit-inclusion}
    The \textbf{unit inclusion of syntax}
    \[
        \iota_K:K\to D_K
    \]
    is
    \[
        \kappa\longmapsto(1_{p_K(\kappa)},\kappa).
    \]
    It satisfies
    \[
        d_K\iota_K=k.
    \]
\end{definition}

\subsection{Derivative images}
\label{subsec:derivative-image}

\begin{lemma}[Regular-epi descent for subobjects]
\label{lem:regular-epi-descent-subobjects}
    Let \(\mathcal E\) be regular. If
    \[
        e:X\twoheadrightarrow Y
    \]
    is a regular epimorphism and
    \[
        m:S\hookrightarrow Z
    \]
    is a subobject, then a map
    \[
        g:Y\to Z
    \]
    factors through \(m\) whenever
    \[
        ge:X\to Z
    \]
    factors through \(m\).
\end{lemma}
\begin{proof}
    \leavevmode
    \begin{enumerate}
        \item \textbf{Use the kernel pair of the regular epimorphism.} Suppose \(ge=mh\). Let
        \[
            K_e\rightrightarrows X
        \]
        be the kernel pair of \(e\). Since \(e\) is the coequalizer of its kernel pair,
        \[
            ge\pi_1=ge\pi_2.
        \]
        Hence
        \[
            mh\pi_1=mh\pi_2.
        \]

        \item \textbf{Descend through the monomorphism.} Since \(m\) is monic,
        \[
            h\pi_1=h\pi_2.
        \]
        Thus \(h\) descends uniquely through \(e\), giving
        \[
            \bar h:Y\to S
        \]
        with
        \[
            \bar h e=h.
        \]

        \item \textbf{Recover the factorization of \(g\).} Since \(e\) is epi,
        \[
            m\bar h=g.
        \]
        The descent diagram is
        \[
            \begin{tikzcd}[column sep=large, row sep=large]
                    K_e
                        \arrow[r, shift left=0.5ex, "\pi_1"]
                        \arrow[r, shift right=0.5ex, "\pi_2"']
                &
                    X
                        \arrow[r, two heads, "e"]
                        \arrow[d, "h"']
                &
                    Y
                        \arrow[dl, dashed, "\bar h"]
                        \arrow[d, "g"]
                \\
                &
                    S
                        \arrow[r, hook, "m"']
                &
                    Z
            \end{tikzcd}
        \]
    \end{enumerate}
\end{proof}

Since \(\mathcal E\) is regular, the map
\[
    d_K:D_K\to\mathbf{Beh}_0
\]
has a regular epi--mono factorization in the slice
\[
    \mathcal E/Z.
\]
Write it as
\[
    D_K
        \xrightarrow{e_K}
    \operatorname{Der}_0(K)
        \xrightarrow{j_K}
    \mathbf{Beh}_0,
\]
where \(e_K\) is a regular epimorphism and \(j_K\) is monic in \(\mathcal E/Z\):
\[
    \begin{tikzcd}[column sep=large, row sep=large]
            D_K
                \arrow[rr, "d_K"]
                \arrow[dr, two heads, "e_K"']
        &
        &
            \mathbf{Beh}_0
        \\
        &
            \operatorname{Der}_0(K)
                \arrow[ur, hook, "j_K"']
        &
    \end{tikzcd}
    \qquad
    \text{in }\mathcal E/Z.
\]

\begin{definition}[Derivative image]
\label{def:derivative-image}
    The object
    \[
        \operatorname{Der}_0(K)
    \]
    is the \textbf{derivative image}, or \textbf{object-level residual closure}, of \(K\).
\end{definition}

\subsection{Derivative closure}
\label{subsec:derivative-closure}

\begin{theorem}[Derivative closure]
\label{thm:derivative-closure}
    The subobject
    \[
        j_K:\operatorname{Der}_0(K)\hookrightarrow\mathbf{Beh}_0
    \]
    is the smallest derivative-closed subobject of \(\mathbf{Beh}_0\) in
    \(\mathcal E/Z\) containing the image of
    \[
        k:K\to\mathbf{Beh}_0.
    \]
\end{theorem}
\begin{proof}
    Let
    \[
        D_K
            \xrightarrow{e_K}
        \operatorname{Der}_0(K)
            \xrightarrow{j_K}
        \mathbf{Beh}_0
    \]
    be the regular epi--mono factorization of \(d_K\) in \(\mathcal E/Z\).

    \begin{enumerate}
        \item \textbf{The image contains the specified behaviours.}

        The unit inclusion
        \[
            \iota_K:K\to D_K
        \]
        satisfies
        \[
            d_K\iota_K=k.
        \]
        Since \(d_K=j_Ke_K\), one has
        \[
            k=j_Ke_K\iota_K.
        \]
        Hence the image of \(k\) is contained in the subobject \(j_K\).

        \item \textbf{Form the derivative domain of the image.}

        Let
        \[
            \mathsf D(\operatorname{Der}_0(K))
            =
            A_1
            \times_{t_A,A_0,p_{\operatorname{Der}}}
            \operatorname{Der}_0(K),
        \]
        where
        \[
            p_{\operatorname{Der}}
            :=
            p_{\mathbf{Beh}}j_K.
        \]
        Consider the derivative composite
        \[
            g
            :=
            \partial(1_{A_1}\times j_K):
            \mathsf D(\operatorname{Der}_0(K))
            \to
            \mathbf{Beh}_0.
        \]

        Pull back the underlying regular epimorphism
        \[
            e_K:D_K\twoheadrightarrow\operatorname{Der}_0(K)
        \]
        along the projection
        \[
            \mathsf D(\operatorname{Der}_0(K))
            \to
            \operatorname{Der}_0(K)
        \]
        in \(\mathcal E\). Stability of regular epimorphisms under pullback
        gives a regular epimorphism
        \[
            e'_K:
            \widetilde D_K
            \twoheadrightarrow
            \mathsf D(\operatorname{Der}_0(K)).
        \]
        By the universal property of the relevant pullbacks, there is a
        canonical isomorphism
        \[
            \widetilde D_K
            \cong
            A_1\times_{t_A,A_0,p_{D_K}}D_K,
        \]
        where
        \[
            p_{D_K}:=p_{\mathbf{Beh}}d_K.
        \]
        We use this canonical identification below.

        \item \textbf{Define residual concatenation on the cover.}

        Define
        \[
            \chi_K:
            A_1\times_{t_A,A_0,p_{D_K}}D_K
            \to
            D_K
        \]
        by
        \[
            \chi_K(c,(a,\kappa))
            :=
            (a\circ c,\kappa).
        \]
        In morphism notation, its first component is induced by the
        multiplication
        \[
            m_A:
            A_1\times_{t_A,A_0,s_A}A_1\to A_1,
        \]
        since
        \[
            a\circ c=m_A(c,a),
        \]
        and its second component is \(\kappa\).

        \item \textbf{Compute the derivative on the regular-epimorphic cover.}

        The split derivative law gives the morphism equation
        \[
            g e'_K
            =
            j_Ke_K\chi_K.
        \]
        Indeed, in generalized-element notation,
        \[
            \begin{aligned}
                g e'_K(c,a,\kappa)
                &=
                \partial_c(j_Ke_K(a,\kappa))
                \\
                &=
                \partial_c(d_K(a,\kappa))
                \\
                &=
                \partial_c(\partial_a k(\kappa))
                \\
                &=
                \partial_{a\circ c}k(\kappa)
                \\
                &=
                d_K(a\circ c,\kappa)
                \\
                &=
                j_Ke_K\chi_K(c,a,\kappa).
            \end{aligned}
        \]
        Therefore \(ge'_K\) factors through \(j_K\).

        \item \textbf{Descend the factorization.}

        Since \(e'_K\) is a regular epimorphism and \(j_K\) is monic,
        Lemma~\ref{lem:regular-epi-descent-subobjects} gives a unique map
        \[
            \bar\partial_K:
            \mathsf D(\operatorname{Der}_0(K))
            \to
            \operatorname{Der}_0(K)
        \]
        such that
        \[
            j_K\bar\partial_K=g.
        \]
        This map lies over \(Z\): after composing with the monomorphism
        \(j_K\), both composites to \(Z\) agree, and \(j_K\) is a morphism
        over \(Z\). Thus \(j_K\) is derivative-closed.

        \item \textbf{Prove minimality.}

        Let
        \[
            i:S\hookrightarrow\mathbf{Beh}_0
        \]
        be derivative-closed and suppose that \(k\) factors through \(i\):
        \[
            k=i\bar k.
        \]
        Derivative-closedness gives a factorization
        \[
            \partial(1_{A_1}\times i)
            =
            i\bar\partial_S.
        \]
        Consequently,
        \[
            \begin{aligned}
                d_K
                &=
                \partial(1_{A_1}\times k)
                \\
                &=
                \partial(1_{A_1}\times i)
                (1_{A_1}\times\bar k)
                \\
                &=
                i\bar\partial_S
                (1_{A_1}\times\bar k).
            \end{aligned}
        \]
        Hence \(d_K\) factors through \(i\). Since \(j_K\) is the regular
        image of \(d_K\), the image subobject \(j_K\) factors through \(i\).

        Therefore \(j_K\) is the smallest derivative-closed subobject
        containing the image of \(k\).
    \end{enumerate}
\end{proof}

\section{The Categorical Myhill--Nerode Theorem}
\label{sec:categorical-myhill-nerode-theorem}

We now assemble the equality-level Nerode relation, the derivative image, the minimal residual realization, and the reachability--separatedness criterion. The result is a categorical Myhill--Nerode theorem for the residual semantics developed above. Related categorical minimization theorems for automata and languages appear in functorial and topos-theoretic settings \cite{ColcombetPetrisan2020,Iwaniack2023}. The theorem below is stated for Mealy morphisms between internal categories and uses the equality of represented residual behaviours.

\subsection{The Nerode relation}
\label{subsec:nerode-relation}

\begin{definition}[Nerode relation]
\label{def:nerode-relation}
    The \textbf{Nerode relation}
    \[
        {\equiv_K}\rightrightarrows D_K
    \]
    is the kernel pair of
    \[
        d_K:D_K\to\mathbf{Beh}_0
    \]
    in the slice
    \[
        \mathcal E/Z,
        \qquad
        Z=A_0\times B_0.
    \]
    Thus
    \[
        (a,\kappa)\equiv_K(a',\kappa')
    \]
    precisely when
    \[
        \partial_a k(\kappa)=\partial_{a'}k(\kappa')
    \]
    as represented residual behaviours.
\end{definition}

The kernel-pair square is
\[
    \begin{tikzcd}[column sep=large, row sep=large]
            {\equiv_K}
                \arrow[r, "\pi_2"]
                \arrow[d, "\pi_1"']
                \arrow[dr, phantom, "\lrcorner", very near start]
        &
            D_K
                \arrow[d, "d_K"]
        \\
            D_K
                \arrow[r, "d_K"']
        &
            \mathbf{Beh}_0.
    \end{tikzcd}
\]

\begin{remark}
\label{rem:equality-level-nerode}
    This is an equality-level Nerode relation. It identifies formal residual derivatives exactly when they determine the same represented residual functor, including both object and arrow assignments. It does not quotient by isomorphism, by behaviour comparison, or by equivalence inside the vertical behaviour category.
\end{remark}

\subsection{The Nerode quotient}
\label{subsec:nerode-quotient}

\begin{theorem}[Categorical Myhill--Nerode quotient]
\label{thm:categorical-myhill-nerode-quotient}
    The quotient
    \[
        D_K/{\equiv_K}
    \]
    exists and is effective in
    \[
        \mathcal E/Z.
    \]
    Moreover, there is a canonical isomorphism
    \[
        D_K/{\equiv_K}
        \cong
        \operatorname{Der}_0(K)
    \]
    in \(\mathcal E/Z\).
\end{theorem}
\begin{proof}
    \leavevmode
    \begin{enumerate}
        \item \textbf{Use exactness of the slice.} Since \(\mathcal E\) is Barr-exact, the slice \(\mathcal E/Z\) is Barr-exact. Hence the kernel pair of \(d_K\) has an effective quotient.

        \item \textbf{Identify the quotient with the image.} In a Barr-exact category, the quotient of the kernel pair of a morphism is its regular image. Therefore
        \[
            D_K/{\equiv_K}
            \cong
            \operatorname{Im}(d_K)
            =
            \operatorname{Der}_0(K).
        \]

        \item \textbf{Display the comparison.} The comparison is represented by the diagram
        \[
            \begin{tikzcd}[column sep=large, row sep=large]
                    {\equiv_K}
                        \arrow[r, shift left=0.5ex]
                        \arrow[r, shift right=0.5ex]
                &
                    D_K
                        \arrow[r, two heads]
                        \arrow[dr, two heads, "e_K"']
                &
                    D_K/{\equiv_K}
                        \arrow[d, dashed, "\cong"]
                \\
                &
                &
                    \operatorname{Der}_0(K).
            \end{tikzcd}
        \]
    \end{enumerate}
\end{proof}

\subsection{The minimal residual realization}
\label{subsec:minimal-residual-realization}

\begin{definition}[Minimal residual realization]
\label{def:minimal-residual-realization}
    The \textbf{minimal residual realization} of a specification
    \[
        k:K\to\mathbf{Beh}_0
    \]
    is the restriction of the canonical behaviour Mealy morphism to the derivative-closed image
    \[
        \operatorname{Der}_0(K)\hookrightarrow\mathbf{Beh}_0.
    \]
    It is denoted
    \[
        \mathsf{Min}(K):\mathbb A\rightsquigarrow\mathbb B.
    \]
\end{definition}

The minimal realization carries a canonical realization map
\[
    i_{\min}:K\to\operatorname{Der}_0(K)
\]
defined by
\[
    i_{\min}:=e_K\iota_K.
\]
It satisfies
\[
    j_K i_{\min}
    =
    j_Ke_K\iota_K
    =
    d_K\iota_K
    =
    k.
\]

\subsection{Realizations and reachability}
\label{subsec:realizations-reachability}

\begin{definition}[Realization]
\label{def:realization}
    A \textbf{realization} of a behaviour specification
    \[
        k:K\to\mathbf{Beh}_0
    \]
    is a Mealy morphism
    \[
        P:\mathbb A\rightsquigarrow\mathbb B
    \]
    together with a map
    \[
        i:K\to P
    \]
    such that
    \[
        \beta_P i=k.
    \]
\end{definition}

\begin{remark}
\label{rem:realization-over-base}
    Since \(\beta_P\) and \(k\) both lie over \(Z=A_0\times B_0\), the equation
    \[
        \beta_P i=k
    \]
    implies that \(i\) is a map over \(Z\).
\end{remark}

\begin{definition}[Reachability map]
\label{def:reachability-map}
    Given a realization
    \[
        i:K\to P,
    \]
    the \textbf{reachability map}
    \[
        \rho_{P,K}:D_K\to P
    \]
    is defined by
    \[
        \rho_{P,K}(a,\kappa)
        =
        \delta_P(a,i(\kappa)).
    \]
\end{definition}

\begin{lemma}
\label{lem:reachability-map-over-base}
    The reachability map satisfies
    \[
        \beta_P\rho_{P,K}=d_K.
    \]
    Consequently, \(\rho_{P,K}\) is a map over \(Z\), where \(D_K\) is viewed over \(Z\) through \(d_K\).
\end{lemma}
\begin{proof}
    \leavevmode
    \begin{enumerate}
        \item \textbf{Compute the behaviour of a reachable state.} For \((a,\kappa)\in D_K\),
        \[
            \begin{aligned}
                \beta_P(\rho_{P,K}(a,\kappa))
                &=
                \beta_P(\delta_P(a,i(\kappa)))
                \\
                &=
                \partial_a(\beta_P(i(\kappa)))
                \\
                &=
                \partial_a k(\kappa)
                \\
                &=
                d_K(a,\kappa).
            \end{aligned}
        \]

        \item \textbf{Display the commutative triangle.} This is the commutativity of
        \[
            \begin{tikzcd}[column sep=large, row sep=large]
                    D_K
                        \arrow[r, "\rho_{P,K}"]
                        \arrow[dr, "d_K"']
                &
                    P
                        \arrow[d, "\beta_P"]
                \\
                &
                    \mathbf{Beh}_0.
            \end{tikzcd}
        \]
    \end{enumerate}
\end{proof}

\begin{definition}[Transition-closed subobject]
\label{def:transition-closed-subobject}
    Let
    \[
        P:\mathbb A\rightsquigarrow\mathbb B
    \]
    be a Mealy morphism, and let
    \[
        m:S\hookrightarrow P
    \]
    be a subobject in \(\mathcal E/Z\), where \(Z=A_0\times B_0\). Write
    \[
        p_S:=p_Pm.
    \]
    The \textbf{transition domain} of \(S\) is
    \[
        \mathsf T(S):=
        A_1\times_{t_A,A_0,p_S}S.
    \]
    We regard \(\mathsf T(S)\) as an object of \(\mathcal E/Z\) through the composite
    \[
        \mathsf T(S)
            \xrightarrow{1_{A_1}\times m}
        A_1\times_{t_A,A_0,p_P}P
            \xrightarrow{\delta_P}
        P
            \xrightarrow{(p_P,q_P)}
        Z.
    \]

    We say that \(S\) is \textbf{transition-closed} if the transition map
    \[
        \delta_P(1_{A_1}\times m):
        \mathsf T(S)\to P
    \]
    factors through \(m\) in \(\mathcal E/Z\). Equivalently, there is a unique map
    \[
        \bar\delta_S:\mathsf T(S)\to S
    \]
    over \(Z\) such that
    \[
        m\bar\delta_S
        =
        \delta_P(1_{A_1}\times m).
    \]
\end{definition}

The factorization condition is
\[
    \begin{tikzcd}[column sep=large, row sep=large]
            \mathsf T(S)
                \arrow[rr, "{\delta_P(1_{A_1}\times m)}"]
                \arrow[dr, "\bar\delta_S"']
        &
        &
            P
        \\
        &
            S
                \arrow[ur, hook, "m"']
        &
    \end{tikzcd}
    \qquad
    \text{in }\mathcal E/Z.
\]

\begin{lemma}[Restriction to transition-closed subobjects]
\label{lem:restriction-transition-closed-subobject}
    Let
    \[
        m:S\hookrightarrow P
    \]
    be a transition-closed subobject of a Mealy morphism
    \[
        P:\mathbb A\rightsquigarrow\mathbb B
    \]
    in \(\mathcal E/Z\). Then \(S\), with the restricted carrier
    \[
        A_0\xleftarrow{p_Pm}S\xrightarrow{q_Pm}B_0,
    \]
    carries a canonical restricted Mealy morphism structure.

    Its transition is the factorization
    \[
        \bar\delta_S:
        A_1\times_{t_A,A_0,p_Pm}S\to S
    \]
    supplied by transition-closedness, and its output is the restriction of
    \[
        \omega_P:
        A_1\times_{t_A,A_0,p_P}P\to B_1
    \]
    along
    \[
        1_{A_1}\times m.
    \]
\end{lemma}
\begin{proof}
    Write
    \[
        p_S:=p_Pm,
        \qquad
        q_S:=q_Pm.
    \]
    Transition-closedness supplies a unique map
    \[
        \bar\delta_S:
        A_1\times_{t_A,A_0,p_S}S
        \to
        S
    \]
    satisfying
    \[
        m\bar\delta_S
        =
        \delta_P(1_{A_1}\times m).
    \]
    Define
    \[
        \delta_S:=\bar\delta_S.
    \]

    Define the restricted output by
    \[
        \omega_S
        :=
        \omega_P\circ(1_{A_1}\times m).
    \]
    In generalized-element notation,
    \[
        \delta_S(a,s)=\bar\delta_S(a,s),
    \]
    and
    \[
        \omega_S(a,s)=\omega_P(a,m(s)).
    \]

    Since \(\bar\delta_S\) is a morphism over the transition-induced
    \(Z\)-structure,
    \[
        p_S(\delta_S(a,s))=s_A(a).
    \]
    The output source equation is
    \[
        \begin{aligned}
            s_B\omega_S(a,s)
            &=
            s_B\omega_P(a,m(s))
            \\
            &=
            q_P(\delta_P(a,m(s)))
            \\
            &=
            q_P(m(\delta_S(a,s)))
            \\
            &=
            q_S(\delta_S(a,s)).
        \end{aligned}
    \]
    The output target equation is
    \[
        \begin{aligned}
            t_B\omega_S(a,s)
            &=
            t_B\omega_P(a,m(s))
            \\
            &=
            q_P(m(s))
            \\
            &=
            q_S(s).
        \end{aligned}
    \]

    For the transition unit law, compose with \(m\):
    \[
        \begin{aligned}
            m(\delta_S(1_{p_S(s)},s))
            &=
            \delta_P(1_{p_P(m(s))},m(s))
            \\
            &=
            m(s).
        \end{aligned}
    \]
    Since \(m\) is monic,
    \[
        \delta_S(1_{p_S(s)},s)=s.
    \]

    For the transition multiplication law, let
    \[
        a:u\to v,
        \qquad
        b:v\to w,
        \qquad
        p_S(s)=w.
    \]
    After composing with \(m\),
    \[
        \begin{aligned}
            m(\delta_S(b\circ a,s))
            &=
            \delta_P(b\circ a,m(s))
            \\
            &=
            \delta_P(a,\delta_P(b,m(s)))
            \\
            &=
            \delta_P(a,m(\delta_S(b,s)))
            \\
            &=
            m(\delta_S(a,\delta_S(b,s))).
        \end{aligned}
    \]
    Hence, by monicity,
    \[
        \delta_S(b\circ a,s)
        =
        \delta_S(a,\delta_S(b,s)).
    \]

    The output unit law is the restriction of the output unit law of \(P\):
    \[
        \begin{aligned}
            \omega_S(1_{p_S(s)},s)
            &=
            \omega_P(1_{p_P(m(s))},m(s))
            \\
            &=
            1_{q_P(m(s))}
            \\
            &=
            1_{q_S(s)}.
        \end{aligned}
    \]

    For the output multiplication law,
    \[
        \begin{aligned}
            \omega_S(b\circ a,s)
            &=
            \omega_P(b\circ a,m(s))
            \\
            &=
            \omega_P(b,m(s))
            \circ
            \omega_P(a,\delta_P(b,m(s)))
            \\
            &=
            \omega_P(b,m(s))
            \circ
            \omega_P(a,m(\delta_S(b,s)))
            \\
            &=
            \omega_S(b,s)
            \circ
            \omega_S(a,\delta_S(b,s)).
        \end{aligned}
    \]

    Therefore the restricted carrier, transition, and output define a
    Mealy morphism
    \[
        \mathbb A\rightsquigarrow\mathbb B.
    \]
    By construction, the inclusion
    \[
        m:S\to P
    \]
    preserves transition and output strictly.
\end{proof}

\begin{definition}[Reachable part]
\label{def:reachable-part}
    The \textbf{reachable part}
    \[
        \operatorname{Reach}_0(P,K)\hookrightarrow P
    \]
    is the regular image of
    \[
        \rho_{P,K}:D_K\to P
    \]
    in the slice
    \[
        \mathcal E/Z.
    \]
\end{definition}

\begin{proposition}[Reachable part]
\label{prop:reachable-part}
    Let
    \[
        i:K\to P
    \]
    be a realization of
    \[
        k:K\to\mathbf{Beh}_0.
    \]
    Let
    \[
        D_K
            \xrightarrow{e_R}
        \operatorname{Reach}_0(P,K)
            \xrightarrow{m_R}
        P
    \]
    be the regular epi--mono factorization of the reachability map
    \[
        \rho_{P,K}:D_K\to P
    \]
    in \(\mathcal E/Z\).

    Then:

    \begin{enumerate}
        \item the subobject
        \[
            m_R:
            \operatorname{Reach}_0(P,K)\hookrightarrow P
        \]
        is transition-closed;

        \item the restricted Mealy structure makes \(m_R\) a strict semantic
        Mealy homomorphism;

        \item the map
        \[
            i_R:=e_R\iota_K:
            K\to\operatorname{Reach}_0(P,K)
        \]
        is a realization of \(k\);

        \item the realization \(i_R\) is reachable;

        \item there is a regular epimorphic strict semantic homomorphism
        \[
            h:
            \operatorname{Reach}_0(P,K)
            \twoheadrightarrow
            \operatorname{Der}_0(K)
        \]
        satisfying
        \[
            j_Kh=\beta_Pm_R.
        \]
    \end{enumerate}
\end{proposition}
\begin{proof}
    \begin{enumerate}
        \item \textbf{Prove transition-closedness.}

        Put
        \[
            p_{\operatorname{Reach}}
            :=
            p_Pm_R.
        \]
        Consider the transition composite
        \[
            g:
            A_1
            \times_{t_A,A_0,p_{\operatorname{Reach}}}
            \operatorname{Reach}_0(P,K)
            \to
            P
        \]
        defined by
        \[
            g
            :=
            \delta_P(1_{A_1}\times m_R).
        \]

        Pull back the underlying regular epimorphism
        \[
            e_R:D_K\twoheadrightarrow\operatorname{Reach}_0(P,K)
        \]
        along the projection
        \[
            A_1
            \times_{t_A,A_0,p_{\operatorname{Reach}}}
            \operatorname{Reach}_0(P,K)
            \to
            \operatorname{Reach}_0(P,K)
        \]
        in \(\mathcal E\). This gives a regular epimorphism
        \[
            e'_R:
            \widetilde D_R
            \twoheadrightarrow
            A_1
            \times_{t_A,A_0,p_{\operatorname{Reach}}}
            \operatorname{Reach}_0(P,K).
        \]
        There is a canonical pullback isomorphism
        \[
            \widetilde D_R
            \cong
            A_1\times_{t_A,A_0,p_{D_K}}D_K.
        \]
        We use this identification below.

        Define
        \[
            \chi_{P,K}:
            A_1\times_{t_A,A_0,p_{D_K}}D_K
            \to
            D_K
        \]
        by
        \[
            \chi_{P,K}(c,(a,\kappa))
            =
            (a\circ c,\kappa).
        \]
        The Mealy multiplication law for \(\delta_P\) gives
        \[
            g e'_R
            =
            m_Re_R\chi_{P,K}.
        \]
        In generalized-element notation,
        \[
            \begin{aligned}
                g e'_R(c,a,\kappa)
                &=
                \delta_P(c,m_Re_R(a,\kappa))
                \\
                &=
                \delta_P(c,\rho_{P,K}(a,\kappa))
                \\
                &=
                \delta_P(c,\delta_P(a,i(\kappa)))
                \\
                &=
                \delta_P(a\circ c,i(\kappa))
                \\
                &=
                \rho_{P,K}(a\circ c,\kappa)
                \\
                &=
                m_Re_R\chi_{P,K}(c,a,\kappa).
            \end{aligned}
        \]
        Hence \(ge'_R\) factors through \(m_R\). Since \(e'_R\) is a regular
        epimorphism and \(m_R\) is monic,
        Lemma~\ref{lem:regular-epi-descent-subobjects} supplies a unique map
        \[
            \bar\delta_R:
            A_1
            \times_{t_A,A_0,p_{\operatorname{Reach}}}
            \operatorname{Reach}_0(P,K)
            \to
            \operatorname{Reach}_0(P,K)
        \]
        such that
        \[
            m_R\bar\delta_R=g.
        \]
        Thus \(m_R\) is transition-closed.

        \item \textbf{Restrict the Mealy structure.}

        By Lemma~\ref{lem:restriction-transition-closed-subobject}, the
        reachable part carries the restricted Mealy structure. Its transition
        is \(\bar\delta_R\), and its output is the restriction of
        \(\omega_P\). By construction,
        \[
            m_R:
            \operatorname{Reach}_0(P,K)\to P
        \]
        preserves transition and output strictly and lies over \(Z\).
        Hence it is a strict semantic Mealy homomorphism.

        \item \textbf{Construct the realization map.}

        Define
        \[
            i_R:=e_R\iota_K.
        \]
        Since
        \[
            m_Re_R=\rho_{P,K},
        \]
        the Mealy unit law gives
        \[
            \begin{aligned}
                m_Ri_R
                &=
                m_Re_R\iota_K
                \\
                &=
                \rho_{P,K}\iota_K
                \\
                &=
                i.
            \end{aligned}
        \]

        The composite
        \[
            \beta_Pm_R:
            \operatorname{Reach}_0(P,K)\to\mathbf{Beh}_0
        \]
        is a strict semantic homomorphism. By terminality of the canonical
        behaviour machine,
        \[
            \beta_{\operatorname{Reach}}
            =
            \beta_Pm_R.
        \]
        Therefore
        \[
            \begin{aligned}
                \beta_{\operatorname{Reach}}i_R
                &=
                \beta_Pm_Re_R\iota_K
                \\
                &=
                \beta_P\rho_{P,K}\iota_K
                \\
                &=
                d_K\iota_K
                \\
                &=
                k.
            \end{aligned}
        \]
        Thus \(i_R\) is a realization of \(k\).

        \item \textbf{Show that the restricted realization is reachable.}

        Let
        \[
            \rho_{R,K}:
            D_K\to\operatorname{Reach}_0(P,K)
        \]
        be the reachability map of the realization \(i_R\). After composing
        with \(m_R\), one has
        \[
            \begin{aligned}
                m_R\rho_{R,K}(a,\kappa)
                &=
                \delta_P(a,m_Ri_R(\kappa))
                \\
                &=
                \delta_P(a,i(\kappa))
                \\
                &=
                \rho_{P,K}(a,\kappa)
                \\
                &=
                m_Re_R(a,\kappa).
            \end{aligned}
        \]
        Since \(m_R\) is monic,
        \[
            \rho_{R,K}=e_R.
        \]
        The map \(e_R\) is a regular epimorphism, so its image is all of
        \(\operatorname{Reach}_0(P,K)\). Hence \(i_R\) is reachable.

        \item \textbf{Construct the quotient map to the derivative image.}

        Let
        \[
            D_K
                \xrightarrow{e_K}
            \operatorname{Der}_0(K)
                \xrightarrow{j_K}
            \mathbf{Beh}_0
        \]
        be the image factorization of \(d_K\). By
        Lemma~\ref{lem:reachability-map-over-base},
        \[
            \beta_Pm_Re_R
            =
            \beta_P\rho_{P,K}
            =
            d_K
            =
            j_Ke_K.
        \]

        Let
        \[
            R_{e_R}\rightrightarrows D_K
        \]
        be the kernel pair of \(e_R\), with projections
        \(\pi_1,\pi_2\). Then
        \[
            \begin{aligned}
                j_Ke_K\pi_1
                &=
                \beta_Pm_Re_R\pi_1
                \\
                &=
                \beta_Pm_Re_R\pi_2
                \\
                &=
                j_Ke_K\pi_2.
            \end{aligned}
        \]
        Since \(j_K\) is monic,
        \[
            e_K\pi_1=e_K\pi_2.
        \]
        Since \(e_R\) is the coequalizer of its kernel pair, there is a unique
        map
        \[
            h:
            \operatorname{Reach}_0(P,K)
            \to
            \operatorname{Der}_0(K)
        \]
        such that
        \[
            he_R=e_K.
        \]
        Since both \(j_Kh\) and \(\beta_Pm_R\) agree after precomposition with
        the epimorphism \(e_R\),
        \[
            j_Kh=\beta_Pm_R.
        \]

        Factor \(h\) as
        \[
            \operatorname{Reach}_0(P,K)
                \twoheadrightarrow
            I
                \rightarrowtail
            \operatorname{Der}_0(K)
        \]
        in \(\mathcal E/Z\). The regular epimorphism
        \[
            e_K=he_R
        \]
        factors through the displayed monomorphism \(I\rightarrowtail
        \operatorname{Der}_0(K)\). Since regular epimorphisms are extremal,
        that monomorphism is an isomorphism. Hence \(h\) is a regular
        epimorphism.

        Finally, \(h\) is a strict semantic homomorphism because
        \[
            j_Kh=\beta_Pm_R,
        \]
        the right-hand side is strict, and \(j_K\) is the strict inclusion of
        a derivative-closed canonical submachine.
    \end{enumerate}
\end{proof}

\begin{definition}[Reachable realization]
\label{def:reachable-realization}
    A realization \(i:K\to P\) is \textbf{reachable} if
    \[
        \operatorname{Reach}_0(P,K)\hookrightarrow P
    \]
    is an isomorphism.
\end{definition}

\subsection{Semantic quotient and minimality}
\label{subsec:semantic-quotient-minimality}

\begin{theorem}[Semantic quotient theorem]
\label{thm:semantic-quotient-theorem}
    Let
    \[
        i:K\to P
    \]
    be a realization of
    \[
        k:K\to\mathbf{Beh}_0.
    \]
    The behaviour assignment restricts on the reachable part to a strict semantic homomorphism
    \[
        \overline{\beta}_{P,K}:
        \operatorname{Reach}_0(P,K)\to\mathsf{Min}(K).
    \]
    This map is a regular epimorphism in \(\mathcal E/Z\). If \(P\) is reachable, then
    \[
        \overline{\beta}_{P,K}:P\twoheadrightarrow\mathsf{Min}(K)
    \]
    is a regular epimorphic strict semantic homomorphism.
\end{theorem}
\begin{proof}
    \leavevmode
    \begin{enumerate}
        \item \textbf{Define the restricted semantic map.} By Proposition~\ref{prop:reachable-part}, there is a regular epimorphism
        \[
            h:\operatorname{Reach}_0(P,K)\twoheadrightarrow\operatorname{Der}_0(K)
        \]
        satisfying
        \[
            j_Kh=\beta_Pm_R.
        \]
        Define
        \[
            \overline{\beta}_{P,K}:=h.
        \]

        \item \textbf{Restrict the Mealy structures.} Since the reachable part is transition-closed and \(\operatorname{Der}_0(K)\) is derivative-closed, the source carries its restricted Mealy structure by Lemma~\ref{lem:restriction-transition-closed-subobject}, and the target carries its restricted Mealy structure by Lemma~\ref{lem:restriction-derivative-closed-subobject}.

        \item \textbf{Check strict preservation.} The map \(\overline{\beta}_{P,K}\) is a strict semantic homomorphism because it is the restriction of the strict semantic homomorphism
        \[
            \beta_P:P\to
            \mathbf{Beh}^{\mathrm{can}}_{\mathbb A,\mathbb B}.
        \]
        For example, transition preservation follows after composing with \(j_K\): both sides become the corresponding transition-preservation equation for \(\beta_P\). Since \(j_K\) is monic, the equality holds in \(\operatorname{Der}_0(K)\). The output equation is inherited directly from the restricted output structure.

        \item \textbf{Conclude regular epimorphy.} Its regular epimorphy is the regular epimorphy of \(h\). If \(P\) is reachable, then \(\operatorname{Reach}_0(P,K)\cong P\), giving the final statement.
    \end{enumerate}
\end{proof}

\begin{definition}[Separated realization]
\label{def:separated-realization}
    A realization \(i:K\to P\) is \textbf{behaviourally separated} if the Mealy morphism \(P\) is behaviourally separated, equivalently if
    \[
        \beta_P:P\to\mathbf{Beh}_0
    \]
    is monic.
\end{definition}

\begin{theorem}[Minimal realization theorem]
\label{thm:minimal-realization}
    If a realization
    \[
        i:K\to P
    \]
    is reachable and behaviourally separated, then the canonical quotient map
    \[
        \overline{\beta}_{P,K}:P\to\mathsf{Min}(K)
    \]
    is an isomorphism of Mealy morphisms. Moreover,
    \[
        \overline{\beta}_{P,K}i=i_{\min}.
    \]
    Hence every reachable separated realization of \(K\) is isomorphic to
    \[
        \mathsf{Min}(K)
    \]
    as a realization.
\end{theorem}
\begin{proof}
    \leavevmode
    \begin{enumerate}
        \item \textbf{Use reachability to get a regular epimorphism.} By Theorem~\ref{thm:semantic-quotient-theorem}, reachability gives a regular epimorphic strict semantic homomorphism
        \[
            \overline{\beta}_{P,K}:P\twoheadrightarrow\mathsf{Min}(K).
        \]

        \item \textbf{Use separatedness to get a monomorphism.} Let
        \[
            j_K:\operatorname{Der}_0(K)\hookrightarrow\mathbf{Beh}_0
        \]
        be the image inclusion. Then
        \[
            j_K\overline{\beta}_{P,K}=\beta_P.
        \]
        Since \(P\) is behaviourally separated, \(\beta_P\) is monic. Since \(j_K\) is monic, \(\overline{\beta}_{P,K}\) is monic. A morphism that is both regular epimorphic and monic in a regular category is an isomorphism.

        \item \textbf{Promote the isomorphism to Mealy morphisms.} Since \(\overline{\beta}_{P,K}\) is a strict semantic homomorphism and an isomorphism over \(Z\), its inverse preserves transitions and outputs by applying the inverse to the strict preservation equations. Hence it is an isomorphism of Mealy morphisms.

        \item \textbf{Compare realization maps.} Finally,
        \[
            j_K\overline{\beta}_{P,K}i
            =
            \beta_P i
            =
            k
            =
            j_K i_{\min}.
        \]
        Since \(j_K\) is monic,
        \[
            \overline{\beta}_{P,K}i=i_{\min}.
        \]
    \end{enumerate}
\end{proof}

\subsection{The categorical theorem}
\label{subsec:categorical-myhill-nerode-theorem}

\begin{theorem}[Categorical Myhill--Nerode theorem]
\label{thm:categorical-myhill-nerode}
    Let
    \[
        k:K\to\mathbf{Beh}_0
    \]
    be a behaviour specification. Then:
    \begin{enumerate}
        \item the equality-level Nerode quotient
        \[
            D_K/{\equiv_K}
        \]
        exists effectively in
        \[
            \mathcal E/(A_0\times B_0);
        \]

        \item there is a canonical isomorphism
        \[
            D_K/{\equiv_K}\cong\operatorname{Der}_0(K);
        \]

        \item under this isomorphism, the quotient carries the restricted canonical Mealy structure;

        \item this restricted Mealy morphism is the minimal residual realization
        \[
            \mathsf{Min}(K);
        \]

        \item every realization \(i:K\to P\) has a reachable part mapping by a regular epimorphic strict semantic homomorphism onto
        \[
            \mathsf{Min}(K);
        \]

        \item every reachable behaviourally separated realization of \(K\) is isomorphic to
        \[
            \mathsf{Min}(K)
        \]
        as a realization.
    \end{enumerate}
\end{theorem}
\begin{proof}
    \leavevmode
    \begin{enumerate}
        \item \textbf{Apply the quotient theorem.} The first two claims are Theorem~\ref{thm:categorical-myhill-nerode-quotient}.

        \item \textbf{Restrict the canonical Mealy structure.} By Theorem~\ref{thm:derivative-closure},
        \[
            \operatorname{Der}_0(K)\hookrightarrow\mathbf{Beh}_0
        \]
        is derivative-closed. Therefore, by Lemma~\ref{lem:restriction-derivative-closed-subobject}, the canonical behaviour Mealy morphism restricts to \(\operatorname{Der}_0(K)\). This restricted Mealy morphism is \(\mathsf{Min}(K)\) by Definition~\ref{def:minimal-residual-realization}.

        \item \textbf{Apply the semantic quotient and minimality theorems.} The regular epimorphic quotient from the reachable part of any realization onto \(\mathsf{Min}(K)\) is Theorem~\ref{thm:semantic-quotient-theorem}. The uniqueness of reachable separated realizations is Theorem~\ref{thm:minimal-realization}.
    \end{enumerate}
\end{proof}

\subsection{Finite-state corollary}
\label{subsec:finite-state-corollary}

\begin{definition}[Finite-state structure]
\label{def:finite-state-doctrine}
    A \textbf{finite-state structure} on \(\mathcal E\) is a choice, for every \(Z\in\mathcal E\), of a replete full subcategory
    \[
        \mathsf{Fin}_Z\subseteq\mathcal E/Z
    \]
    closed under:
    \begin{enumerate}
        \item isomorphisms;

        \item subobjects;

        \item regular images;

        \item effective quotients of internal equivalence relations.
    \end{enumerate}
\end{definition}

For the remainder of this subsection, put
\[
    Z=A_0\times B_0.
\]

\begin{definition}[Finite-state realization]
\label{def:finite-state-realization}
    A realization with carrier
    \[
        P\to Z
    \]
    is \textbf{finite-state} if
    \[
        P\to Z
    \]
    lies in
    \[
        \mathsf{Fin}_Z.
    \]
\end{definition}

\begin{corollary}[Finite-state Myhill--Nerode theorem]
\label{cor:finite-state-myhill-nerode}
    Let
    \[
        k:K\to\mathbf{Beh}_0
    \]
    be a behaviour specification. The following are equivalent:
    \begin{enumerate}
        \item \(K\) has a finite-state realization.

        \item \(K\) has a finite-state reachable behaviourally separated realization.

        \item The carrier
        \[
            \operatorname{Der}_0(K)\to Z
        \]
        of
        \[
            \mathsf{Min}(K)
        \]
        lies in \(\mathsf{Fin}_Z\).

        \item The Nerode quotient
        \[
            D_K/{\equiv_K}
        \]
        lies in \(\mathsf{Fin}_Z\).
    \end{enumerate}
\end{corollary}
\begin{proof}
    \leavevmode
    \begin{enumerate}
        \item \textbf{Derive finiteness of the minimal carrier from a finite realization.} Suppose \(K\) has a finite-state realization \(i:K\to P\). The reachable part
        \[
            \operatorname{Reach}_0(P,K)\hookrightarrow P
        \]
        is a subobject of \(P\) in \(\mathcal E/Z\), hence lies in \(\mathsf{Fin}_Z\). By Theorem~\ref{thm:semantic-quotient-theorem}, there is a regular epimorphism
        \[
            \operatorname{Reach}_0(P,K)\twoheadrightarrow\mathsf{Min}(K)
        \]
        in \(\mathcal E/Z\). Closure under regular images gives
        \[
            \operatorname{Der}_0(K)\in\mathsf{Fin}_Z.
        \]
        Thus (1) implies (3).

        \item \textbf{Show the minimal realization is reachable and separated.}
        If
        \[
            \operatorname{Der}_0(K)\in\mathsf{Fin}_Z,
        \]
        then \(\mathsf{Min}(K)\) is finite-state. Its realization map is
        \[
            i_{\min}:K\to\operatorname{Der}_0(K).
        \]
        The reachability map of this realization has image
        \[
            \operatorname{Der}_0(K),
        \]
        because its composite with the inclusion
        \[
            \operatorname{Der}_0(K)\hookrightarrow\mathbf{Beh}_0
        \]
        is
        \[
            (a,\kappa)
            \longmapsto
            \partial_a k(\kappa)
            =
            d_K(a,\kappa).
        \]
        Hence \(\mathsf{Min}(K)\) is reachable. It is behaviourally separated because, by Lemma~\ref{lem:canonical-behaviour-map-identity}, its behaviour map is the monomorphism
        \[
            \operatorname{Der}_0(K)\hookrightarrow\mathbf{Beh}_0.
        \]
        Thus (3) implies (2).

        \item \textbf{Conclude the remaining implications.} The implication (2) implies (1) is immediate. Finally, Theorem~\ref{thm:categorical-myhill-nerode-quotient} gives
        \[
            D_K/{\equiv_K}\cong\operatorname{Der}_0(K)
        \]
        in \(\mathcal E/Z\), so (3) and (4) are equivalent.
    \end{enumerate}
\end{proof}

\begin{remark}
\label{rem:finite-state-doctrine-quotients}
    In the proof above, finite realizability implies finiteness of the minimal carrier using closure under subobjects and regular images. Closure under effective quotients ensures that, when a formal derivative object lies in the finite-state structure, its Nerode quotient remains in the same finite-state structure.
\end{remark}

\begin{remark}
\label{rem:finite-arrow-object}
    For finite categorical semantics, one may also require the arrow object
    \[
        \mathbf{Der}(K)_1
    \]
    to lie in the chosen finite class over
    \[
        \operatorname{Der}_0(K)\times_{A_0}\operatorname{Der}_0(K).
    \]
    This records finiteness of comparison structure as well as finiteness of residual states.
\end{remark}

\section{Generated residual behaviour categories}
\label{sec:generated-residual-behaviour-categories}

The minimal Mealy carrier records the object image of the derivative orbit. The vertical behaviour category also records behaviour comparisons between generated residual behaviours.

\subsection{Derivative-closed full internal subcategories}
\label{subsec:derivative-closed-full-internal-subcategories}

\begin{definition}[Derivative-closed full internal subcategory]
\label{def:derivative-closed-full-subcategory}
    A full internal subcategory
    \[
        \mathbb S\subseteq
        \mathbf{Beh}^{\mathrm{vert}}_{\mathbb A,\mathbb B}
    \]
    is \textbf{derivative-closed} if its object subobject
    \[
        S_0\hookrightarrow\mathbf{Beh}_0
    \]
    is derivative-closed in the sense of Definition~\ref{def:derivative-closed-subobject}.
\end{definition}

\subsection{The generated residual behaviour category}
\label{subsec:generated-residual-behaviour-category}

\begin{definition}[Generated residual behaviour category]
\label{def:generated-residual-behaviour-category}
    The \textbf{generated residual behaviour category} of a specification
    \[
        k:K\to\mathbf{Beh}_0
    \]
    is the full internal subcategory
    \[
        \mathbf{Der}(K)
        \subseteq
        \mathbf{Beh}^{\mathrm{vert}}_{\mathbb A,\mathbb B}
    \]
    on the object subobject
    \[
        \operatorname{Der}_0(K)\hookrightarrow\mathbf{Beh}_0.
    \]
    Thus
    \[
        \mathbf{Der}(K)_0=\operatorname{Der}_0(K),
    \]
    and
    \[
        \mathbf{Der}(K)_1
        =
        \operatorname{Der}_0(K)
        \times_{\mathbf{Beh}_0,s}
        \mathbf{Beh}_1
        \times_{t,\mathbf{Beh}_0}
        \operatorname{Der}_0(K).
    \]
\end{definition}

The arrow object is defined by the pullback rectangle
\[
    \begin{tikzcd}[column sep=large, row sep=large]
            \mathbf{Der}(K)_1
                \arrow[r]
                \arrow[d]
                \arrow[dr, phantom, "\lrcorner", very near start]
        &
            \mathbf{Beh}_1
                \arrow[d, "{(s,t)}"]
        \\
            \operatorname{Der}_0(K)\times_{A_0}\operatorname{Der}_0(K)
                \arrow[r, "{(j_K,j_K)}"']
        &
            \mathbf{Beh}_0\times_{A_0}\mathbf{Beh}_0.
    \end{tikzcd}
\]

\begin{theorem}[Generated residual behaviour category]
\label{thm:generated-residual-behaviour-category}
    The internal category
    \[
        \mathbf{Der}(K)
    \]
    is the smallest full derivative-closed internal subcategory of
    \[
        \mathbf{Beh}^{\mathrm{vert}}_{\mathbb A,\mathbb B}
    \]
    containing the specified behaviours.
\end{theorem}
\begin{proof}
    \leavevmode
    \begin{enumerate}
        \item \textbf{Use the object-level closure.} By Theorem~\ref{thm:derivative-closure}, the object part
        \[
            \operatorname{Der}_0(K)
        \]
        is the smallest derivative-closed subobject of \(\mathbf{Beh}_0\) containing the image of \(K\).

        \item \textbf{Pass to the full subcategory.} The category \(\mathbf{Der}(K)\) is the full internal subcategory on this object subobject. Hence any full derivative-closed internal subcategory containing \(K\) contains \(\operatorname{Der}_0(K)\) on objects. Since it is full, it contains all arrows between those objects, and therefore contains the full subcategory \(\mathbf{Der}(K)\). Thus \(\mathbf{Der}(K)\) is the smallest full derivative-closed internal subcategory containing the specified behaviours.
    \end{enumerate}
\end{proof}

\subsection{The syntactic residual comparison category}
\label{subsec:syntactic-residual-comparison-category}

\begin{definition}[Syntactic residual comparison category]
\label{def:syntactic-residual-comparison-category}
    The \textbf{syntactic residual comparison category} of \(K\) is the full inverse-image internal category along the object map \(d_K\):
    \[
        \mathbf{Syn}(K):=d_K^*
        \mathbf{Beh}^{\mathrm{vert}}_{\mathbb A,\mathbb B}.
    \]
    Thus
    \[
        \mathbf{Syn}(K)_0=D_K,
    \]
    and
    \[
        \mathbf{Syn}(K)_1
        =
        D_K
        \times_{\mathbf{Beh}_0,s}
        \mathbf{Beh}_1
        \times_{t,\mathbf{Beh}_0}
        D_K.
    \]
    A generalized arrow
    \[
        (a,\kappa)\to(a',\kappa')
    \]
    is a behaviour comparison
    \[
        \partial_a k(\kappa)
        \Rightarrow
        \partial_{a'} k(\kappa').
    \]
\end{definition}

The inverse-image square on arrow objects may be displayed fibrewise as
\[
    \begin{tikzcd}[column sep=large, row sep=large]
            \mathbf{Syn}(K)_1
                \arrow[r]
                \arrow[d]
                \arrow[dr, phantom, "\lrcorner", very near start]
        &
            \mathbf{Beh}_1
                \arrow[d, "{(s,t)}"]
        \\
            D_K\times_{A_0}D_K
                \arrow[r, "{(d_K,d_K)}"']
        &
            \mathbf{Beh}_0\times_{A_0}\mathbf{Beh}_0.
    \end{tikzcd}
\]

The evident projection functor
\[
    \nu_K:
    \mathbf{Syn}(K)\to
    \mathbf{Beh}^{\mathrm{vert}}_{\mathbb A,\mathbb B}
\]
has object map \(d_K\) and sends each arrow, already a behaviour comparison in \(\mathbf{Beh}^{\mathrm{vert}}_{\mathbb A,\mathbb B}\), to that same comparison.

\begin{remark}
\label{rem:comparison-level-image-full}
    The comparison-level image considered below is the full subcategory on the object image. It is not the pointwise arrow-image of the residual evaluation functor. The map from \(\mathbf{Syn}(K)\) to this full subcategory is regular epimorphic on arrows because \(\mathbf{Syn}(K)\) is defined by pulling back the whole vertical behaviour category along \(d_K\).
\end{remark}

\subsection{The full residual image}
\label{subsec:full-residual-image}

\begin{theorem}[Full residual image on the derivative object image]
\label{thm:comparison-level-residual-image}
    The residual evaluation functor
    \[
        \nu_K:
        \mathbf{Syn}(K)\to
        \mathbf{Beh}^{\mathrm{vert}}_{\mathbb A,\mathbb B}
    \]
    factors as
    \[
        \mathbf{Syn}(K)
        \twoheadrightarrow
        \mathbf{Der}(K)
        \hookrightarrow
        \mathbf{Beh}^{\mathrm{vert}}_{\mathbb A,\mathbb B},
    \]
    where the first functor is regular epimorphic on objects and arrows, and the second functor is the full inclusion.
\end{theorem}
\begin{proof}
    \leavevmode
    \begin{enumerate}
        \item \textbf{Identify the object image.} The object map of \(\nu_K\) is
        \[
            d_K:D_K\to\mathbf{Beh}_0,
        \]
        whose image is
        \[
            \operatorname{Der}_0(K).
        \]

        \item \textbf{Compare the arrow objects.} The arrow object of \(\mathbf{Syn}(K)\) is
        \[
            D_K
            \times_{\mathbf{Beh}_0,s}
            \mathbf{Beh}_1
            \times_{t,\mathbf{Beh}_0}
            D_K,
        \]
        while the arrow object of \(\mathbf{Der}(K)\) is
        \[
            \operatorname{Der}_0(K)
            \times_{\mathbf{Beh}_0,s}
            \mathbf{Beh}_1
            \times_{t,\mathbf{Beh}_0}
            \operatorname{Der}_0(K).
        \]
        The induced arrow map factors as
        \[
            \begin{aligned}
                &
                D_K
                \times_{\mathbf{Beh}_0,s}
                \mathbf{Beh}_1
                \times_{t,\mathbf{Beh}_0}
                D_K
                \\
                &\qquad\twoheadrightarrow
                \operatorname{Der}_0(K)
                \times_{\mathbf{Beh}_0,s}
                \mathbf{Beh}_1
                \times_{t,\mathbf{Beh}_0}
                D_K
                \\
                &\qquad\twoheadrightarrow
                \operatorname{Der}_0(K)
                \times_{\mathbf{Beh}_0,s}
                \mathbf{Beh}_1
                \times_{t,\mathbf{Beh}_0}
                \operatorname{Der}_0(K).
            \end{aligned}
        \]

        \item \textbf{Use stability of regular epimorphisms.} Each map is a pullback of
        \[
            e_K:D_K\twoheadrightarrow\operatorname{Der}_0(K).
        \]
        Regular epimorphisms are stable under pullback and closed under composition, so the arrow map is a regular epimorphism.

        \item \textbf{Display the arrow-level image.} The arrow-level image diagram is
        \[
            \begin{tikzcd}[column sep=large, row sep=large]
                    \mathbf{Syn}(K)_1
                        \arrow[r, two heads]
                        \arrow[d]
                &
                    \mathbf{Der}(K)_1
                        \arrow[d, hook]
                \\
                    \mathbf{Beh}_1
                        \arrow[r, equal]
                &
                    \mathbf{Beh}_1.
            \end{tikzcd}
        \]
    \end{enumerate}
\end{proof}

\section{The one-object specialization}
\label{sec:behaviour-one-object-specialization}

Let \(M\) and \(N\) be monoid objects in \(\mathcal E\), regarded as
one-object internal categories
\[
    \mathbb A=M,
    \qquad
    \mathbb B=N,
\]
using the one-object convention of
Chapter~\ref{ch:spans-internal-categories-mealy}:
\[
    b\circ a=ba.
\]
Thus a Mealy morphism
\[
    M\rightsquigarrow N
\]
is an internal right \(M\)-object \(P\), written
\[
    x\cdot m:=\delta(m,x),
\]
together with an internal \(N\)-valued output cocycle
\[
    \omega:M\times P\to N
\]
satisfying
\[
    x\cdot 1=x,
    \qquad
    x\cdot(mn)=(x\cdot m)\cdot n,
\]
and
\[
    \omega(1,x)=1,
    \qquad
    \omega(mn,x)=\omega(m,x)\omega(n,x\cdot m).
\]

\subsection{Residual right-Cayley syntax}
\label{subsec:one-object-residual-cayley-syntax}

The residual opfibration has a single base object, but its fibre is the
internal residual right-Cayley category of \(M\). Its object object is \(M\).
An object is written \(r\), and an arrow labelled by \(m\) has the form
\[
    rm\longrightarrow r.
\]
Equivalently, the arrow object is \(M\times M\), with source and target maps
\[
    (r,m)\longmapsto rm,
    \qquad
    (r,m)\longmapsto r.
\]
The identity at \(r\) is \((r,1)\). Composition is determined by
\[
    rmn \xrightarrow{\;n\;} rm \xrightarrow{\;m\;} r
    \qquad\longmapsto\qquad
    rmn \xrightarrow{\;mn\;} r,
\]
or, in morphism notation,
\[
    (rm,n),(r,m)\longmapsto (r,mn).
\]

Transport along an input \(s\in M\) is postcomposition in the one-object
category. Hence
\[
    s_*r=sr,
\]
where \(sr\) is the monoid product corresponding to the categorical composite
\[
    s\circ r.
\]
Thus the derivative of a residual behaviour will be given by left translation
of the residual-history variable.

\subsection{Residual behaviours and comparisons}
\label{subsec:one-object-residual-behaviours}

A residual behaviour is an internal functor from the residual right-Cayley
category of \(M\) to \(N\). Since \(N\) has one object, such a functor is
equivalently a map
\[
    \lambda:M\times M\to N
\]
satisfying
\[
    \lambda(r,1)=1,
    \qquad
    \lambda(r,mn)=\lambda(r,m)\lambda(rm,n).
\]
The element \(\lambda(r,m)\) is the output attached to the residual arrow
\[
    rm\longrightarrow r.
\]

A behaviour comparison
\[
    \alpha:\lambda\Rightarrow\mu
\]
is an internal natural transformation between the corresponding residual
cocycle functors. Equivalently, it is a map
\[
    \alpha:M\to N
\]
satisfying
\[
    \alpha_r\lambda(r,m)=\mu(r,m)\alpha_{rm}.
\]
This is the naturality square for the residual arrow \(rm\to r\):
\[
    \begin{tikzcd}[column sep=huge, row sep=large]
            \bullet
                \arrow[r, "{\lambda(r,m)}"]
                \arrow[d, "{\alpha_{rm}}"']
        &
            \bullet
                \arrow[d, "{\alpha_r}"]
        \\
            \bullet
                \arrow[r, "{\mu(r,m)}"']
        &
            \bullet .
    \end{tikzcd}
\]
If \(N=G\) is an internal group, this may equivalently be written as the
nonabelian coboundary formula
\[
    \mu(r,m)=\alpha_r\lambda(r,m)\alpha_{rm}^{-1}.
\]

\subsection{The canonical residual-cocycle machine}
\label{subsec:one-object-canonical-behaviour-machine}

The states of the canonical behaviour Mealy morphism are residual cocycles
\[
    \lambda:M\times M\to N.
\]
For an input \(s\in M\), the transition is derivative reindexing:
\[
    \lambda\longmapsto \partial_s\lambda,
    \qquad
    (\partial_s\lambda)(r,m)=\lambda(sr,m).
\]
The output on input \(s\) is the anchor comparison associated to the residual
arrow
\[
    s\longrightarrow 1
\]
in the residual right-Cayley category. Under the one-object identification of
\(\mathbb B=N\), this comparison is represented by
\[
    \lambda(1,s)\in N.
\]
Thus the canonical behaviour Mealy morphism specializes to the residual-cocycle
machine
\[
    \lambda
        \xmapsto{\;s\;/\;\lambda(1,s)\;}
    \partial_s\lambda .
\]

\subsection{Residual semantics of a one-object Mealy morphism}
\label{subsec:one-object-residual-semantics}

Let
\[
    P:M\rightsquigarrow N
\]
be a one-object Mealy morphism. Its execution category is the internal action
category of the right \(M\)-object \(P\). Its objects are states \(x\in P\),
and an arrow labelled by \(m\in M\) has the form
\[
    x\longrightarrow x\cdot m.
\]
The output functor to \(N^{\mathrm{op}}\) sends this execution arrow to the
element
\[
    \omega(m,x)\in N.
\]
The cocycle law says exactly that this assignment preserves composition:
\[
    x
        \xrightarrow{\;m\;}
    x\cdot m
        \xrightarrow{\;n\;}
    (x\cdot m)\cdot n
\]
has total output
\[
    \omega(m,x)\omega(n,x\cdot m)=\omega(mn,x).
\]

For a state \(x\in P\), the residual behaviour of \(x\) is the residual cocycle
\[
    \lambda_x:M\times M\to N
\]
defined by
\[
    \lambda_x(r,m)=\omega(m,x\cdot r).
\]
Thus \(\lambda_x(r,m)\) is the output produced after first moving from \(x\)
to the residual state \(x\cdot r\), and then reading the residual input \(m\).
The cocycle equations for \(\omega\) imply
\[
    \lambda_x(r,1)=1,
    \qquad
    \lambda_x(r,mn)=\lambda_x(r,m)\lambda_x(rm,n).
\]
The behaviour assignment is therefore
\[
    \beta_P:P\to\mathbf{Beh}_0,
    \qquad
    \beta_P(x)=\lambda_x.
\]

Terminality of the canonical behaviour machine specializes to the statement
that this is the unique strict semantic map from \(P\) to the residual-cocycle
machine. Behavioural separatedness says that
\[
    \beta_P:P\to\mathbf{Beh}_0
\]
is monic, equivalently that states are separated by their future residual
cocycles:
\[
    \lambda_x=\lambda_y
    \quad\Longrightarrow\quad
    x=y.
\]

\subsection{Specifications, Nerode quotients, and generated comparisons}
\label{subsec:one-object-nerode-specialization}

A behaviour specification
\[
    k:K\to\mathbf{Beh}_0
\]
is an internal family of residual \(N\)-valued cocycles. Write
\[
    k(\kappa)=\lambda_\kappa.
\]
The formal derivative orbit is
\[
    D_K=M\times K,
\]
and the residual evaluation map is
\[
    d_K:M\times K\to\mathbf{Beh}_0,
    \qquad
    d_K(s,\kappa)=\partial_s\lambda_\kappa.
\]
Explicitly,
\[
    \bigl(\partial_s\lambda_\kappa\bigr)(r,m)
    =
    \lambda_\kappa(sr,m).
\]

The Nerode relation is the kernel pair of \(d_K\). Hence
\[
    (s,\kappa)\equiv_K(s',\kappa')
\]
precisely when the corresponding residual cocycles are equal:
\[
    \partial_s\lambda_\kappa
    =
    \partial_{s'}\lambda_{\kappa'}.
\]
Equivalently,
\[
    \lambda_\kappa(sr,m)
    =
    \lambda_{\kappa'}(s'r,m)
\]
as internal residual behaviours.

The Nerode quotient
\[
    D_K/{\equiv_K}
\]
is canonically the derivative image
\[
    \operatorname{Der}_0(K)\rightarrowtail\mathbf{Beh}_0.
\]
Thus, in the one-object specialization, the state object of the minimal
residual realization is the residual derivative closure of the specified
cocycles. Its transition is
\[
    \lambda\longmapsto\partial_s\lambda,
\]
and its output on input \(s\) is
\[
    \lambda(1,s)\in N.
\]
Hence \(\mathsf{Min}(K)\) is the smallest residual-cocycle machine containing
the specified cocycles and closed under all residual derivatives.

The generated residual behaviour category
\[
    \mathbf{Der}(K)
\]
is the full residual comparison category on these generated residual cocycles.
Its objects are residual derivatives of the specified cocycles. An arrow
\[
    \lambda\to\mu
\]
is a residual behaviour comparison, equivalently a map
\[
    \alpha:M\to N
\]
satisfying
\[
    \alpha_r\lambda(r,m)=\mu(r,m)\alpha_{rm}.
\]
If \(N=G\) is an internal group, these arrows may be read as nonabelian
residual coboundaries:
\[
    \mu(r,m)=\alpha_r\lambda(r,m)\alpha_{rm}^{-1}.
\]
Thus \(\mathbf{Der}(K)\) records not only the residual derivatives generated by
\(K\), but also all residual coboundary comparisons between them.

In the ordinary external finite-state case, the finite-state theorem says that
a finite realization exists precisely when there are only finitely many
distinct residual cocycles in \(\operatorname{Der}_0(K)\).

\section{The stateless specialization}
\label{sec:behaviour-stateless-specialization}

Let
\[
    F:\mathbb A\to\mathbb B
\]
be an internal functor, regarded as the map-representable Mealy morphism
associated to its object map
\[
    F_0:A_0\to B_0.
\]
Its carrier is
\[
    A_0
        \xleftarrow{\;1_{A_0}\;}
    A_0
        \xrightarrow{\;F_0\;}
    B_0.
\]
Thus a state is simply an input object \(v\in A_0\). There is no additional
state object beyond the object-of-objects of \(\mathbb A\).

For an input arrow
\[
    a:u\to v,
\]
the transition sends the state \(v\) to the state \(u\), and the emitted output
arrow is
\[
    F_1(a):F_0u\to F_0v.
\]
Consequently the execution category has the same arrows as \(\mathbb A\), but
read in the target-to-source execution direction. It is canonically identified
with
\[
    \mathbb A^{\mathrm{op}},
\]
and the output functor is
\[
    F^{\mathrm{op}}:\mathbb A^{\mathrm{op}}\to\mathbb B^{\mathrm{op}}.
\]

\subsection{Slice-restricted residual semantics}
\label{subsec:stateless-slice-residual-semantics}

For an object \(v\in A_0\), the residual behaviour of the stateless state \(v\)
is the internal functor
\[
    H^F_v:\mathbb A/v\to\mathbb B
\]
defined by restriction of \(F\) along the domain functor
\[
    \operatorname{dom}_v:\mathbb A/v\to\mathbb A.
\]
Thus
\[
    H^F_v=F\circ\operatorname{dom}_v.
\]
Explicitly,
\[
    H^F_v(r:u\to v)=F_0(u),
\]
and, for a residual triangle
\[
    w\xrightarrow{a}u\xrightarrow{r}v,
\]
one has
\[
    H^F_v(a,r)=F_1(a):F_0(w)\to F_0(u).
\]

Hence \(F\) determines a canonical map into the behaviour object,
\[
    J_F:A_0\to\mathbf{Beh}_0,
    \qquad
    J_F(v)=H^F_v=F\circ\operatorname{dom}_v.
\]
This is the stateless residual semantics of \(F\). It assigns to an input
object \(v\) the ordinary slice restriction of \(F\) to arrows with codomain
\(v\).

For an input arrow \(a:u\to v\), residual derivative is reindexing along
postcomposition
\[
    a_*:\mathbb A/u\to\mathbb A/v.
\]
The defining square is
\[
    \begin{tikzcd}[column sep=huge, row sep=large]
            \mathbb A/u
                \arrow[r, "a_*"]
                \arrow[d, "\operatorname{dom}_u"']
        &
            \mathbb A/v
                \arrow[d, "\operatorname{dom}_v"]
        \\
            \mathbb A
                \arrow[r, equal]
        &
            \mathbb A.
    \end{tikzcd}
\]
Therefore
\[
    \partial_aH^F_v
    =
    H^F_v\circ a_*
    =
    F\circ\operatorname{dom}_v\circ a_*
    =
    F\circ\operatorname{dom}_u
    =
    H^F_u.
\]
Equivalently,
\[
    \partial_aJ_F(v)=J_F(u).
\]

The emitted output arrow of the canonical behaviour machine at \(J_F(v)\) is
the anchor comparison for \(H^F_v\) along \(a\). Since \(H^F_v\) is the slice
restriction of \(F\), this anchor comparison is exactly
\[
    F_1(a):F_0u\to F_0v.
\]
Thus the canonical behaviour machine contains the stateless submachine
\[
    H^F_v
        \xmapsto{\;a\;/\;F_1(a)\;}
    H^F_u.
\]

The behaviour assignment of the map-representable Mealy morphism associated to
\(F\) is precisely
\[
    \beta_{F_!}=J_F.
\]
Moreover,
\[
    p_{\mathbf{Beh}}J_F=1_{A_0},
\]
so \(J_F\) is monic. Hence stateless map-representable machines are
behaviourally separated.

For \(F=1_{\mathbb A}\), the stateless residual semantics is
\[
    v\longmapsto
    \operatorname{dom}_v:\mathbb A/v\to\mathbb A,
\]
so the canonical behaviour machine recovers the codomain-slice geometry of
\(\mathbb A\) itself.

\subsection{Simulations and residual natural transformations}
\label{subsec:stateless-simulations-specialization}

Let
\[
    F,G:\mathbb A\to\mathbb B
\]
be internal functors, regarded as map-representable Mealy morphisms
\[
    F_!,G_!:\mathbb A\rightsquigarrow\mathbb B.
\]
A Mealy simulation
\[
    F_!\Rightarrow G_!
\]
is exactly an internal natural transformation
\[
    \alpha:F\Rightarrow G.
\]
Thus \(\alpha\) consists of a map
\[
    \alpha:A_0\to B_1
\]
satisfying
\[
    s_B\alpha=F_0,
    \qquad
    t_B\alpha=G_0,
\]
together with the internal naturality equation.

The induced behaviour comparison is obtained by restricting this natural
transformation to each slice. For each \(v\in A_0\), it is a vertical natural
transformation
\[
    \widehat{\alpha}_v:
    H^F_v\Rightarrow H^G_v
\]
between
\[
    H^F_v=F\circ\operatorname{dom}_v,
    \qquad
    H^G_v=G\circ\operatorname{dom}_v.
\]
At an object
\[
    r:u\to v
\]
of the slice \(\mathbb A/v\), its component is
\[
    (\widehat{\alpha}_v)_r
    =
    \alpha_u:F_0(u)\to G_0(u).
\]
For a residual triangle
\[
    w\xrightarrow{a}u\xrightarrow{r}v,
\]
naturality of \(\widehat{\alpha}_v\) is the ordinary naturality square for
\(\alpha\) at \(a:w\to u\):
\[
    \begin{tikzcd}[column sep=huge, row sep=large]
            F_0(w)
                \arrow[r, "F_1(a)"]
                \arrow[d, "\alpha_w"']
        &
            F_0(u)
                \arrow[d, "\alpha_u"]
        \\
            G_0(w)
                \arrow[r, "G_1(a)"']
        &
            G_0(u).
    \end{tikzcd}
\]
Equivalently,
\[
    \widehat{\alpha}_v
    =
    \alpha\ast\operatorname{dom}_v:
    F\circ\operatorname{dom}_v
        \Rightarrow
    G\circ\operatorname{dom}_v.
\]
Thus simulation-induced behaviour comparison specializes to ordinary
restriction of a natural transformation along slices.

\subsection{Stateless specifications and their Nerode quotients}
\label{subsec:stateless-specifications-nerode}

Let
\[
    i:K\to A_0
\]
be a family of specified input objects. The corresponding stateless
specification for \(F\) is
\[
    k:=J_Fi:K\to\mathbf{Beh}_0.
\]
Thus the specified behaviours are
\[
    H^F_{i(\kappa)}
    =
    F\circ\operatorname{dom}_{i(\kappa)}.
\]

The formal derivative orbit is
\[
    D_K
    =
    A_1\times_{t_A,A_0,i}K.
\]
A generalized element is an arrow
\[
    a:u\to i(\kappa)
\]
into one of the specified objects.

The residual evaluation map is
\[
    d_K:D_K\to\mathbf{Beh}_0,
    \qquad
    d_K(a,\kappa)
    =
    \partial_aJ_F(i(\kappa)).
\]
Since
\[
    \partial_aJ_F(i(\kappa))
    =
    J_F(s_Aa),
\]
one has the factorization
\[
    d_K
    =
    J_Fs_A\pi_{A_1}.
\]

Put
\[
    Z:=A_0\times B_0.
\]
Regard \(A_0\) as an object of \(\mathcal E/Z\) through
\[
    \bar F
    :=
    (1_{A_0},F_0):
    A_0\to Z.
\]
The map
\[
    J_F:A_0\to\mathbf{Beh}_0
\]
is a morphism in \(\mathcal E/Z\), because
\[
    p_{\mathbf{Beh}}J_F=1_{A_0},
\]
and
\[
    q_{\mathbf{Beh}}J_F=F_0.
\]
Moreover, \(J_F\) is monic, since \(p_{\mathbf{Beh}}\) is a left inverse:
\[
    p_{\mathbf{Beh}}J_F=1_{A_0}.
\]

The object \(D_K\) is regarded as an object of \(\mathcal E/Z\) through
\[
    (p_{\mathbf{Beh}},q_{\mathbf{Beh}})d_K.
\]
The map
\[
    s_A\pi_{A_1}:D_K\to A_0
\]
is a morphism in \(\mathcal E/Z\). Indeed,
\[
    \begin{aligned}
        (1_{A_0},F_0)s_A\pi_{A_1}
        &=
        (p_{\mathbf{Beh}},q_{\mathbf{Beh}})
        J_Fs_A\pi_{A_1}
        \\
        &=
        (p_{\mathbf{Beh}},q_{\mathbf{Beh}})d_K.
    \end{aligned}
\]

Take the regular epi--mono factorization of
\[
    s_A\pi_{A_1}:D_K\to A_0
\]
in the slice \(\mathcal E/Z\):
\[
    D_K
        \xrightarrow{e_S}
    S
        \xrightarrow{m_S}
    A_0.
\]
Thus \(e_S\) is a regular epimorphism and \(m_S\) is monic in
\(\mathcal E/Z\), and
\[
    m_Se_S=s_A\pi_{A_1}.
\]

The composite
\[
    J_Fm_S:S\to\mathbf{Beh}_0
\]
is monic. Moreover,
\[
    \begin{aligned}
        d_K
        &=
        J_Fs_A\pi_{A_1}
        \\
        &=
        J_Fm_Se_S.
    \end{aligned}
\]
Hence
\[
    D_K
        \xrightarrow{e_S}
    S
        \xrightarrow{J_Fm_S}
    \mathbf{Beh}_0
\]
is a regular epi--mono factorization of \(d_K\) in \(\mathcal E/Z\).

Let
\[
    D_K
        \xrightarrow{e_K}
    \operatorname{Der}_0(K)
        \xrightarrow{j_K}
    \mathbf{Beh}_0
\]
be the regular epi--mono factorization used in
Definition~\ref{def:derivative-image}. By uniqueness of regular images, there
is a unique isomorphism
\[
    \varphi_K:
    S\overset{\cong}{\longrightarrow}\operatorname{Der}_0(K)
\]
in \(\mathcal E/Z\) satisfying
\[
    \varphi_Ke_S=e_K,
\]
and
\[
    j_K\varphi_K=J_Fm_S.
\]
We henceforth use this canonical isomorphism to identify
\[
    \operatorname{Der}_0(K)\cong S.
\]
Under this identification,
\[
    e_K=e_S,
    \qquad
    j_K=J_Fm_S.
\]

The Nerode relation is the kernel pair of \(d_K\). Since
\[
    d_K=J_Fs_A\pi_{A_1}
\]
and \(J_F\) is monic, the kernel pair of \(d_K\) is the kernel pair of
\[
    s_A\pi_{A_1}.
\]
Since
\[
    s_A\pi_{A_1}=m_Se_S
\]
and \(m_S\) is monic, this is in turn the kernel pair of \(e_S\).

Consequently,
\[
    (a,\kappa)\equiv_K(a',\kappa')
\]
if and only if
\[
    s_Aa=s_Aa'.
\]
Since \(e_S\) is the effective quotient of its kernel pair,
\[
    D_K/{\equiv_K}
    \cong
    S
    \cong
    \operatorname{Der}_0(K)
\]
in \(\mathcal E/Z\).

The subobject
\[
    m_S:S\hookrightarrow A_0
\]
is downward closed under domains of arrows. More precisely, there is a unique
map
\[
    \bar s:
    A_1\times_{t_A,A_0,m_S}S
    \to
    S
\]
such that
\[
    m_S\bar s=s_A\pi_{A_1}.
\]

To construct this map, transport the derivative-closure transition
\[
    \bar\partial_K:
    A_1
    \times_{t_A,A_0,p_{\operatorname{Der}}}
    \operatorname{Der}_0(K)
    \to
    \operatorname{Der}_0(K)
\]
of Theorem~\ref{thm:derivative-closure} across the canonical isomorphism
\[
    \varphi_K:S\cong\operatorname{Der}_0(K).
\]
For
\[
    a:u\to v,
    \qquad
    m_S(s)=v,
\]
derivative closure gives
\[
    J_Fm_S(\bar s(a,s))
    =
    \partial_aJ_F(m_Ss).
\]
But
\[
    \partial_aJ_F(m_Ss)
    =
    J_F(s_Aa).
\]
Since \(J_F\) is monic,
\[
    m_S(\bar s(a,s))
    =
    s_Aa.
\]
Thus
\[
    m_S\bar s=s_A\pi_{A_1}.
\]

The minimal residual realization generated by the stateless specification has
carrier
\[
    A_0
    \xleftarrow{m_S}
    S
    \xrightarrow{F_0m_S}
    B_0.
\]
For an arrow
\[
    a:u\to v
\]
with \(v\) represented by a state \(s\in S\), its transition is
\[
    s\longmapsto\bar s(a,s),
\]
whose image in \(A_0\) is \(u\), and its emitted output is
\[
    F_1(a):F_0u\to F_0v.
\]

Thus the minimal residual realization is the restriction of the
map-representable machine associated to \(F\) to the generated
downward-closed object subobject
\[
    m_S:S\hookrightarrow A_0.
\]
Its carrier is map-representable in the sense of
Chapter~\ref{ch:spans-internal-categories-mealy} if and only if
\[
    m_S:S\to A_0
\]
is an isomorphism.

Let
\[
    \mathbb A|_S
\]
denote the full internal subcategory of \(\mathbb A\) on \(S\). Its object
object is \(S\), and its arrow object is
\[
    S
    \times_{m_S,A_0,s_A}
    A_1
    \times_{t_A,A_0,m_S}
    S.
\]
The downward-closure equation
\[
    m_S\bar s=s_A\pi_{A_1}
\]
means that every arrow of \(\mathbb A\) whose target lies in \(S\) also has
its source in \(S\).

The execution category of the minimal residual realization is canonically
\[
    (\mathbb A|_S)^{\mathrm{op}},
\]
and its output functor is the restriction of
\[
    F^{\mathrm{op}}:
    \mathbb A^{\mathrm{op}}\to\mathbb B^{\mathrm{op}}.
\]

The subobject \(S\) is the smallest downward-closed subobject of \(A_0\)
containing the image of \(i\).

First, identities show that the image of \(i\) is contained in \(S\). Indeed,
the unit inclusion
\[
    \iota_K:K\to D_K
\]
satisfies
\[
    s_A\pi_{A_1}\iota_K=i,
\]
and therefore
\[
    i=m_Se_S\iota_K.
\]

Now let
\[
    n:T\hookrightarrow A_0
\]
be any downward-closed subobject through which \(i\) factors. For every
generalized element
\[
    (a,\kappa)\in D_K,
\]
the target \(t_A(a)=i(\kappa)\) lies in \(T\). Downward closure therefore
implies that the source \(s_A(a)\) also lies in \(T\). Consequently the map
\[
    s_A\pi_{A_1}:D_K\to A_0
\]
factors through \(n\). Since \(m_S\) is its regular image, \(m_S\) factors
through \(n\). Thus \(S\) is the smallest such subobject.

For the full specification
\[
    i=1_{A_0}:A_0\to A_0,
\]
there is a canonical isomorphism
\[
    D_K\cong A_1,
\]
under which the map defining \(S\) is
\[
    s_A:A_1\to A_0.
\]
Since
\[
    s_Ae_A=1_{A_0},
\]
the map \(s_A\) is a split epimorphism, and its regular image is all of
\(A_0\). Therefore
\[
    S\cong A_0,
\]
and the minimal residual realization is the original map-representable
machine associated to \(F\).

\section{Examples}
\label{sec:behaviour-myhill-nerode-examples}

\begin{example}[Classical deterministic Mealy machines]
\label{ex:classical-deterministic-mealy-machines}
    Let
    \[
        M=I^*
    \]
    be the free monoid on an input alphabet \(I\), and let \(N\) be an output
    monoid.

    The residual syntax fibre is the residual right-Cayley category of
    \(I^*\). A residual behaviour is a residual cocycle
    \[
        \lambda:I^*\times I^*\to N
    \]
    satisfying
    \[
        \lambda(w,1)=1,
    \]
    and
    \[
        \lambda(w,uv)
        =
        \lambda(w,u)\lambda(wu,v).
    \]

    Since \(I^*\) is free, such a cocycle is determined by its values on
    generators:
    \[
        \bar\lambda:I^*\times I\to N.
    \]
    Equivalently, define
    \[
        o:I^+\to N
    \]
    by
    \[
        o(wi):=\bar\lambda(w,i).
    \]
    The associated residual cocycle is
    \[
        \lambda_o(w,i_1\cdots i_n)
        =
        o(wi_1)
        o(wi_1i_2)
        \cdots
        o(wi_1\cdots i_n),
    \]
    with
    \[
        \lambda_o(w,1)=1.
    \]

    The derivative by \(w\in I^*\) is left translation of the residual
    history:
    \[
        (\partial_w o)(u)=o(wu),
        \qquad
        u\in I^+.
    \]
    The minimal residual realization has state set
    \[
        \{\partial_w o\mid w\in I^*\}.
    \]
    On a generator \(i\in I\), its transition is
    \[
        \partial_w o
        \longmapsto
        \partial_{wi}o,
    \]
    and its output is
    \[
        o(wi)\in N.
    \]

    This describes deterministic Mealy machines whose single-step outputs
    take values in an arbitrary output monoid \(N\).

    To recover the usual letter-output formulation, let \(O\) be an output
    alphabet, take
    \[
        N=O^*,
    \]
    and impose the generator-level condition
    \[
        \bar\lambda(w,i)\in O\subseteq O^*
    \]
    for every \(w\in I^*\) and \(i\in I\). The cumulative output on a word is
    then the concatenation of the individual output letters emitted during
    its sequential processing.

    In either formulation, the finite-state criterion is
    \[
        \{\partial_w o\mid w\in I^*\}
    \]
    finite. In the letter-output specialization, this is the deterministic
    Mealy analogue of the residual finiteness criterion in the classical
    Myhill--Nerode theorem
    \cite{Myhill1957,Nerode1958,Brzozowski1964}.
\end{example}

\begin{example}[Group-valued cumulative output languages]
\label{ex:group-valued-cumulative-output-languages}
    Let \(G\) be a group and let
    \[
        \ell:\Sigma^*\to G
    \]
    be a cumulative output function. It determines an incremental residual cocycle by
    \[
        \lambda_\ell(r,m)=\ell(r)^{-1}\ell(rm).
    \]
    Then
    \[
        \lambda_\ell(r,1)=1,
    \]
    and
    \[
        \lambda_\ell(r,mn)
        =
        \lambda_\ell(r,m)\lambda_\ell(rm,n).
    \]
    Thus group-valued cumulative output languages fit the present Mealy-type semantics after passing to prefix quotients.
\end{example}

\begin{example}[Base-\texorpdfstring{\(r\)}{r} carry machine]
\label{ex:base-r-carry-machine}
    Let
    \[
        M=N=(\mathbb N,+),
        \qquad
        P=\{0,\dots,r-1\}.
    \]
    Define
    \[
        x\cdot n=(x+n)\bmod r,
    \]
    and
    \[
        \omega(n,x)=\left\lfloor\frac{x+n}{r}\right\rfloor.
    \]
    The behaviour of \(x\) is
    \[
        \lambda_x(s,n)
        =
        \omega(n,x\cdot s)
        =
        \left\lfloor
        \frac{(x+s)\bmod r+n}{r}
        \right\rfloor.
    \]
    The derivative closure is
    \[
        \{\lambda_x\mid x\in\{0,\dots,r-1\}\}.
    \]
    The Nerode quotient identifies residual histories with the same carry residue.
\end{example}

\begin{example}[Typed path-category machines]
\label{ex:typed-path-category-machines}
    Let \(G\) and \(H\) be directed graphs, and consider a Mealy morphism
    \[
        \mathbf{Path}(G)\rightsquigarrow\mathbf{Path}(H).
    \]
    For a vertex \(v\) of \(G\), the residual syntax fibre is
    \[
        \mathbf{Path}(G)/v.
    \]
    A residual behaviour is a functor
    \[
        \mathbf{Path}(G)/v\to\mathbf{Path}(H).
    \]
    It assigns coherent output paths in \(H\) to input paths ending at \(v\). Derivatives are reindexing along residual transport by path concatenation. The categorical Myhill--Nerode object is the derivative image of these path-output behaviours. Simulations induce natural transformations between the corresponding residual path-output functors.
\end{example}

\begin{example}[Context restriction and register allocation]
\label{ex:register-allocation-myhill-nerode}
    Consider the register-allocation Mealy morphism whose states are triples
    \[
        (\Gamma,R,\rho),
    \]
    where \(\Gamma\) is a finite variable context, \(R\) is a finite register context, and
    \[
        \rho:\Gamma\cong R
    \]
    is a bijection. For an inclusion
    \[
        \Gamma'\hookrightarrow\Gamma,
    \]
    the residual state is
    \[
        (\Gamma',\rho(\Gamma'),\rho|_{\Gamma'}).
    \]
    The residual behaviour sends the residual object
    \[
        \Gamma'\hookrightarrow\Gamma
    \]
    to the restricted register context
    \[
        \rho(\Gamma')\subseteq R.
    \]
    For a residual triangle
    \[
        \Gamma''\hookrightarrow\Gamma'\hookrightarrow\Gamma,
    \]
    the behaviour arrow is the inclusion
    \[
        \rho(\Gamma'')\hookrightarrow\rho(\Gamma').
    \]
    The generated residual behaviour category records residual restriction behaviours and coherent register-interface comparisons.

    The residual behaviour square is
    \[
        \begin{tikzcd}[column sep=large, row sep=large]
                \Gamma''
                    \arrow[r, hook]
                    \arrow[d, "{\rho|_{\Gamma''}}"', "\sim"]
            &
                \Gamma'
                    \arrow[d, "{\rho|_{\Gamma'}}", "\sim"']
            \\
                \rho(\Gamma'')
                    \arrow[r, hook]
            &
                \rho(\Gamma').
        \end{tikzcd}
    \]
\end{example}

\begin{example}[Poset-valued systems]
\label{ex:poset-valued-systems}
    If \(\mathbb A\) and \(\mathbb B\) are internal posets, then the residual fibre over \(v\) is the downset of objects mapping to \(v\). A residual behaviour over \(v\) is a monotone assignment from this residual downset into \(\mathbb B\). Behaviour comparisons are pointwise inequalities. Thus
    \[
        \mathbf{Beh}^{\mathrm{vert}}_{\mathbb A,\mathbb B}
    \]
    is an internal preorder, and \(\mathbf{Der}(K)\) is the derivative-closed residual preorder generated by the specification.
\end{example}

\begin{example}[Action groupoids and gauge comparisons]
\label{ex:action-groupoids-gauge-comparisons}
    Let
    \[
        X/\!/G
        \qquad\text{and}\qquad
        Y/\!/H
    \]
    be action groupoids. The residual syntax fibre over \(x\in X\) is the category of arrows ending at \(x\). A residual behaviour records coherent \(H\)-transport along residual \(G\)-histories. Behaviour comparisons are gauge transformations between such transport assignments. The generated residual behaviour category is the full residual image of the syntactic residual comparison category in the behaviour category.
\end{example}

\begin{example}[Group-valued coboundary outputs]
\label{ex:group-valued-coboundary-outputs}
    Let \(N=G\) be a group. A behaviour comparison
    \[
        \alpha:\lambda\Rightarrow\mu
    \]
    between residual cocycles is a family
    \[
        \alpha_r\in G
    \]
    satisfying
    \[
        \alpha_r\lambda(r,m)=\mu(r,m)\alpha_{rm}.
    \]
    Equivalently,
    \[
        \mu(r,m)=\alpha_r\lambda(r,m)\alpha_{rm}^{-1}.
    \]
    Thus the vertical behaviour category is a category of residual cocycles and nonabelian coboundary comparisons.
\end{example}

\section{Chapter summary}
\label{sec:behaviour-myhill-nerode-summary}

The constructions of this chapter may be summarized as follows:
\[
\begin{array}{c|l}
\text{fibred structure} & \text{Mealy-theoretic role} \\
\hline
\operatorname{Res}_{\mathbb A}\to\mathbb A
&
\text{split codomain opfibration of residual input syntax}
\\
\mathbb A/v
&
\text{residual histories ending at }v
\\
a_*:\mathbb A/u\to\mathbb A/v
&
\text{cocartesian residual transport by postcomposition}
\\
\mathbf{Beh}_{\mathbb A,\mathbb B}\to\mathbb A
&
\text{behaviour fibration}
\\
\mathbf{Beh}_{\mathbb A,\mathbb B}(v)
&
[\mathbb A/v,\mathbb B]
\\
\partial_aH=a^*H=H\circ a_*
&
\text{cartesian derivative reindexing}
\\
\mathbf{Beh}_0
&
\text{total object of residual behaviours}
\\
\mathbf{Beh}_1
&
\text{vertical behaviour comparisons}
\\
q_{\mathbf{Beh}}(H)=H(1_v)
&
\text{anchor evaluation}
\\
\text{cleavage of }\mathbf{Beh}
&
\text{canonical behaviour Mealy morphism}
\\
\beta_P(x)=\Omega_P^{\mathrm{op}}\rho_x
&
\text{classifying map of residual orbit semantics}
\\
D_K
&
\text{formal derivative orbit of a specification}
\\
d_K:D_K\to\mathbf{Beh}_0
&
\text{fibred residual evaluation}
\\
{\equiv_K}
&
\text{kernel-pair Nerode relation}
\\
D_K/{\equiv_K}\cong \operatorname{Der}_0(K)
&
\text{categorical Myhill--Nerode quotient}
\\
\operatorname{Der}_0(K)
&
\text{object image of the derivative orbit}
\\
\mathsf{Min}(K)
&
\text{minimal residual realization}
\\
\operatorname{Reach}_0(P,K)\twoheadrightarrow \mathsf{Min}(K)
&
\text{semantic quotient of a reachable realization}
\\
\mathbf{Der}(K)
&
\text{generated full residual behaviour category}
\end{array}
\]

The residual syntax of the input category is organized by the split codomain
opfibration
\[
    \operatorname{Res}_{\mathbb A}\to\mathbb A.
\]
Its fibre over \(v\) is the residual slice
\[
    \mathbb A/v,
\]
and transport along \(a:u\to v\) is postcomposition
\[
    a_*:\mathbb A/u\to\mathbb A/v.
\]
Behaviours form the corresponding contravariant behaviour fibration:
\[
    \mathbf{Beh}_{\mathbb A,\mathbb B}(v)
    =
    [\mathbb A/v,\mathbb B].
\]
Derivatives are cartesian reindexing:
\[
    \partial_aH=H\circ a_*.
\]

The cleavage of the behaviour fibration makes the object of behaviours into
the state object of the canonical behaviour Mealy morphism
\[
    \mathbf{Beh}^{\mathrm{can}}_{\mathbb A,\mathbb B}:
    \mathbb A\rightsquigarrow\mathbb B.
\]
The transition is derivative reindexing, and the output is the anchor
comparison
\[
    H(a,1_v):H(a)\to H(1_v).
\]

For any Mealy morphism \(P\), each state \(x\) determines an orbit functor
\[
    \rho_x:\mathbb A/p(x)\to
    \left(\int_{\mathbb A}P\right)^{\mathrm{op}},
\]
and composing this orbit with the output functor gives the residual behaviour
\[
    H_x=\Omega_P^{\mathrm{op}}\rho_x.
\]
Thus the semantic map
\[
    \beta_P:P\to\mathbf{Beh}_0
\]
is determined by fibred residual orbit semantics. The canonical behaviour
machine is terminal in the ordinary category of strict semantic Mealy
homomorphisms.

For a specification
\[
    k:K\to\mathbf{Beh}_0,
\]
the residual evaluation map
\[
    d_K:D_K\to\mathbf{Beh}_0
\]
sends a formal derivative \((a,\kappa)\) to
\[
    \partial_a k(\kappa).
\]
Its kernel pair is the equality-level Nerode relation, and exactness gives
\[
    D_K/{\equiv_K}
    \cong
    \operatorname{Der}_0(K).
\]
The right-hand side is derivative-closed, so the canonical behaviour machine
restricts to it. This restricted machine is the minimal residual realization
\[
    \mathsf{Min}(K).
\]
Every reachable realization maps regularly epimorphically onto
\(\mathsf{Min}(K)\), and every reachable behaviourally separated realization
is isomorphic to it. With a finite-state structure, a finite realization exists
exactly when the Nerode quotient, equivalently the minimal residual
realization, is finite.

The generated residual behaviour category
\[
    \mathbf{Der}(K)
\]
records the vertical behaviour comparisons between generated residual
behaviours. It is the full internal subcategory of
\[
    \mathbf{Beh}^{\mathrm{vert}}_{\mathbb A,\mathbb B}
\]
on the derivative image
\[
    \operatorname{Der}_0(K)\hookrightarrow\mathbf{Beh}_0.
\]
Thus the categorical Myhill--Nerode theorem has an object-level Mealy
component, namely \(\mathsf{Min}(K)\), and a comparison-level fibred component,
namely \(\mathbf{Der}(K)\).
